\documentclass[11pt,openany]{book}
\usepackage[letterpaper,margin=0.78in,headheight=24pt,footskip=28pt]{geometry}
% Inlined BetterGrades environment package

\RequirePackage[T1]{fontenc}
\RequirePackage[utf8]{inputenc}
\RequirePackage{lmodern}
\RequirePackage{microtype}
\RequirePackage{amsmath,amssymb,amsthm,mathtools}
\RequirePackage{booktabs,tabularx,array,longtable,multicol}
\RequirePackage{enumitem}
\RequirePackage{fancyhdr}
\RequirePackage{titlesec}
\RequirePackage{xcolor}
\RequirePackage{tikz}
\usetikzlibrary{arrows.meta,calc,positioning,backgrounds,patterns,decorations.pathreplacing,angles,quotes,intersections,shapes.geometric}
\RequirePackage{pgfplots}
\pgfplotsset{compat=1.18}
\usepgfplotslibrary{groupplots,fillbetween}
\RequirePackage{hyperref}
\RequirePackage[capitalize,noabbrev]{cleveref}
\RequirePackage{needspace}
\RequirePackage{etoolbox}
\RequirePackage{setspace}
\RequirePackage{ragged2e}
\RequirePackage{makecell}
\RequirePackage{graphicx}
\RequirePackage{comment}
\RequirePackage{xparse}
\RequirePackage{framed}
\RequirePackage{pgfkeys}
\RequirePackage{ifthen}
\RequirePackage{marginnote}

\definecolor{BGInk}{HTML}{18242C}
\definecolor{BGGold}{HTML}{B37A22}
\definecolor{BGBlue}{HTML}{2F6178}
\definecolor{BGGreen}{HTML}{477A65}
\definecolor{BGRed}{HTML}{9B4A3F}
\definecolor{BGPurple}{HTML}{735B82}
\definecolor{BGLight}{HTML}{F5F2EB}
\definecolor{BGLine}{HTML}{D8D2C7}
\definecolor{BGGray}{HTML}{5D6870}
\definecolor{BGPaleBlue}{HTML}{EDF4F7}
\definecolor{BGPaleGreen}{HTML}{EFF6F1}
\definecolor{BGPaleGold}{HTML}{FBF5E8}
\definecolor{BGPaleRed}{HTML}{FAF0EE}
\definecolor{BGPalePurple}{HTML}{F5F0F7}

\hypersetup{
  colorlinks=true,
  linkcolor=BGBlue,
  citecolor=BGGreen,
  urlcolor=BGBlue,
  pdftitle={BetterGrades Calculus I: Limits and Continuity},
  pdfauthor={BetterGrades},
  pdfsubject={A web-ready, self-contained tutorial for the first unit of Calculus I}
}

\pagestyle{fancy}
\fancyhf{}
\renewcommand{\chaptername}{Section}
\fancyhead[LE]{\small\color{BGGray}\leftmark}
\fancyhead[RO]{\small\color{BGGray}\rightmark}
\fancyfoot[C]{\small\color{BGGray}\thepage}
\renewcommand{\headrulewidth}{0.35pt}
\renewcommand{\headrule}{\hbox to\headwidth{\color{BGLine}\leaders\hrule height \headrulewidth\hfill}}

\titleformat{\chapter}[display]
  {\normalfont\bfseries\color{BGInk}}
  {\filleft\fontsize{13}{15}\selectfont CALCULUS I \;\textcolor{BGGold}{/}\; SECTION \thechapter}
  {0.4ex}
  {\titlerule[1.1pt]\vspace{1.1ex}\Huge}
  [\vspace{0.8ex}\titlerule]
\titleformat{\section}{\Large\bfseries\color{BGInk}}{\thesection}{0.65em}{}
\titleformat{\subsection}{\large\bfseries\color{BGBlue}}{\thesubsection}{0.6em}{}
\titleformat{\subsubsection}{\normalsize\bfseries\color{BGGreen}}{}{0em}{}

\setlist[itemize]{leftmargin=1.45em,itemsep=0.25em,topsep=0.35em}
\setlist[enumerate]{leftmargin=2.1em,itemsep=0.42em,topsep=0.5em}
\setlist[description]{leftmargin=1.5em,labelindent=0em,itemsep=0.25em}
\setlength{\parindent}{0pt}
\setlength{\parskip}{0.62em}
\setlength{\emergencystretch}{3em}
\setstretch{1.03}

\newtheoremstyle{bgplain}{0.8em}{0.8em}{\itshape}{}{}{.}{0.5em}{}
\theoremstyle{bgplain}
\newtheorem{theorem}{Theorem}[chapter]
\newtheorem{proposition}[theorem]{Proposition}

\newif\ifbgshowwebmarkers
\newif\ifbgshowreveals
\newif\ifbgshowcheckmetadata
\bgshowwebmarkersfalse
\bgshowrevealstrue
\bgshowcheckmetadatafalse

\pgfkeys{/bgblock/.is family,/bgblock,
 title/.store in=\BGBlockTitle,
 title/.default={},
 .unknown/.code={}}
\newcommand{\bgparsetitle}[2]{%
  \def\BGBlockTitle{#2}%
  \pgfkeys{/bgblock/.cd,title={#2},#1}%
}

% A breakable one-sided accent, deliberately not a rectangular box.
\newenvironment{bgleftbar}[1]{%
  \def\FrameCommand{{\color{#1}\vrule width 2.2pt}\hspace{0.8em}}%
  \MakeFramed{\advance\hsize-\width\FrameRestore}%
}{\endMakeFramed}

\newcommand{\bgblockheading}[2]{\noindent{\bfseries\color{#1}#2}\par\smallskip}

\NewDocumentEnvironment{bgconcept}{O{}}{%
  \bgparsetitle{#1}{Plain-language idea}\par\Needspace{4\baselineskip}\begin{bgleftbar}{BGBlue}\bgblockheading{BGBlue}{\BGBlockTitle}%
}{\end{bgleftbar}\par\addvspace{0.45em}}

\NewDocumentEnvironment{bgmethod}{O{}}{%
  \bgparsetitle{#1}{Method}\par\Needspace{4\baselineskip}\begin{bgleftbar}{BGGold}\bgblockheading{BGGold}{\BGBlockTitle}%
}{\end{bgleftbar}\par\addvspace{0.45em}}

\NewDocumentEnvironment{bgdefinition}{O{}}{%
  \bgparsetitle{#1}{Definition}\par\Needspace{4\baselineskip}\begin{bgleftbar}{BGBlue}\bgblockheading{BGBlue}{\BGBlockTitle}%
}{\end{bgleftbar}\par\addvspace{0.45em}}

\NewDocumentEnvironment{bgtheorem}{O{}}{%
  \bgparsetitle{#1}{Theorem}\par\Needspace{4\baselineskip}\begin{bgleftbar}{BGPurple}\bgblockheading{BGPurple}{\BGBlockTitle}%
}{\end{bgleftbar}\par\addvspace{0.45em}}

\newcounter{bgexample}[chapter]
\renewcommand{\thebgexample}{\thechapter.\arabic{bgexample}}
\NewDocumentEnvironment{bgexample}{O{}}{%
  \bgparsetitle{#1}{Worked Example}\refstepcounter{bgexample}\par\Needspace{6\baselineskip}\addvspace{0.9em}%
  \noindent{\color{BGGreen}\rule{\textwidth}{0.8pt}}\par\smallskip
  {\bfseries\color{BGInk}Example \thebgexample. \BGBlockTitle}\par\smallskip
}{\par\smallskip\noindent{\color{BGLine}\rule{\textwidth}{0.35pt}}\par\addvspace{0.75em}}

\NewDocumentEnvironment{bgguided}{O{}}{%
  \bgparsetitle{#1}{Guided Walkthrough}\refstepcounter{bgexample}\par\Needspace{7\baselineskip}\addvspace{0.9em}%
  \noindent{\color{BGGold}\rule{\textwidth}{0.8pt}}\par\smallskip
  {\bfseries\color{BGInk}Guided Example \thebgexample. \BGBlockTitle}\par
  {\small\color{BGGray}Every decision is stated explicitly; later examples become more compact.}\par\smallskip
}{\par\smallskip\noindent{\color{BGLine}\rule{\textwidth}{0.35pt}}\par\addvspace{0.75em}}

\NewDocumentEnvironment{bgsolution}{O{}}{%
  \bgparsetitle{#1}{Solution}\par\Needspace{3\baselineskip}\begin{bgleftbar}{BGGreen}\bgblockheading{BGGreen}{\BGBlockTitle}%
}{\end{bgleftbar}\par\addvspace{0.35em}}

\NewDocumentEnvironment{bgtry}{O{}}{%
  \bgparsetitle{#1}{Try this before reading on}\par\Needspace{3\baselineskip}\begin{bgleftbar}{BGGold}\bgblockheading{BGGold}{\BGBlockTitle}%
}{\end{bgleftbar}\par\addvspace{0.35em}}

\NewDocumentEnvironment{bgmistake}{O{}}{%
  \bgparsetitle{#1}{Common error}\par\Needspace{3\baselineskip}\begin{bgleftbar}{BGRed}\bgblockheading{BGRed}{\BGBlockTitle}%
}{\end{bgleftbar}\par\addvspace{0.35em}}

\NewDocumentEnvironment{bgexamnote}{O{}}{%
  \bgparsetitle{#1}{Exam note}\par\Needspace{3\baselineskip}\begin{bgleftbar}{BGRed}\bgblockheading{BGRed}{\BGBlockTitle}%
}{\end{bgleftbar}\par\addvspace{0.35em}}

\NewDocumentEnvironment{bgquickcheck}{O{}}{%
  \bgparsetitle{#1}{Quick check}\par\Needspace{3\baselineskip}\begin{bgleftbar}{BGGold}\bgblockheading{BGGold}{\BGBlockTitle}%
}{\end{bgleftbar}\par\addvspace{0.35em}}

\NewDocumentEnvironment{bgsummary}{O{}}{%
  \bgparsetitle{#1}{Keep these ideas}\par\Needspace{4\baselineskip}\addvspace{0.7em}%
  \noindent{\color{BGPurple}\rule{\textwidth}{0.65pt}}\par\smallskip
  {\bfseries\color{BGPurple}\BGBlockTitle}\par\smallskip
}{\par\smallskip\noindent{\color{BGPurple}\rule{\textwidth}{0.3pt}}\par\addvspace{0.65em}}

\NewDocumentEnvironment{bghomework}{O{}}{%
  \bgparsetitle{#1}{Homework pattern}\par\Needspace{4\baselineskip}\begin{bgleftbar}{BGBlue}\bgblockheading{BGBlue}{\BGBlockTitle}%
}{\end{bgleftbar}\par\addvspace{0.35em}}

\NewDocumentEnvironment{bggraphspec}{O{}}{%
  \bgparsetitle{#1}{Digital graph specification}\par\Needspace{4\baselineskip}\addvspace{0.65em}%
  {\small\bfseries\color{BGGray}\BGBlockTitle}\par\smallskip\begingroup\small\color{BGGray}
}{\par\endgroup\addvspace{0.6em}}

\NewDocumentEnvironment{bgsource}{O{}}{%
  \bgparsetitle{#1}{Source note}\par\Needspace{3\baselineskip}\begin{bgleftbar}{BGGray}\bgblockheading{BGGray}{\BGBlockTitle}\small
}{\end{bgleftbar}\par\addvspace{0.35em}}

% Exercises and problem pages are semantic rather than merely visual.
\NewDocumentEnvironment{bgexercises}{O{}}{%
  \begin{enumerate}[leftmargin=2.4em,itemsep=0.62em,topsep=0.45em,#1]
}{\end{enumerate}}

\newcounter{bgproblem}[chapter]
\renewcommand{\thebgproblem}{\thechapter.P\arabic{bgproblem}}
\NewDocumentEnvironment{bgproblem}{O{} m}{%
  \refstepcounter{bgproblem}\bgparsetitle{#1}{Practice problem}\par\Needspace{4\baselineskip}%
  \noindent{\bfseries\color{BGInk}Problem \thebgproblem. \BGBlockTitle}\hfill{\scriptsize\color{BGGray}ID: #2}\par\smallskip
}{\par\addvspace{0.7em}}

% Interactive checks: the website maps id/type/answer to a checker. PDF retains the learning content.
\newcounter{bgcheck}[section]
\renewcommand{\thebgcheck}{\thesection.C\arabic{bgcheck}}
\NewDocumentEnvironment{bgcheck}{m m m O{Quick check}}{%
  \refstepcounter{bgcheck}\def\BGCheckId{#1}\def\BGCheckType{#2}\def\BGCheckAnswer{#3}%
  \par\Needspace{6\baselineskip}\addvspace{0.8em}
  \noindent{\color{BGGold}\rule{\textwidth}{0.65pt}}\par\smallskip
  {\bfseries\color{BGInk}#4 \thebgcheck}\hfill{\scriptsize\color{BGGray}Try before revealing}\par\smallskip
  \ifbgshowcheckmetadata{\scriptsize\color{BGGray}Web ID: \texttt{#1}; input: \texttt{#2}; canonical answer: \texttt{#3}.\par\smallskip}\fi
}{\par\smallskip\noindent{\color{BGLine}\rule{\textwidth}{0.3pt}}\par\addvspace{0.7em}}

\NewDocumentEnvironment{bghint}{}{%
  \par\begin{bgleftbar}{BGBlue}\bgblockheading{BGBlue}{Hint (hidden on the website until requested)}\small
}{\end{bgleftbar}\par}

\NewDocumentEnvironment{bgreveal}{}{%
  \ifbgshowreveals\par\begin{bgleftbar}{BGGreen}\bgblockheading{BGGreen}{Answer and reasoning (revealed after an attempt)}\small\fi
}{\ifbgshowreveals\end{bgleftbar}\par\fi}

% Explicit web-page segmentation. Source parsers can split on these environments.
\NewDocumentEnvironment{webpage}{m m m}{%
  \clearpage\phantomsection\hypertarget{web:#1}{}%
  \def\BGCurrentWebSlug{#1}\def\BGCurrentWebTitle{#2}\def\BGCurrentWebType{#3}%
  \ifbgshowwebmarkers
    \noindent{\scriptsize\color{BGGray}\textsc{Web #3} \quad \texttt{/#1}}\par\smallskip
  \fi
}{\par}

\newcommand{\webnext}[2]{%
  \par\addvspace{1.2em}\noindent{\color{BGLine}\rule{\textwidth}{0.35pt}}\par\smallskip
  {\small\color{BGGray}\textbf{Continue in sequence:} \hyperlink{web:#1}{#2}}\par
}

\newcommand{\lessonobjective}[1]{%
  \par\Needspace{4\baselineskip}\addvspace{0.6em}%
  {\bfseries\color{BGInk}After this lesson, you will be able to}\par\smallskip
  #1\par\addvspace{0.5em}%
}
\newcommand{\whyitworks}[1]{\par\smallskip\textbf{Why this works.} #1}
\newcommand{\whattolearn}[1]{\par\smallskip\textbf{What to learn from this.} #1}
\newcommand{\step}[2]{\textbf{Step #1:} #2\par}
\newcommand{\workedreason}[2]{\begin{tabularx}{\textwidth}{@{}>{\raggedright\arraybackslash}p{0.20\textwidth}X@{}}\textbf{Move} & \textbf{Reason}\\\midrule #1 & #2\end{tabularx}}
\newcommand{\difficulty}[1]{\hfill{\small\color{BGGray}Difficulty: #1/5}}
\newcommand{\sourceadapted}[1]{\hfill{\scriptsize\color{BGGray}Problem type inspired by #1.}}
\newcommand{\R}{\mathbb{R}}
\newcommand{\DNE}{\mathrm{DNE}}
\newcommand{\eps}{\varepsilon}
\newcommand{\answerline}{\vspace{1.2em}\hrulefill}
\newcommand{\smallanswerline}{\rule{2.8cm}{0.4pt}}


% Editorial switches. Set web markers false for a clean student PDF.
\bgshowwebmarkersfalse
\bgshowrevealstrue
\bgshowcheckmetadatafalse

\begin{document}
\frontmatter

% --- BEGIN INLINE chapters/frontmatter ---
\begin{titlepage}
\thispagestyle{empty}
\begin{tikzpicture}[remember picture,overlay]
  \fill[BGInk] (current page.south west) rectangle (current page.north east);
  \fill[BGGold] (current page.south west) rectangle ([xshift=0.24in]current page.north west);
  \draw[BGGold,line width=1.2pt] ([xshift=0.85in,yshift=-1.20in]current page.north west) -- ([xshift=-0.85in,yshift=-1.20in]current page.north east);
  \foreach \x/\y in {1.1/1.5,1.6/1.8,2.1/1.4,2.6/2.0,3.1/1.6,3.6/2.15,4.1/1.75,4.6/2.3,5.1/1.9,5.6/2.45,6.1/2.0,6.6/2.55,7.1/2.05}{
    \fill[BGGold!80] ([xshift=\x in,yshift=\y in]current page.south west) circle (1.4pt);
  }
\end{tikzpicture}
\vspace*{1.55in}
{\color{white}\fontsize{17}{21}\selectfont BETTERGRADES CALCULUS I}\par
\vspace{0.28in}
{\color{white}\fontsize{39}{45}\selectfont\bfseries Limits and Continuity}\par
\vspace{0.22in}
{\color{BGGold}\fontsize{21}{25}\selectfont A complete first-unit tutorial}\par
\vspace{0.55in}
{\color{white}\large Learn the idea. See every step. Practice until it sticks.}\par
\vfill
{\color{white}\normalsize Expanded tutorial edition}\par
{\color{white}\small BetterGrades manuscript}\par
\vspace{0.75in}
\end{titlepage}

\begin{webpage}{calculus/limits-and-continuity}{Limits and Continuity Course Overview}{hub}
\chapter*{What This Book Is For}
\addcontentsline{toc}{chapter}{What This Book Is For}

This is not a packet of lecture notes. It is meant to be the thing you read when the lecture moved too quickly, the textbook skipped three algebra steps, the homework answer says only ``6,'' and every explanation online assumes you already understand the part you are trying to learn.

A student who can do the prerequisite algebra in the diagnostic should be able to use this unit to learn the first major topic of Calculus I from beginning to end. The unit is designed to support three levels of understanding at the same time:

\begin{description}
  \item[Plain-language understanding.] What the idea means before the notation begins.
  \item[Homework competence.] A repeatable method for solving the kinds of problems assigned in a first calculus course.
  \item[Exam readiness.] Enough conceptual and algebraic control to handle unfamiliar versions of familiar problems without copying a pattern blindly.
\end{description}

Every major concept therefore appears in four forms:

\begin{enumerate}
  \item a plain-language explanation that assumes no prior calculus vocabulary;
  \item a formal mathematical statement;
  \item a fully worked example with no important algebra steps omitted;
  \item practice organized from warm-up to exam level.
\end{enumerate}

\begin{bgconcept}
Think of this book as a patient tutor who never gets annoyed when you ask, ``Why did that denominator disappear?'' The answer should be visible on the page. If a solution jumps from line one to line five, it has failed at its job, regardless of how elegant the author felt while writing it.
\end{bgconcept}

\section*{How to Study With This Unit}

Do not read mathematics the way you read a news article. A worked solution only becomes useful when you attempt the next move before seeing it.

For each worked example:

\begin{enumerate}
  \item Cover the solution.
  \item Write down what direct substitution gives.
  \item Name the obstacle: ordinary value, \(0/0\), division by zero, infinity over infinity, a jump, or something else.
  \item Choose a method before uncovering the next line.
  \item Compare your work with the solution and identify the first point where your reasoning differs.
\end{enumerate}

For each exercise set, use three passes:

\begin{description}
  \item[Pass 1: Learn.] Use the section and check answers after small groups of problems.
  \item[Pass 2: Remember.] Return later and work without notes.
  \item[Pass 3: Perform.] Work a mixed set under a time limit, with no labels telling you which method to use.
\end{description}

\begin{bgsummary}
Looking at a solution and thinking ``that makes sense'' is not the same as being able to produce the solution. Recognition feels like knowledge because the page is doing half the thinking for you. Practice removes the page.
\end{bgsummary}

\section*{Suggested Pacing}

The unit can be used as a two-week intensive review or a four-week first course unit.

\begin{center}
\begin{tabularx}{\textwidth}{@{}p{0.13\textwidth}X X@{}}
\toprule
\textbf{Session} & \textbf{Main work} & \textbf{Minimum practice}\\
\midrule
1 & Average change, instantaneous change, basic limit language & Section 1 warm-ups 1--10\\
2 & Tables, graphs, function value versus limit & Section 1 core 11--24\\
3 & One-sided limits and failure modes & Section 1 exam practice 25--40\\
4 & Direct substitution and limit laws & Section 2 warm-ups 1--16\\
5 & Factoring and cancellation & Section 2 core 17--34\\
6 & Radicals, complex fractions, absolute values & Section 2 core 35--54\\
7 & Squeeze Theorem and trig limits & Section 3 exercises 1--38\\
8 & Infinite limits and vertical asymptotes & Section 4 exercises 1--26\\
9 & Limits at infinity and radicals & Section 4 exercises 27--50\\
10 & Continuity and discontinuity & Section 5 exercises 1--26\\
11 & Parameter problems and IVT & Section 5 exercises 27--52\\
12 & Formal definition or optional honors track & Section 6 selected exercises\\
13 & Cumulative review & Section 7 review set\\
14 & Practice exam and correction & Practice Exam A or B\\
\bottomrule
\end{tabularx}
\end{center}
\webnext{calculus/limits/prerequisite-diagnostic}{Calculus Limits Prerequisite Diagnostic}
\end{webpage}

\begin{webpage}{calculus/limits/prerequisite-diagnostic}{Calculus Limits Prerequisite Diagnostic}{diagnostic}
\chapter*{Prerequisite Diagnostic}
\addcontentsline{toc}{chapter}{Prerequisite Diagnostic}

Calculus introduces new ideas, but many wrong calculus answers are caused by old algebra. Complete this diagnostic without a calculator. A score below 15 out of 20 does not mean you cannot learn calculus. It means you should repair the identified skill before the notation becomes crowded enough to hide the original mistake.

\section*{Diagnostic Questions}

\begin{bgexercises}
  \item Factor completely: \(x^2-25\).
  \item Factor completely: \(x^3-8\).
  \item Factor completely: \(2x^2-5x-3\).
  \item Simplify and state all excluded values:
  \[
    \frac{x^2-9}{x-3}.
  \]
  \item Simplify:
  \[
    \frac{\frac{1}{x+2}-\frac12}{x}.
  \]
  \item Rationalize the numerator:
  \[
    \frac{\sqrt{x+9}-3}{x}.
  \]
  \item Solve \(|x-4|<0.3\).
  \item Rewrite \(\sqrt{x^2}\) correctly for every real \(x\).
  \item State the domain of \(f(x)=\dfrac{1}{x^2-4}\).
  \item State the domain of \(g(x)=\sqrt{5-x}\).
  \item Evaluate \(\sin(\pi/6)\), \(\cos(\pi/3)\), and \(\tan(\pi/4)\).
  \item Simplify \(1-\cos^2\theta\).
  \item If \(f(x)=x^2-3x\), find \(f(2+h)\).
  \item For the same function, simplify
  \[
    \frac{f(2+h)-f(2)}{h},\qquad h\ne0.
  \]
  \item Solve \(3x+1=10\).
  \item Solve \(2a+1=4-a\).
  \item Find the slope through \((-1,2)\) and \((3,10)\).
  \item Explain the difference between \(x=2\) and \(x\to2\).
  \item Explain why \(\dfrac00\) is not equal to zero.
  \item Sketch any function with a hole at \((1,3)\) and a filled point at \((1,5)\).
\end{bgexercises}

Answers and skill diagnoses appear in \cref{app:diagnostic-answers}.
\webnext{calculus/limits/notation-guide}{Limit Notation Guide}
\end{webpage}

\begin{webpage}{calculus/limits/notation-guide}{Limit Notation Guide}{reference}
\chapter*{Notation You Will Meet}
\addcontentsline{toc}{chapter}{Notation You Will Meet}

\begin{longtable}{@{}p{0.26\textwidth}p{0.66\textwidth}@{}}
\toprule
\textbf{Notation} & \textbf{Meaning}\\
\midrule
\endhead
\(f(a)\) & The actual output of \(f\) at the input \(a\).\\
\(x\to a\) & The input \(x\) approaches \(a\); it need not equal \(a\).\\
\(x\to a^-\) & The input approaches \(a\) using values less than \(a\).\\
\(x\to a^+\) & The input approaches \(a\) using values greater than \(a\).\\
\(\displaystyle\lim_{x\to a}f(x)=L\) & Nearby outputs approach \(L\) as nearby inputs approach \(a\).\\
\(\DNE\) & The requested limit does not exist as one real number.\\
\(+\infty,-\infty\) & Unbounded positive or negative behavior. These symbols are not real numbers.\\
\([a,b]\) & The closed interval including both endpoints.\\
\((a,b)\) & The open interval excluding both endpoints.\\
\(\forall\) & ``For every.'' Used in the formal definition of a limit.\\
\(\exists\) & ``There exists.''\\
\(\eps,\delta\) & Output tolerance and input tolerance in the formal definition.\\
\bottomrule
\end{longtable}

\begin{bgsource}
The exposition in this unit is original. Its scope and sequencing were checked against \emph{Active Calculus}, 2nd edition, especially its treatment of numerical, graphical, and algebraic perspectives on limits. Historical problem types and applications were also checked against the public-domain texts \emph{Calculus Made Easy} by Silvanus P. Thompson and \emph{Elements of the Differential and Integral Calculus} by Granville and Smith. No Active Calculus exercise is reproduced verbatim. Public-domain examples are modernized, renumbered, and reworked when used as inspiration.
\end{bgsource}
\webnext{calculus/limits/why-limits-matter}{Why Limits Matter in Calculus}
\end{webpage}


% --- END INLINE chapters/frontmatter ---

\tableofcontents
\mainmatter

% --- BEGIN INLINE chapters/ch01_meaning ---
\begin{webpage}{calculus/limits/why-limits-matter}{Why Limits Matter in Calculus}{lesson}
\chapter{What a Limit Means}

\section{The Problem Calculus Is Trying to Solve}

\lessonobjective{
Explain the difference between average and instantaneous change; compute average rates on shrinking intervals; interpret a limit as a number approached by those rates.
}

Suppose you walk 12 feet in 3 seconds. Your average speed is
\[
  \frac{12\text{ feet}}{3\text{ seconds}}=4\text{ feet per second}.
\]
Nothing mysterious has happened. Two different times were compared, so the elapsed time in the denominator was not zero.

Now ask a harder question: how fast were you moving at exactly \(t=2\) seconds? A single instant has no duration. If we try to compare the position at \(t=2\) with itself, we get
\[
  \frac{\text{position at }2-\text{position at }2}{2-2}
  =\frac00,
\]
which is undefined.

Calculus solves this problem by refusing to measure over an interval of zero length. Instead, it measures over ordinary intervals and makes those intervals shorter and shorter. If the resulting average rates settle toward one number, that number is the instantaneous rate.

\begin{bgconcept}
Imagine a video of a moving bicycle. One frame does not show speed. But if you compare two frames that are very close together, you can estimate the speed. Use frames closer and closer together, and the estimates may settle toward the speed at the chosen moment. A limit records the number those estimates approach.
\end{bgconcept}

\subsection{Average rate of change}

For a function \(f\), the average rate of change from \(x=a\) to \(x=b\) is
\[
  \boxed{\frac{f(b)-f(a)}{b-a}}.
\]
This is also the slope of the secant line through the points
\[
  (a,f(a))\qquad\text{and}\qquad(b,f(b)).
\]

\begin{bgguided}[title={A straight-moving cart}]
A toy cart has position
\[
  s(t)=3t+1
\]
in feet after \(t\) seconds. Find its average velocity from \(t=2\) to \(t=5\).

\begin{bgsolution}
First find the two positions:
\[
  s(2)=3(2)+1=7,
\]
\[
  s(5)=3(5)+1=16.
\]
The position changed by
\[
  16-7=9\text{ feet}.
\]
The time changed by
\[
  5-2=3\text{ seconds}.
\]
Therefore,
\[
  \frac{s(5)-s(2)}{5-2}
  =\frac{9}{3}
  =\boxed{3\text{ ft/s}}.
\]
The cart moves at a constant rate, so every average velocity is \(3\) ft/s.
\end{bgsolution}
\end{bgguided}

\begin{bgexample}[title={A ball thrown upward}]
A ball has height
\[
  s(t)=64t-16t^2,
\]
where \(s\) is measured in feet and \(t\) in seconds. Find the average velocity from \(t=1\) to \(t=2\).

\begin{bgsolution}
\step{1}{Find the starting height.}
\[
  s(1)=64(1)-16(1)^2=64-16=48.
\]
\step{2}{Find the ending height.}
\[
  s(2)=64(2)-16(2)^2=128-64=64.
\]
\step{3}{Divide the change in height by the change in time.}
\[
  \frac{s(2)-s(1)}{2-1}
  =\frac{64-48}{1}
  =\boxed{16\text{ ft/s}}.
\]
The positive sign says that, on average, the ball moved upward over this interval.
\end{bgsolution}
\end{bgexample}

\subsection{Shrinking the interval}

To estimate the velocity at exactly \(t=1\), compare \(t=1\) with \(t=1+h\), where \(h\) is a small nonzero number. The average velocity is
\[
  \frac{s(1+h)-s(1)}{h}.
\]
Now simplify very carefully:
\begin{align*}
\frac{s(1+h)-s(1)}{h}
&=\frac{64(1+h)-16(1+h)^2-48}{h}\\
&=\frac{64+64h-16(1+2h+h^2)-48}{h}\\
&=\frac{64+64h-16-32h-16h^2-48}{h}\\
&=\frac{32h-16h^2}{h}\\
&=\frac{h(32-16h)}{h}\\
&=32-16h,\qquad h\ne0.
\end{align*}

Notice why \(h\ne0\) matters. We may cancel \(h\) because the shrinking intervals always have nonzero length. We are studying what happens as \(h\) gets close to zero, not substituting zero before simplifying.

\begin{center}
\begin{tabular}{@{}r r@{}}
\toprule
\(h\) & Average velocity \(32-16h\)\\
\midrule
\(1\) & \(16\)\\
\(0.5\) & \(24\)\\
\(0.1\) & \(30.4\)\\
\(0.01\) & \(31.84\)\\
\(0.001\) & \(31.984\)\\
\(-0.001\) & \(32.016\)\\
\(-0.01\) & \(32.16\)\\
\bottomrule
\end{tabular}
\end{center}

The values approach \(32\). We write
\[
  \boxed{\lim_{h\to0}\frac{s(1+h)-s(1)}{h}=32}.
\]
The instantaneous velocity at \(t=1\) is \(32\) ft/s.

\begin{figure}[htbp]
\centering
\begin{tikzpicture}
\begin{axis}[
  width=0.9\textwidth,height=0.50\textwidth,
  xmin=0,xmax=4,ymin=0,ymax=72,
  axis lines=middle,
  xlabel={$t$ (seconds)},ylabel={$s(t)$ (feet)},
  grid=both,minor tick num=1,
  major grid style={draw=BGLine},minor grid style={draw=BGLine!45},
  tick label style={font=\small},label style={font=\small}
]
\addplot[BGInk,very thick,domain=0:4,samples=120]{64*x-16*x^2};
\addplot[only marks,mark=*,mark size=2.5pt,BGGold] coordinates {(1,48)};
\node[anchor=south east,font=\small] at (axis cs:1,48) {$P$};
\addplot[BGRed!55,thick,domain=0.25:2.25]{16*x+32};
\addplot[BGRed!70,thick,domain=0.45:1.8]{24*x+24};
\addplot[BGRed!90,thick,domain=0.60:1.55]{28.8*x+19.2};
\addplot[BGGreen,very thick,domain=0.55:1.55]{32*x+16};
\node[BGGreen,anchor=west,font=\small] at (axis cs:1.55,65.6) {tangent slope $32$};
\end{axis}
\end{tikzpicture}
\caption{Secant lines through \(P=(1,48)\) approach the tangent line as the second point moves toward \(P\).}
\label{fig:secants}
\end{figure}

\begin{bggraphspec}
\textbf{Function:} \(s(t)=64t-16t^2\). \textbf{Window:} \(0\le t\le4\), \(0\le s\le72\). Fix \(P=(1,48)\). Allow a movable point \(Q=(1+h,s(1+h))\) with controls for positive and negative \(h\). Display the secant slope \(32-16h\) and the tangent line \(s-48=32(t-1)\). Include units in all labels.
\end{bggraphspec}

\begin{bgmistake}
The instantaneous rate is not obtained by pretending \(0/0=0\), \(1\), or infinity. The expression \(0/0\) tells us that direct substitution has failed. The limit process asks what the quotient approaches for nonzero values that become arbitrarily small.
\end{bgmistake}

\begin{bgquickcheck}
A particle has position \(p(t)=5t^2\). Find its average velocity from \(t=2\) to \(t=2+h\), simplify, and predict the instantaneous velocity at \(t=2\).

\textbf{Answer.}
\[
\frac{5(2+h)^2-20}{h}
=\frac{20h+5h^2}{h}
=20+5h,
\]
so the instantaneous velocity is \(20\).
\end{bgquickcheck}
\webnext{calculus/limits/what-a-limit-means}{What a Limit Means}
\end{webpage}

\begin{webpage}{calculus/limits/what-a-limit-means}{What a Limit Means}{lesson}
\section{The Basic Language of Limits}

\lessonobjective{
Interpret \(\lim_{x\to a}f(x)=L\) in words; distinguish the input being approached from the output being approached; estimate a limit numerically and graphically.
}

\begin{bgconcept}
When we say ``the limit is 7,'' we are not saying the input becomes 7. We are saying the \emph{output} approaches 7 while the \emph{input} approaches some specified number.
\end{bgconcept}

\begin{bgdefinition}
We write
\[
  \boxed{\lim_{x\to a}f(x)=L}
\]
when the values of \(f(x)\) can be made as close as desired to \(L\) by taking \(x\) sufficiently close to \(a\), with \(x\ne a\).
\end{bgdefinition}

Read
\[
  \lim_{x\to3}f(x)=7
\]
as

\begin{quote}
``The limit of \(f(x)\) as \(x\) approaches \(3\) is \(7\).''
\end{quote}

The roles are:
\[
\underbrace{x}_{\text{changing input}}
\to
\underbrace{3}_{\text{input target}},
\qquad
\underbrace{f(x)}_{\text{changing output}}
\to
\underbrace{7}_{\text{output target}}.
\]

\begin{bgguided}[title={A limit with no hole}]
Let \(f(x)=x+2\). What happens to \(f(x)\) as \(x\to3\)?

\begin{bgsolution}
If \(x\) is close to \(3\), then \(x+2\) is close to \(5\):
\[
\begin{array}{c|c}
 x & x+2\\
\hline
2.9&4.9\\
2.99&4.99\\
3.01&5.01\\
3.1&5.1
\end{array}
\]
Therefore,
\[
  \boxed{\lim_{x\to3}(x+2)=5}.
\]
This limit is direct because the function has no break at \(x=3\). Later we will call this continuity.
\end{bgsolution}
\end{bgguided}


% --- BEGIN INLINE checks/limit-continuous-01 ---
\begin{bgcheck}{limit-continuous-01}{integer}{10}[Check your understanding]
Evaluate \(\lim_{x\to3}(x^2+1)\).
\begin{bghint}
This polynomial has no break at \(x=3\).
\end{bghint}
\begin{bgreveal}
Substitute: \(3^2+1=10\).
\end{bgreveal}
\end{bgcheck}

% --- END INLINE checks/limit-continuous-01 ---

\webnext{calculus/limits/limit-at-a-hole}{Limits at Holes and Undefined Points}
\end{webpage}

\begin{webpage}{calculus/limits/limit-at-a-hole}{Limits at Holes and Undefined Points}{lesson}
\subsection{A function can have a limit where it is undefined}

Consider
\[
  f(x)=\frac{x^2-4}{x-2}.
\]
At \(x=2\), the formula is undefined. For \(x\ne2\), however,
\[
  f(x)=\frac{(x-2)(x+2)}{x-2}=x+2.
\]
So the graph is the line \(y=x+2\) with one point removed.

\begin{center}
\begin{tabular}{@{}r r@{}}
\toprule
\(x\) & \(f(x)\)\\
\midrule
1.9 & 3.9\\
1.99 & 3.99\\
1.999 & 3.999\\
2.001 & 4.001\\
2.01 & 4.01\\
2.1 & 4.1\\
\bottomrule
\end{tabular}
\end{center}

Both sides approach \(4\), so
\[
  \boxed{\lim_{x\to2}\frac{x^2-4}{x-2}=4}.
\]

\begin{figure}[htbp]
\centering
\begin{tikzpicture}
\begin{axis}[
  width=0.78\textwidth,height=0.48\textwidth,
  xmin=-1,xmax=5,ymin=0,ymax=7,
  axis lines=middle,xlabel={$x$},ylabel={$f(x)$},
  grid=both,major grid style={draw=BGLine},minor tick num=1,
  minor grid style={draw=BGLine!45},
  xtick={-1,0,1,2,3,4,5},ytick={0,1,2,3,4,5,6,7}
]
\addplot[BGBlue,very thick,domain=-1:1.96,samples=80]{x+2};
\addplot[BGBlue,very thick,domain=2.04:5,samples=80]{x+2};
\addplot[only marks,mark=o,mark size=4pt,very thick,BGRed] coordinates {(2,4)};
\draw[dashed,BGGray] (axis cs:2,0) -- (axis cs:2,4);
\draw[dashed,BGGray] (axis cs:0,4) -- (axis cs:2,4);
\node[anchor=south west,font=\small] at (axis cs:2,4) {hole};
\end{axis}
\end{tikzpicture}
\caption{The function is undefined at \(x=2\), but nearby outputs approach \(4\).}
\label{fig:hole}
\end{figure}

\begin{bgmethod}
To estimate a limit from a table, use inputs on \emph{both} sides of the target. The table should move progressively closer to the target. A table gives evidence, not proof; algebra or a trustworthy graph is usually needed to confirm the pattern.
\end{bgmethod}

\begin{bgexample}[title={Table, graph, and algebra agree}]
Estimate and then evaluate
\[
  \lim_{x\to1}\frac{x^3-1}{x-1}.
\]

\begin{bgsolution}
\textbf{Numerical evidence.}
\[
\begin{array}{c|c}
 x & \dfrac{x^3-1}{x-1}\\
\hline
0.9&2.71\\
0.99&2.9701\\
0.999&2.997001\\
1.001&3.003001\\
1.01&3.0301\\
1.1&3.31
\end{array}
\]
The values appear to approach \(3\).

\textbf{Algebraic confirmation.} Use the difference-of-cubes identity:
\[
  x^3-1=(x-1)(x^2+x+1).
\]
For \(x\ne1\),
\[
  \frac{x^3-1}{x-1}=x^2+x+1.
\]
Therefore,
\[
\begin{aligned}
\lim_{x\to1}\frac{x^3-1}{x-1}
&=\lim_{x\to1}(x^2+x+1)\\
&=1^2+1+1\\
&=\boxed{3}.
\end{aligned}
\]
\end{bgsolution}
\end{bgexample}
\webnext{calculus/limits/function-value-vs-limit}{Function Value vs. Limit}
\end{webpage}

\begin{webpage}{calculus/limits/function-value-vs-limit}{Function Value vs. Limit}{lesson}
\section{The Function Value and the Limit Are Different Questions}

\lessonobjective{
Determine \(f(a)\) and \(\lim_{x\to a}f(x)\) separately; explain why changing one isolated function value does not change the limit.
}

The expression \(f(a)\) asks:

\begin{quote}
What output has the function assigned at the single input \(a\)?
\end{quote}

The expression
\[
  \lim_{x\to a}f(x)
\]
asks:

\begin{quote}
What value do nearby outputs approach when nearby inputs approach \(a\)?
\end{quote}

These answers can agree, disagree, or one may fail to exist.

\begin{bgguided}[title={The dot and the road}]
Imagine a road drawn toward the point \((2,5)\), but someone places a separate dot at \((2,9)\). The road tells you where nearby points are going. The separate dot tells you the actual value at \(x=2\).

Thus,
\[
  \lim_{x\to2}f(x)=5,
  \qquad
  f(2)=9.
\]
The limit follows the road, not the misplaced dot.
\end{bgguided}


% --- BEGIN INLINE checks/limit-hole-01 ---
\begin{bgcheck}{limit-hole-01}{integer}{4}[Check your understanding]
Evaluate \(\lim_{x\to2}\frac{x^2-4}{x-2}\).
\begin{bghint}
Factor the difference of squares.
\end{bghint}
\begin{bgreveal}
For \(x\ne2\), the expression is \(x+2\), so the limit is \(4\).
\end{bgreveal}
\end{bgcheck}

% --- END INLINE checks/limit-hole-01 ---


\begin{bgexample}[title={A piecewise value does not control the limit}]
Let
\[
  g(x)=
  \begin{cases}
    x^2+1,&x\ne2,\\
    9,&x=2.
  \end{cases}
\]
Find \(g(2)\) and \(\lim_{x\to2}g(x)\).

\begin{bgsolution}
The piecewise definition directly states
\[
  \boxed{g(2)=9}.
\]
For all nearby inputs \(x\ne2\), the rule is \(g(x)=x^2+1\). Therefore,
\[
\begin{aligned}
\lim_{x\to2}g(x)
&=\lim_{x\to2}(x^2+1)\\
&=2^2+1\\
&=\boxed{5}.
\end{aligned}
\]
So
\[
  g(2)\ne\lim_{x\to2}g(x).
\]
\end{bgsolution}
\end{bgexample}

\begin{figure}[htbp]
\centering
\begin{tikzpicture}
\begin{axis}[
  width=0.76\textwidth,height=0.49\textwidth,
  xmin=-1,xmax=4,ymin=0,ymax=11,
  axis lines=middle,xlabel={$x$},ylabel={$g(x)$},
  grid=both,major grid style={draw=BGLine},minor tick num=1,
  minor grid style={draw=BGLine!45}
]
\addplot[BGBlue,very thick,domain=-1:1.96,samples=100]{x^2+1};
\addplot[BGBlue,very thick,domain=2.04:3.15,samples=100]{x^2+1};
\addplot[only marks,mark=o,mark size=4pt,very thick,BGRed] coordinates {(2,5)};
\addplot[only marks,mark=*,mark size=3pt,BGGold] coordinates {(2,9)};
\node[anchor=west,font=\small] at (axis cs:2.08,5) {$\lim g(x)=5$};
\node[anchor=west,font=\small] at (axis cs:2.08,9) {$g(2)=9$};
\end{axis}
\end{tikzpicture}
\caption{The nearby parabola determines the limit; the filled point determines the function value.}
\label{fig:value-v-limit}
\end{figure}

\begin{bgsummary}
A single point has no control over a two-sided limit. You may change \(f(a)\), delete it, or define it later, and the limit remains the same as long as all nearby values with \(x\ne a\) remain unchanged.
\end{bgsummary}

\begin{bgquickcheck}
Suppose \(f(x)=3x-1\) for every \(x\ne4\), while \(f(4)=-100\). Find
\[
  f(4)
  \qquad\text{and}\qquad
  \lim_{x\to4}f(x).
\]

\textbf{Answer.} \(f(4)=-100\), but
\[
  \lim_{x\to4}f(x)=3(4)-1=11.
\]
\end{bgquickcheck}
\webnext{calculus/limits/one-sided-limits}{One-Sided Limits From the Left and Right}
\end{webpage}

\begin{webpage}{calculus/limits/one-sided-limits}{One-Sided Limits From the Left and Right}{lesson}
\section{One-Sided Limits}

\lessonobjective{
Read and compute left-hand and right-hand limits; use them to decide whether a two-sided limit exists.
}

An input can approach \(a\) from values smaller than \(a\) or from values greater than \(a\).

\begin{bgdefinition}[title={One-Sided Limit Notation}]
\[
  \lim_{x\to a^-}f(x)=L
\]
means that \(f(x)\) approaches \(L\) as \(x\) approaches \(a\) using values \(x<a\).

\[
  \lim_{x\to a^+}f(x)=L
\]
means that \(f(x)\) approaches \(L\) as \(x\) approaches \(a\) using values \(x>a\).
\end{bgdefinition}

The minus and plus signs are directions, not arithmetic operations.

\begin{bgconcept}
Stand at the doorway \(x=a\). Approaching from the hallway on the left is \(x\to a^-\). Approaching from the hallway on the right is \(x\to a^+\). A two-sided meeting happens only if both groups arrive at the same height.
\end{bgconcept}

\begin{bgtheorem}[title={One-Sided Test for a Two-Sided Limit}]
\[
  \boxed{\lim_{x\to a}f(x)=L}
\]
if and only if
\[
  \boxed{\lim_{x\to a^-}f(x)=L
  \quad\text{and}\quad
  \lim_{x\to a^+}f(x)=L.}
\]
If the one-sided limits disagree, the two-sided limit does not exist.
\end{bgtheorem}

\begin{bgguided}[title={Two staircases}]
Suppose the graph approaches height \(3\) from the left and height \(3\) from the right. Then the two-sided limit is \(3\).

If the graph approaches height \(3\) from the left but height \(4\) from the right, there is no single answer to ``what height does the graph approach?'' The two-sided limit does not exist.
\end{bgguided}


% --- BEGIN INLINE checks/left-right-01 ---
\begin{bgcheck}{left-right-01}{choice}{DNE}[Check your understanding]
The left-hand limit is \(2\) and the right-hand limit is \(5\). What is the two-sided limit?
\begin{bghint}
A two-sided limit exists only when the sides agree.
\end{bghint}
\begin{bgreveal}
The two-sided limit is DNE.
\end{bgreveal}
\end{bgcheck}

% --- END INLINE checks/left-right-01 ---


\begin{bgexample}[title={A piecewise function with matching sides}]
Let
\[
  f(x)=
  \begin{cases}
    x+1,&x<2,\\
    5-x,&x\ge2.
  \end{cases}
\]
Find the left-hand, right-hand, and two-sided limits at \(x=2\).

\begin{bgsolution}
For inputs to the left of \(2\), use \(x+1\):
\[
  \lim_{x\to2^-}f(x)=2+1=3.
\]
For inputs to the right of \(2\), use \(5-x\):
\[
  \lim_{x\to2^+}f(x)=5-2=3.
\]
The one-sided limits agree, so
\[
  \boxed{\lim_{x\to2}f(x)=3}.
\]
The actual value is also \(f(2)=3\), but that equality is not what made the two-sided limit exist. The matching one-sided limits did.
\end{bgsolution}
\end{bgexample}

\begin{bgexample}[title={A jump}]
Let
\[
  g(x)=
  \begin{cases}
    x+1,&x<2,\\
    6-x,&x\ge2.
  \end{cases}
\]
Find the one-sided and two-sided limits at \(x=2\).

\begin{bgsolution}
From the left,
\[
  \lim_{x\to2^-}g(x)=2+1=3.
\]
From the right,
\[
  \lim_{x\to2^+}g(x)=6-2=4.
\]
Because \(3\ne4\),
\[
  \boxed{\lim_{x\to2}g(x)=\DNE}.
\]
This is a jump discontinuity. The graph jumps from a left-hand height of \(3\) to a right-hand height of \(4\).
\end{bgsolution}
\end{bgexample}

\begin{figure}[htbp]
\centering
\begin{tikzpicture}
\begin{axis}[
  width=0.78\textwidth,height=0.48\textwidth,
  xmin=-1,xmax=6,ymin=0,ymax=7,
  axis lines=middle,xlabel={$x$},ylabel={$g(x)$},
  grid=both,major grid style={draw=BGLine},minor tick num=1,
  minor grid style={draw=BGLine!45}
]
\addplot[BGBlue,very thick,domain=-1:1.999,samples=100]{x+1};
\addplot[BGGreen,very thick,domain=2:6,samples=100]{6-x};
\addplot[only marks,mark=o,mark size=4pt,very thick,BGBlue] coordinates {(2,3)};
\addplot[only marks,mark=*,mark size=3pt,BGGreen] coordinates {(2,4)};
\draw[dashed,BGGray] (axis cs:2,0) -- (axis cs:2,4);
\node[anchor=east,font=\small] at (axis cs:1.94,3) {left $\to3$};
\node[anchor=west,font=\small] at (axis cs:2.06,4) {right $\to4$};
\end{axis}
\end{tikzpicture}
\caption{Unequal one-sided limits produce a jump and no two-sided limit.}
\label{fig:jump}
\end{figure}

\begin{bgexamnote}
A filled point at \(x=a\) does not repair unequal one-sided limits. Even if you define \(g(2)=3.5\), the left side still approaches \(3\) and the right side still approaches \(4\). No single dot can persuade two disagreeing neighborhoods to cooperate.
\end{bgexamnote}
\webnext{calculus/limits/when-a-limit-does-not-exist}{When a Limit Does Not Exist}
\end{webpage}

\begin{webpage}{calculus/limits/when-a-limit-does-not-exist}{When a Limit Does Not Exist}{lesson}
\section{How a Limit Can Fail}

\lessonobjective{
Recognize and explain the major reasons a limit fails to exist: unequal one-sided limits, unbounded behavior, oscillation, or missing domain on one side.
}

Writing \(\DNE\) is not a full explanation. A good solution identifies the behavior responsible.

\subsection{Failure mode 1: the sides disagree}

This is the jump behavior already seen:
\[
  \lim_{x\to a^-}f(x)\ne\lim_{x\to a^+}f(x).
\]

\subsection{Failure mode 2: unbounded behavior}

For
\[
  f(x)=\frac1x,
\]
as \(x\to0^+\), the outputs become arbitrarily large and positive. As \(x\to0^-\), they become arbitrarily negative. We write
\[
  \lim_{x\to0^+}\frac1x=+\infty,
  \qquad
  \lim_{x\to0^-}\frac1x=-\infty.
\]
The two-sided limit does not exist as a real number.

\subsection{Failure mode 3: endless oscillation}

For
\[
  f(x)=\sin\left(\frac1x\right),
\]
the argument \(1/x\) becomes arbitrarily large as \(x\to0\). The sine function cycles between \(-1\) and \(1\) infinitely often, so the outputs do not approach one number.

\begin{figure}[htbp]
\centering
\begin{tikzpicture}
\begin{axis}[
  width=0.88\textwidth,height=0.42\textwidth,
  xmin=-0.25,xmax=0.25,ymin=-1.25,ymax=1.25,
  axis lines=middle,xlabel={$x$},ylabel={$\sin(1/x)$},
  samples=220,domain=-0.25:-0.01,
  grid=major,major grid style={draw=BGLine}
]
\addplot[BGBlue,thick,domain=-0.25:-0.008,samples=260]{sin(deg(1/x))};
\addplot[BGBlue,thick,domain=0.008:0.25,samples=260]{sin(deg(1/x))};
\end{axis}
\end{tikzpicture}
\caption{The function \(\sin(1/x)\) oscillates increasingly rapidly near zero and has no limit there.}
\label{fig:oscillation}
\end{figure}

\subsection{Failure mode 4: the function exists only on one side}

The function \(f(x)=\sqrt{x}\) has no real values for \(x<0\). Therefore, the right-hand limit
\[
  \lim_{x\to0^+}\sqrt{x}=0
\]
exists, but an ordinary two-sided real limit at \(0\) is not available because there is no left-side domain near zero.

At an endpoint of a domain or interval, a one-sided limit is often exactly the correct question. Later, continuity at endpoints will also use one-sided limits.

\begin{bghomework}
When asked why a limit does not exist, use one of these sentence frames:
\begin{itemize}
  \item ``The left-hand limit is \(A\) and the right-hand limit is \(B\), and \(A\ne B\).''
  \item ``The function is unbounded near the target input.''
  \item ``The function oscillates without approaching one output.''
  \item ``The function has no domain values on one required side of the target.''
\end{itemize}
\end{bghomework}
\webnext{calculus/limits/read-limits-from-graphs}{How to Read Limits From a Graph}
\end{webpage}

\begin{webpage}{calculus/limits/read-limits-from-graphs}{How to Read Limits From a Graph}{lesson}
\section{A Reliable Method for Reading a Graph}

Graph questions are often rushed and surprisingly easy to misread. At each marked input \(x=a\), use the same order.

\begin{bgmethod}[title={Graph-Reading Routine}]
\begin{enumerate}
  \item Trace the curve toward \(a\) from the left. Record the height approached.
  \item Trace the curve toward \(a\) from the right. Record the height approached.
  \item Compare the one-sided limits.
  \item Only then inspect the filled point to find \(f(a)\).
\end{enumerate}
\end{bgmethod}

\begin{bgexample}[title={Reading four quantities from one graph}]
Suppose a graph has an open circle at \((3,6)\), the curve approaches that circle from both sides, and a filled point appears at \((3,-1)\). Find
\[
  \lim_{x\to3^-}f(x),\quad
  \lim_{x\to3^+}f(x),\quad
  \lim_{x\to3}f(x),\quad
  f(3).
\]

\begin{bgsolution}
Following the curve from the left gives
\[
  \lim_{x\to3^-}f(x)=6.
\]
Following the curve from the right gives
\[
  \lim_{x\to3^+}f(x)=6.
\]
The two sides agree, so
\[
  \lim_{x\to3}f(x)=6.
\]
The filled point gives the actual value:
\[
  f(3)=-1.
\]
Thus,
\[
  \boxed{6,\ 6,\ 6,\ -1}.
\]
\end{bgsolution}
\end{bgexample}
\webnext{calculus/limits/direct-substitution}{Direct Substitution for Limits}
\end{webpage}

\begin{webpage}{calculus/limits/meaning-review}{Limit Meaning Review}{review}
\section*{Section 1 Summary}
\addcontentsline{toc}{section}{Section 1 Summary}

\begin{bgsummary}
\begin{itemize}
  \item A limit describes nearby behavior, not necessarily the value at the target point.
  \item \(f(a)\) and \(\lim_{x\to a}f(x)\) are separate questions.
  \item A two-sided limit exists exactly when both one-sided limits exist and agree.
  \item The major failure modes are a jump, unbounded behavior, oscillation, or missing domain on one side.
  \item Tables suggest. Graphs display. Algebra confirms.
\end{itemize}
\end{bgsummary}
\webnext{calculus/limits/direct-substitution}{Direct Substitution for Limits}
\end{webpage}

\begin{webpage}{calculus/limits/meaning-concept-quiz}{Limit Meaning Concept Quiz}{quiz}
\section*{Section 1 Concept Quiz}
\addcontentsline{toc}{section}{Section 1 Concept Quiz}

Use this as a short retrieval check after finishing the section. On the website, each item should accept an answer before revealing the hint and worked explanation.

\begin{bgcheck}{limit-meaning-q1}{integer}{7}[Concept check]
If \(\lim_{x\to 4}f(x)=7\), what output value is approached?
\begin{bghint}
The number after the equals sign is the output target.
\end{bghint}
\begin{bgreveal}
\(7\). The input approaches \(4\), while the output approaches \(7\).
\end{bgreveal}
\end{bgcheck}

\begin{bgcheck}{one-sided-q1}{integer}{3}[Concept check]
A graph approaches \(3\) from the left at \(x=2\). What is \(\lim_{x\to2^-}f(x)\)?
\begin{bghint}
The superscript minus means use inputs smaller than \(2\).
\end{bghint}
\begin{bgreveal}
The left-hand limit is \(3\).
\end{bgreveal}
\end{bgcheck}

\begin{bgcheck}{hole-q1}{integer}{5}[Concept check]
A curve approaches the open point \((1,5)\), while a filled point is at \((1,9)\). What is the limit as \(x\to1\)?
\begin{bghint}
The limit follows the nearby curve, not the filled point.
\end{bghint}
\begin{bgreveal}
The limit is \(5\); the function value is \(9\).
\end{bgreveal}
\end{bgcheck}

\begin{bgcheck}{rate-q1}{integer}{12}[Concept check]
An average-rate expression simplifies to \(12-2h\). What value does it approach as \(h\to0\)?
\begin{bghint}
Set only the limiting value of \(h\) after simplification.
\end{bghint}
\begin{bgreveal}
\(12-2(0)=12\).
\end{bgreveal}
\end{bgcheck}

\begin{bgsummary}[title={Quiz use on the website}]
This page should report completion by concept, not merely a total score. A wrong answer should link back to the exact lesson that teaches the missing method. The same questions may appear in randomized numerical variants, but the canonical explanation should remain stable.
\end{bgsummary}
\webnext{calculus/limits/direct-substitution}{Direct Substitution for Limits}
\end{webpage}

\begin{webpage}{calculus/limits/meaning-practice}{Limit Meaning Practice Problems}{practice}
\section*{Section 1 Exercises}
\addcontentsline{toc}{section}{Section 1 Exercises}

\subsection*{A. Warm-Up: meaning and notation}

\begin{bgexercises}[label=\textbf{1.\arabic*}]
  \item In the statement \(\lim_{x\to4}f(x)=9\), identify the input target and output target.
  \item Write in symbols: ``The limit of \(g(t)\) as \(t\) approaches \(2\) is \(-5\).''
  \item Write in words: \(\lim_{h\to0}Q(h)=12\).
  \item True or false: if \(\lim_{x\to a}f(x)=L\), then \(f(a)=L\). Explain.
  \item True or false: a function must be defined at \(a\) for \(\lim_{x\to a}f(x)\) to exist. Explain.
  \item What does the superscript minus mean in \(x\to3^-\)?
  \item What does the superscript plus mean in \(x\to3^+\)?
  \item Explain why \(x\to2\) does not mean \(x=2\).
  \item A graph approaches \(7\) from both sides at \(x=1\), but \(f(1)=10\). Find the limit.
  \item A graph has no filled point at \(x=-2\), but both sides approach \(4\). Find \(f(-2)\) and the limit.
\end{bgexercises}

\subsection*{B. Average and instantaneous change}

\begin{bgexercises}[label=\textbf{1.\arabic*},resume]
  \item A car's position is \(s(t)=5t+2\). Find its average velocity on \([1,4]\).
  \item A particle's position is \(s(t)=t^2\). Find its average velocity on \([3,3+h]\), simplify, and predict the instantaneous velocity at \(t=3\).
  \item A ball's height is \(s(t)=80t-16t^2\). Find its average velocity on \([2,2+h]\) and its instantaneous velocity at \(t=2\).
  \item A population model is \(P(t)=1000+50t+2t^2\). Find the average rate of change on \([4,4+h]\) and predict \(P'(4)\) informally from the limit.
  \item Explain geometrically what the average rate of change represents on a graph.
  \item Explain geometrically what the limiting secant line becomes when the limit exists.
\end{bgexercises}

\subsection*{C. Tables and formulas}

\begin{bgexercises}[label=\textbf{1.\arabic*},resume]
  \item Complete a table near \(x=3\) for \(f(x)=\dfrac{x^2-9}{x-3}\), then estimate the limit.
  \item Complete a table near \(x=2\) for \(g(x)=\dfrac{x^3-8}{x-2}\), then estimate the limit.
  \item For \(h(x)=\dfrac{x^2+x-6}{x-2}\), simplify for \(x\ne2\) and find the limit at \(2\).
  \item For \(p(x)=\dfrac{x^2-1}{x-1}\), find \(p(1)\) and \(\lim_{x\to1}p(x)\).
  \item Define \(q(1)=100\) and \(q(x)=x+1\) for \(x\ne1\). Find \(q(1)\) and the limit at \(1\).
  \item Create two functions with different values at \(x=0\) but the same limit as \(x\to0\).
\end{bgexercises}

\subsection*{D. One-sided limits}

\begin{bgexercises}[label=\textbf{1.\arabic*},resume]
  \item Let \(f(x)=\begin{cases}2x+1,&x<1,\\x+3,&x\ge1.\end{cases}\) Find both one-sided limits and the two-sided limit at \(1\).
  \item Let \(g(x)=\begin{cases}x^2,&x\le2,\\3x-1,&x>2.\end{cases}\) Find both one-sided limits at \(2\).
  \item Let \(h(x)=\begin{cases}4,&x<0,\\4+x,&x\ge0.\end{cases}\) Does the limit at \(0\) exist?
  \item Let \(k(x)=\begin{cases}-1,&x<3,\\2,&x\ge3.\end{cases}\) Explain why the limit at \(3\) does not exist.
  \item Find \(\lim_{x\to0^-}|x|/x\) and \(\lim_{x\to0^+}|x|/x\).
  \item Find \(\lim_{x\to2^-}|x-2|/(x-2)\) and the corresponding right-hand limit.
\end{bgexercises}

\subsection*{E. Failure modes and reasoning}

\begin{bgexercises}[label=\textbf{1.\arabic*},resume]
  \item State the reason \(\lim_{x\to0}1/x\) does not exist as a real number.
  \item State the reason \(\lim_{x\to0}\sin(1/x)\) does not exist.
  \item Does \(\lim_{x\to0^+}\sqrt{x}\) exist? Does the ordinary two-sided limit exist in the real domain?
  \item A student says, ``The limit is \(7\) because the filled dot is at \((2,7)\).'' What information is missing?
  \item A student says, ``The limit does not exist because the function is undefined at the point.'' Give a counterexample.
  \item Construct a function whose left-hand limit at \(1\) is \(2\), right-hand limit is \(5\), and function value is \(-3\).
  \item Construct a function that is undefined at \(4\) but has limit \(9\) there.
  \item Explain why a finite table can suggest but cannot prove the existence of a limit for every possible function.
\end{bgexercises}

\subsection*{F. Exam-Level Mixed Questions}

\begin{bgexercises}[label=\textbf{1.\arabic*},resume]
  \item Suppose
  \[
    f(x)=\begin{cases}
      x^2-1,&x<2,\\
      4,&x=2,\\
      2x-1,&x>2.
    \end{cases}
  \]
  Find both one-sided limits, the two-sided limit, and \(f(2)\).
  \item Suppose \(f(x)=x+4\) for every \(x\ne-1\), and \(f(-1)\) is not given. Determine the limit and list every possible value of \(f(-1)\) consistent with that limit.
  \item A moving object's average velocity from \(t=5\) to \(t=5+h\) simplifies to \(18-3h\). Find its instantaneous velocity at \(t=5\), and explain the role of the limit.
  \item Give a graph description for a function satisfying
  \[
    f(0)=2,\qquad
    \lim_{x\to0^-}f(x)=1,\qquad
    \lim_{x\to0^+}f(x)=1.
  \]
  \item Give a graph description for a function satisfying
  \[
    f(0)=2,\qquad
    \lim_{x\to0^-}f(x)=1,\qquad
    \lim_{x\to0^+}f(x)=3.
  \]
  \item Explain why the first function in the previous problem has a two-sided limit and the second does not.
\end{bgexercises}

Answers begin in \cref{app:chapter-answers}.
\webnext{calculus/limits/direct-substitution}{Direct Substitution for Limits}
\end{webpage}


% --- END INLINE chapters/ch01_meaning ---


% --- BEGIN INLINE chapters/ch02_finite ---
\begin{webpage}{calculus/limits/direct-substitution}{Direct Substitution for Limits}{lesson}
\chapter{Evaluating Finite Limits}

\section{Start With Direct Substitution}

\lessonobjective{
Evaluate limits of familiar continuous functions by substitution; interpret the result of substitution before choosing a more complicated method.
}

Once you understand what a limit means, most homework problems become a method-selection problem. The first move should almost always be the simplest one:

\begin{bgmethod}[title={First Move for Nearly Every Limit}]
Substitute the target input into the expression.

\begin{itemize}
  \item If you get an ordinary real number, that is usually the limit.
  \item If you get \(0/0\), simplify the expression and try again.
  \item If you get a nonzero number divided by zero, analyze one-sided infinite behavior.
  \item If the input approaches infinity and you get an indeterminate form such as \(\infty/\infty\), compare dominant terms.
\end{itemize}
\end{bgmethod}

Why does substitution work for so many functions? Because polynomials, rational functions away from zero denominators, root functions on their domains, and trigonometric functions on their domains are continuous. Continuity will be developed carefully in Section 5. For now, think of a continuous graph as one with no break at the target point.

\begin{bgguided}[title={Substitute and stop}]
Evaluate
\[
  \lim_{x\to2}(x+5).
\]

\begin{bgsolution}
Replace \(x\) by \(2\):
\[
  2+5=7.
\]
Nothing breaks, no denominator becomes zero, and no special method is needed. Therefore,
\[
  \boxed{\lim_{x\to2}(x+5)=7}.
\]
\end{bgsolution}
\end{bgguided}


% --- BEGIN INLINE checks/direct-sub-01 ---
\begin{bgcheck}{direct-sub-01}{integer}{7}[Check your understanding]
Evaluate \(\lim_{x\to2}(3x+1)\).
\begin{bghint}
Substitute because the linear function is continuous.
\end{bghint}
\begin{bgreveal}
\(3(2)+1=7\).
\end{bgreveal}
\end{bgcheck}

% --- END INLINE checks/direct-sub-01 ---


\begin{bgexample}[title={Polynomial substitution}]
Evaluate
\[
  \lim_{x\to-2}(3x^3-4x^2+5x-1).
\]

\begin{bgsolution}
Polynomials are continuous everywhere, so substitute \(x=-2\):
\begin{align*}
3(-2)^3-4(-2)^2+5(-2)-1
&=3(-8)-4(4)-10-1\\
&=-24-16-10-1\\
&=\boxed{-51}.
\end{align*}
\end{bgsolution}
\end{bgexample}

\begin{bgexample}[title={Rational function with a safe denominator}]
Evaluate
\[
  \lim_{x\to3}\frac{x^2+2x-1}{x+4}.
\]

\begin{bgsolution}
Substitute first:
\[
  \frac{3^2+2(3)-1}{3+4}
  =\frac{9+6-1}{7}
  =\frac{14}{7}
  =\boxed{2}.
\]
The denominator approaches \(7\), not zero, so the quotient law is valid.
\end{bgsolution}
\end{bgexample}

\begin{bgexample}[title={A root function}]
Evaluate
\[
  \lim_{x\to5}\sqrt{3x+1}.
\]

\begin{bgsolution}
The expression under the root approaches
\[
  3(5)+1=16.
\]
Therefore,
\[
  \lim_{x\to5}\sqrt{3x+1}=\sqrt{16}=\boxed{4}.
\]
\end{bgsolution}
\end{bgexample}

\begin{bgmistake}
Students sometimes overcomplicate a limit because the section is about limits. If substitution gives a legal real number, do not factor six polynomials, draw a sign chart, or summon l'Hospital's Rule from a future section. Stop.
\end{bgmistake}
\webnext{calculus/limits/limit-laws}{Limit Laws and How to Use Them}
\end{webpage}

\begin{webpage}{calculus/limits/limit-laws}{Limit Laws and How to Use Them}{lesson}
\section{Limit Laws}

\lessonobjective{
Use sums, differences, constant multiples, products, quotients, powers, and roots to build limits from simpler limits.
}

Suppose
\[
  \lim_{x\to a}f(x)=L
  \qquad\text{and}\qquad
  \lim_{x\to a}g(x)=M.
\]
The basic limit laws are:

\begin{bgtheorem}[title={Limit Laws}]
\begin{align*}
\lim_{x\to a}[f(x)+g(x)]&=L+M,\\
\lim_{x\to a}[f(x)-g(x)]&=L-M,\\
\lim_{x\to a}[cf(x)]&=cL,\\
\lim_{x\to a}[f(x)g(x)]&=LM,\\
\lim_{x\to a}\frac{f(x)}{g(x)}&=\frac{L}{M},\qquad M\ne0,\\
\lim_{x\to a}[f(x)]^n&=L^n.
\end{align*}
When the relevant root is defined,
\[
  \lim_{x\to a}\sqrt[n]{f(x)}=\sqrt[n]{L}.
\]
\end{bgtheorem}

\begin{bgconcept}
The laws say that limits cooperate with ordinary arithmetic, provided the arithmetic itself remains legal. If one part approaches \(3\) and another approaches \(5\), their sum approaches \(8\). The only major warning in the basic laws is division: the denominator's limit may not be zero.
\end{bgconcept}

\begin{bgguided}[title={Building a limit from known pieces}]
Suppose
\[
  \lim_{x\to2}f(x)=3
  \qquad\text{and}\qquad
  \lim_{x\to2}g(x)=5.
\]
Find
\[
  \lim_{x\to2}[2f(x)+g(x)].
\]

\begin{bgsolution}
The constant multiple law gives
\[
  \lim_{x\to2}2f(x)=2(3)=6.
\]
Then the sum law gives
\[
  6+5=11.
\]
Therefore,
\[
  \boxed{\lim_{x\to2}[2f(x)+g(x)]=11}.
\]
\end{bgsolution}
\end{bgguided}

\begin{bgexample}[title={Several laws at once}]
Suppose
\[
  \lim_{x\to1}f(x)=-2,
  \qquad
  \lim_{x\to1}g(x)=4.
\]
Evaluate
\[
  \lim_{x\to1}\frac{[f(x)]^2+3g(x)}{g(x)-f(x)}.
\]

\begin{bgsolution}
Apply the laws to each part:
\[
  [f(x)]^2\to(-2)^2=4,
\]
\[
  3g(x)\to3(4)=12,
\]
so the numerator approaches
\[
  4+12=16.
\]
The denominator approaches
\[
  4-(-2)=6.
\]
Because \(6\ne0\), the quotient law applies:
\[
  \boxed{\frac{16}{6}=\frac83}.
\]
\end{bgsolution}
\end{bgexample}

\begin{bgquickcheck}
If \(f(x)\to3\) and \(g(x)\to-1\) as \(x\to4\), evaluate
\[
  \lim_{x\to4}[f(x)g(x)+2f(x)].
\]

\textbf{Answer.} \(3(-1)+2(3)=3\).
\end{bgquickcheck}
\webnext{calculus/limits/zero-over-zero}{What 0/0 Means in a Limit}
\end{webpage}

\begin{webpage}{calculus/limits/zero-over-zero}{What 0/0 Means in a Limit}{lesson}
\section{The Indeterminate Form \texorpdfstring{\(0/0\)}{0/0}}


% --- BEGIN INLINE checks/zero-over-zero-01 ---
\begin{bgcheck}{zero-over-zero-01}{choice}{simplify}[Check your understanding]
Direct substitution gives \(0/0\). Should you report DNE, report 0, or simplify the expression?
\begin{bghint}
\(0/0\) is a diagnosis, not an answer.
\end{bghint}
\begin{bgreveal}
Simplify the expression, then evaluate the limit again.
\end{bgreveal}
\end{bgcheck}

% --- END INLINE checks/zero-over-zero-01 ---


\lessonobjective{
Recognize \(0/0\) as a signal that direct substitution has not determined the limit; choose an algebraic method based on the expression's structure.
}

If direct substitution gives
\[
  \frac00,
\]
then the numerator and denominator both approach zero. The quotient might approach zero, a finite nonzero number, infinity, or no limit at all. The form itself does not decide.

Compare:
\[
  \lim_{x\to0}\frac{x}{x}=1,
\]
\[
  \lim_{x\to0}\frac{x^2}{x}=0,
\]
\[
  \lim_{x\to0}\frac{x}{x^2}
\]
is unbounded, and
\[
  \lim_{x\to0}\frac{|x|}{x}
\]
does not exist. Every expression gives \(0/0\) under direct substitution, but the limits behave differently.

\begin{bgconcept}
The symbol \(0/0\) is like a locked door, not a room number. It says, ``You cannot enter this way.'' It does not tell you what is behind the door. Algebra must change the form of the expression before substitution becomes informative.
\end{bgconcept}

\begin{bghomework}[title={Method Selection After \(0/0\)}]
\begin{center}
\begin{tabularx}{\textwidth}{@{}p{0.35\textwidth}X@{}}
\toprule
\textbf{Visible structure} & \textbf{Method to try}\\
\midrule
Polynomials with common factors & Factor and cancel\\
Square roots or other radicals & Multiply by a conjugate\\
A fraction inside a fraction & Combine the smaller fractions first\\
Absolute values or piecewise rules & Split into one-sided cases\\
Trigonometric expressions near zero & Use identities and fundamental trig limits\\
A bounded oscillating factor & Look for the Squeeze Theorem\\
\bottomrule
\end{tabularx}
\end{center}
\end{bghomework}
\webnext{calculus/limits/factoring-limits}{Evaluating Limits by Factoring}
\end{webpage}

\begin{webpage}{calculus/limits/factoring-limits}{Evaluating Limits by Factoring}{lesson}
\section{Factoring and Cancellation}

\lessonobjective{
Evaluate finite limits by factoring a numerator or denominator, cancelling a common nonzero factor, and then substituting.
}

\begin{bgguided}[title={Why cancellation is legal in a limit}]
Evaluate
\[
  \lim_{x\to2}\frac{x^2-4}{x-2}.
\]

\begin{bgsolution}
Direct substitution gives \(0/0\), so factor:
\[
  x^2-4=(x-2)(x+2).
\]
For every nearby input \(x\ne2\),
\[
  \frac{(x-2)(x+2)}{x-2}=x+2.
\]
Now substitute:
\[
  2+2=4.
\]
Therefore,
\[
  \boxed{\lim_{x\to2}\frac{x^2-4}{x-2}=4}.
\]
We did not claim that the original expression is defined at \(x=2\). We only used the simplified expression for nearby values \(x\ne2\), exactly the values a limit studies.
\end{bgsolution}
\end{bgguided}


% --- BEGIN INLINE checks/factor-flow-01 ---
\begin{bgcheck}{factor-flow-01}{integer}{6}[Check your understanding]
Evaluate \(\lim_{x\to3}\frac{x^2-9}{x-3}\).
\begin{bghint}
Factor \(x^2-9\).
\end{bghint}
\begin{bgreveal}
The expression becomes \(x+3\), so the limit is \(6\).
\end{bgreveal}
\end{bgcheck}

% --- END INLINE checks/factor-flow-01 ---


\begin{bgexample}[title={Factor a quadratic}]
Evaluate
\[
  \lim_{x\to3}\frac{x^2-5x+6}{x-3}.
\]

\begin{bgsolution}
Substitution gives \(0/0\). Factor the numerator:
\[
  x^2-5x+6=(x-2)(x-3).
\]
Cancel the common factor for \(x\ne3\):
\[
  \frac{(x-2)(x-3)}{x-3}=x-2.
\]
Then
\[
  \lim_{x\to3}(x-2)=3-2=\boxed{1}.
\]
\end{bgsolution}
\end{bgexample}

\begin{bgexample}[title={Difference of cubes}]
Evaluate
\[
  \lim_{x\to2}\frac{x^3-8}{x-2}.
\]

\begin{bgsolution}
Use
\[
  a^3-b^3=(a-b)(a^2+ab+b^2).
\]
Thus,
\[
  x^3-8=(x-2)(x^2+2x+4).
\]
For \(x\ne2\),
\[
  \frac{x^3-8}{x-2}=x^2+2x+4.
\]
Now substitute:
\[
  2^2+2(2)+4=4+4+4=\boxed{12}.
\]
\end{bgsolution}
\end{bgexample}

\begin{bgexample}[title={Exam-level: factor more than once}]
Evaluate
\[
  \lim_{x\to1}\frac{x^4-1}{x^2-1}.
\]

\begin{bgsolution}
Substitution gives \(0/0\). Factor both differences of squares:
\[
  x^4-1=(x^2-1)(x^2+1),
\]
so
\[
  \frac{x^4-1}{x^2-1}=x^2+1,\qquad x\ne\pm1.
\]
Therefore,
\[
  \lim_{x\to1}(x^2+1)=1+1=\boxed{2}.
\]
An alternative full factorization is
\[
  x^4-1=(x-1)(x+1)(x^2+1),
\]
\[
  x^2-1=(x-1)(x+1).
\]
Both routes lead to the same simplification.
\end{bgsolution}
\end{bgexample}

\begin{bgexamnote}
Cancel factors, not terms. From
\[
  \frac{x^2+2x}{x}
\]
you may factor the numerator and cancel:
\[
  \frac{x(x+2)}{x}=x+2.
\]
You may not ``cancel the \(x\)'' from only one term and write \(x^2+2\). Addition prevents cancellation until a common factor has been extracted.
\end{bgexamnote}
\webnext{calculus/limits/rationalizing-radicals}{Evaluating Radical Limits With Conjugates}
\end{webpage}

\begin{webpage}{calculus/limits/rationalizing-radicals}{Evaluating Radical Limits With Conjugates}{lesson}
\section{Rationalizing Radicals}

\lessonobjective{
Use a conjugate to remove a radical difference that creates \(0/0\), then evaluate the simplified limit.
}

A conjugate changes the sign between two terms:
\[
  \sqrt{A}-B
  \qquad\longleftrightarrow\qquad
  \sqrt{A}+B.
\]
Multiplying conjugates uses the difference-of-squares identity:
\[
  (u-v)(u+v)=u^2-v^2.
\]
This removes the radical difference.

\begin{bgguided}[title={One conjugate, every line shown}]
Evaluate
\[
  \lim_{x\to0}\frac{\sqrt{x+4}-2}{x}.
\]

\begin{bgsolution}
Direct substitution gives
\[
  \frac{\sqrt4-2}{0}=\frac00.
\]
Multiply the fraction by \(1\) written as the conjugate over itself:
\[
\frac{\sqrt{x+4}-2}{x}
\cdot
\frac{\sqrt{x+4}+2}{\sqrt{x+4}+2}.
\]
Multiply the numerator using difference of squares:
\[
  (\sqrt{x+4})^2-2^2=x+4-4=x.
\]
Therefore,
\[
\frac{\sqrt{x+4}-2}{x}
=\frac{x}{x(\sqrt{x+4}+2)}.
\]
For \(x\ne0\), cancel \(x\):
\[
  \frac{1}{\sqrt{x+4}+2}.
\]
Now substitute \(x=0\):
\[
  \frac1{\sqrt4+2}=\frac1{2+2}=\boxed{\frac14}.
\]
\end{bgsolution}
\end{bgguided}


% --- BEGIN INLINE checks/conjugate-flow-01 ---
\begin{bgcheck}{conjugate-flow-01}{rational}{1/4}[Check your understanding]
Evaluate \(\lim_{x\to0}\frac{\sqrt{x+4}-2}{x}\).
\begin{bghint}
Use the conjugate \(\sqrt{x+4}+2\).
\end{bghint}
\begin{bgreveal}
The expression simplifies to \(1/(\sqrt{x+4}+2)\), giving \(1/4\).
\end{bgreveal}
\end{bgcheck}

% --- END INLINE checks/conjugate-flow-01 ---


\begin{bgexample}[title={Radical in the denominator}]
Evaluate
\[
  \lim_{x\to9}\frac{x-9}{\sqrt{x}-3}.
\]

\begin{bgsolution}
Substitution gives \(0/0\). Notice that
\[
  x-9=(\sqrt{x}-3)(\sqrt{x}+3).
\]
Therefore, for \(x\ne9\),
\[
  \frac{x-9}{\sqrt{x}-3}=\sqrt{x}+3.
\]
Now substitute:
\[
  \sqrt9+3=3+3=\boxed{6}.
\]
You could also multiply by the conjugate. Recognizing the hidden difference of squares is simply faster.
\end{bgsolution}
\end{bgexample}

\begin{bgexample}[title={Exam-level: two radical terms}]
Evaluate
\[
  \lim_{x\to0}\frac{\sqrt{1+3x}-\sqrt{1-x}}{x}.
\]

\begin{bgsolution}
The numerator is a difference of radicals, so use its conjugate:
\begin{align*}
&\frac{\sqrt{1+3x}-\sqrt{1-x}}{x}
\cdot
\frac{\sqrt{1+3x}+\sqrt{1-x}}{\sqrt{1+3x}+\sqrt{1-x}}\\
&=\frac{(1+3x)-(1-x)}{x\left(\sqrt{1+3x}+\sqrt{1-x}\right)}\\
&=\frac{4x}{x\left(\sqrt{1+3x}+\sqrt{1-x}\right)}\\
&=\frac4{\sqrt{1+3x}+\sqrt{1-x}},\qquad x\ne0.
\end{align*}
Now substitute:
\[
  \frac4{1+1}=\boxed{2}.
\]
\end{bgsolution}
\end{bgexample}
\webnext{calculus/limits/complex-fractions}{Limits With Complex Fractions}
\end{webpage}

\begin{webpage}{calculus/limits/complex-fractions}{Limits With Complex Fractions}{lesson}
\section{Complex Fractions}

\lessonobjective{
Combine smaller fractions using a common denominator before simplifying a limit whose numerator or denominator contains fractions.
}

\begin{bgconcept}
A complex fraction is just a large fraction containing smaller fractions. Do not try to ``cancel through'' addition. First turn the top or bottom into one ordinary fraction. Then divide.
\end{bgconcept}

\begin{bgguided}[title={Combine the numerator first}]
Evaluate
\[
  \lim_{x\to0}\frac{\frac1{x+1}-1}{x}.
\]

\begin{bgsolution}
Direct substitution gives \(0/0\). Combine the two terms in the numerator:
\[
  \frac1{x+1}-1
  =\frac1{x+1}-\frac{x+1}{x+1}
  =\frac{1-(x+1)}{x+1}
  =\frac{-x}{x+1}.
\]
Now place that result over \(x\):
\[
  \frac{\frac{-x}{x+1}}{x}
  =\frac{-x}{x+1}\cdot\frac1x
  =-\frac1{x+1},\qquad x\ne0.
\]
Substitute:
\[
  -\frac1{0+1}=\boxed{-1}.
\]
\end{bgsolution}
\end{bgguided}

\begin{bgexample}[title={Difference quotient for a reciprocal}]
Evaluate
\[
  \lim_{h\to0}
  \frac{\frac1{3+h}-\frac13}{h}.
\]

\begin{bgsolution}
Combine the fractions in the numerator using denominator \(3(3+h)\):
\begin{align*}
\frac1{3+h}-\frac13
&=\frac3{3(3+h)}-\frac{3+h}{3(3+h)}\\
&=\frac{3-(3+h)}{3(3+h)}\\
&=\frac{-h}{3(3+h)}.
\end{align*}
Now divide by \(h\):
\[
  \frac{\frac{-h}{3(3+h)}}{h}
  =\frac{-h}{3(3+h)}\cdot\frac1h
  =-\frac1{3(3+h)}.
\]
Take the limit:
\[
  -\frac1{3(3)}=\boxed{-\frac19}.
\]
\end{bgsolution}
\end{bgexample}

\begin{bgexample}[title={Exam-level: two rational terms}]
Evaluate
\[
  \lim_{x\to2}
  \frac{\frac1x-\frac12}{x-2}.
\]

\begin{bgsolution}
Combine the numerator:
\[
  \frac1x-\frac12=\frac{2-x}{2x}=-\frac{x-2}{2x}.
\]
Therefore,
\[
  \frac{-\frac{x-2}{2x}}{x-2}
  =-\frac1{2x},\qquad x\ne2.
\]
Now substitute:
\[
  -\frac1{2(2)}=\boxed{-\frac14}.
\]
\end{bgsolution}
\end{bgexample}
\webnext{calculus/limits/absolute-value-piecewise}{Limits With Absolute Values and Piecewise Functions}
\end{webpage}

\begin{webpage}{calculus/limits/absolute-value-piecewise}{Limits With Absolute Values and Piecewise Functions}{lesson}
\section{Absolute Values and Piecewise Rules}

\lessonobjective{
Rewrite absolute values on each side of a target and evaluate one-sided limits before deciding whether a two-sided limit exists.
}

The absolute value is piecewise:
\[
  |u|=
  \begin{cases}
    u,&u\ge0,\\
    -u,&u<0.
  \end{cases}
\]
The sign of the inside expression determines which rule applies.

\begin{bgguided}[title={\(|x|/x\)}]
Evaluate
\[
  \lim_{x\to0}\frac{|x|}{x}.
\]

\begin{bgsolution}
For \(x>0\), \(|x|=x\), so
\[
  \frac{|x|}{x}=\frac{x}{x}=1.
\]
Therefore,
\[
  \lim_{x\to0^+}\frac{|x|}{x}=1.
\]
For \(x<0\), \(|x|=-x\), so
\[
  \frac{|x|}{x}=\frac{-x}{x}=-1.
\]
Therefore,
\[
  \lim_{x\to0^-}\frac{|x|}{x}=-1.
\]
The sides disagree, so
\[
  \boxed{\lim_{x\to0}\frac{|x|}{x}=\DNE}.
\]
\end{bgsolution}
\end{bgguided}

\begin{bgexample}[title={A shifted absolute value}]
Evaluate
\[
  \lim_{x\to3}\frac{|x-3|}{x-3}.
\]

\begin{bgsolution}
When \(x>3\), the inside \(x-3\) is positive, so
\[
  |x-3|=x-3
\]
and the quotient equals \(1\).

When \(x<3\), the inside \(x-3\) is negative, so
\[
  |x-3|=-(x-3)
\]
and the quotient equals \(-1\).

Thus,
\[
  \lim_{x\to3^-}\frac{|x-3|}{x-3}=-1,
\]
\[
  \lim_{x\to3^+}\frac{|x-3|}{x-3}=1.
\]
Therefore the two-sided limit does not exist.
\end{bgsolution}
\end{bgexample}

\begin{bgexample}[title={An absolute value whose limit does exist}]
Evaluate
\[
  \lim_{x\to2}\frac{|x-2|}{\sqrt{|x-2|}}.
\]

\begin{bgsolution}
For \(x\ne2\), \(|x-2|>0\), and
\[
  \frac{|x-2|}{\sqrt{|x-2|}}=\sqrt{|x-2|}.
\]
As \(x\to2\), the distance \(|x-2|\to0\). Therefore,
\[
  \sqrt{|x-2|}\to0,
\]
so
\[
  \boxed{0}.
\]
Absolute value does not automatically cause a limit to fail. It causes one-sided analysis when the formula changes sign or form across the target.
\end{bgsolution}
\end{bgexample}
\webnext{calculus/limits/finite-limit-decision-tree}{Finite Limit Decision Tree}
\end{webpage}

\begin{webpage}{calculus/limits/finite-limit-decision-tree}{Finite Limit Decision Tree}{reference}
\section{The Finite-Limit Decision Tree}

\begin{bgmethod}[title={A Repeatable Homework Procedure}]
For a finite target \(x\to a\):
\begin{enumerate}
  \item Substitute \(x=a\).
  \item If you get a real number, stop.
  \item If you get \(0/0\), inspect the expression:
  \begin{itemize}
    \item factor polynomial expressions;
    \item rationalize radical differences;
    \item combine complex fractions;
    \item split absolute values or piecewise functions into one-sided cases;
    \item save trig limits for Section 3.
  \end{itemize}
  \item Simplify only for nearby values where the original expression is defined.
  \item Substitute again.
  \item State why the method is valid, especially if one-sided limits are involved.
\end{enumerate}
\end{bgmethod}
\webnext{calculus/limits/squeeze-theorem}{The Squeeze Theorem}
\end{webpage}

\begin{webpage}{calculus/limits/finite-limits-review}{Finite Limits Review}{review}
\section*{Section 2 Summary}
\addcontentsline{toc}{section}{Section 2 Summary}

\begin{bgsummary}
\begin{itemize}
  \item Direct substitution is the correct first move, not an unsophisticated shortcut.
  \item A legal real substitution result usually finishes the problem.
  \item \(0/0\) is indeterminate; it tells you to transform the expression.
  \item Factor polynomial forms, rationalize radicals, combine complex fractions, and use one-sided rules for absolute values.
  \item Cancellation is legal in a limit when it is performed for nearby non-target inputs where the cancelled factor is nonzero.
\end{itemize}
\end{bgsummary}
\webnext{calculus/limits/squeeze-theorem}{The Squeeze Theorem}
\end{webpage}

\begin{webpage}{calculus/limits/finite-limits-concept-quiz}{Finite Limits Concept Quiz}{quiz}
\section*{Section 2 Concept Quiz}
\addcontentsline{toc}{section}{Section 2 Concept Quiz}

Use this as a short retrieval check after finishing the section. On the website, each item should accept an answer before revealing the hint and worked explanation.

\begin{bgcheck}{substitution-q1}{integer}{11}[Concept check]
Evaluate \(\lim_{x\to3}(2x+5)\).
\begin{bghint}
This linear function is continuous, so substitute \(x=3\).
\end{bghint}
\begin{bgreveal}
\(2(3)+5=11\).
\end{bgreveal}
\end{bgcheck}

\begin{bgcheck}{factor-q1}{integer}{8}[Concept check]
Evaluate \(\lim_{x\to4}\frac{x^2-16}{x-4}\).
\begin{bghint}
Factor \(x^2-16=(x-4)(x+4)\).
\end{bghint}
\begin{bgreveal}
Cancel for \(x\ne4\), leaving \(x+4\), which approaches \(8\).
\end{bgreveal}
\end{bgcheck}

\begin{bgcheck}{radical-q1}{rational}{1/6}[Concept check]
Evaluate \(\lim_{x\to0}\frac{\sqrt{x+9}-3}{x}\).
\begin{bghint}
Multiply by the conjugate \(\sqrt{x+9}+3\).
\end{bghint}
\begin{bgreveal}
The expression simplifies to \(1/(\sqrt{x+9}+3)\), so the limit is \(1/6\).
\end{bgreveal}
\end{bgcheck}

\begin{bgcheck}{abs-q1}{choice}{DNE}[Concept check]
Does \(\lim_{x\to0}|x|/x\) exist? Enter DNE if it does not.
\begin{bghint}
Check the left and right sides separately.
\end{bghint}
\begin{bgreveal}
The left-hand limit is \(-1\) and the right-hand limit is \(1\), so the limit is DNE.
\end{bgreveal}
\end{bgcheck}

\begin{bgsummary}[title={Quiz use on the website}]
This page should report completion by concept, not merely a total score. A wrong answer should link back to the exact lesson that teaches the missing method. The same questions may appear in randomized numerical variants, but the canonical explanation should remain stable.
\end{bgsummary}
\webnext{calculus/limits/squeeze-theorem}{The Squeeze Theorem}
\end{webpage}

\begin{webpage}{calculus/limits/finite-limits-practice}{Finite Limit Practice Problems}{practice}
\section*{Section 2 Exercises}
\addcontentsline{toc}{section}{Section 2 Exercises}

\subsection*{A. Direct substitution and limit laws}
\begin{bgexercises}[label=\textbf{2.\arabic*}]
  \item \(\displaystyle\lim_{x\to3}(2x+7)\)
  \item \(\displaystyle\lim_{x\to-2}(x^2-4x+1)\)
  \item \(\displaystyle\lim_{t\to4}(t^3-2t)\)
  \item \(\displaystyle\lim_{x\to1}\frac{x^2+3}{x+2}\)
  \item \(\displaystyle\lim_{x\to-1}\frac{2x^2+x+5}{x^2+4}\)
  \item \(\displaystyle\lim_{u\to8}\sqrt[3]{u+19}\)
  \item \(\displaystyle\lim_{x\to0}(3\cos x+2\sin x)\)
  \item If \(f(x)\to2\) and \(g(x)\to-3\), find \(\lim[4f(x)-g(x)]\).
  \item With the same limits, find \(\lim f(x)g(x)\).
  \item With the same limits, find \(\lim\dfrac{f(x)^2+g(x)}{f(x)-g(x)}\).
  \item State why the quotient law cannot be used when the denominator's limit is zero.
  \item Give an example where direct substitution produces a real number and therefore no algebraic simplification is needed.
\end{bgexercises}

\subsection*{B. Factoring}
\begin{bgexercises}[label=\textbf{2.\arabic*},resume]
  \item \(\displaystyle\lim_{x\to5}\frac{x^2-25}{x-5}\)
  \item \(\displaystyle\lim_{x\to-3}\frac{x^2+5x+6}{x+3}\)
  \item \(\displaystyle\lim_{x\to4}\frac{x^2-7x+12}{x-4}\)
  \item \(\displaystyle\lim_{x\to2}\frac{x^3-8}{x-2}\)
  \item \(\displaystyle\lim_{x\to-2}\frac{x^3+8}{x+2}\)
  \item \(\displaystyle\lim_{x\to1}\frac{x^4-1}{x-1}\)
  \item \(\displaystyle\lim_{x\to-1}\frac{x^4-1}{x+1}\)
  \item \(\displaystyle\lim_{x\to3}\frac{x^2-9}{x^2-4x+3}\)
  \item \(\displaystyle\lim_{x\to1}\frac{x^3-1}{x^2-1}\)
  \item \(\displaystyle\lim_{h\to0}\frac{(2+h)^2-4}{h}\)
  \item \(\displaystyle\lim_{h\to0}\frac{(3+h)^3-27}{h}\)
  \item Explain why the cancelled expression and original expression may differ at the target but have the same limit.
\end{bgexercises}

\subsection*{C. Radical limits}
\begin{bgexercises}[label=\textbf{2.\arabic*},resume]
  \item \(\displaystyle\lim_{x\to0}\frac{\sqrt{x+9}-3}{x}\)
  \item \(\displaystyle\lim_{x\to0}\frac{\sqrt{x+16}-4}{x}\)
  \item \(\displaystyle\lim_{x\to25}\frac{x-25}{\sqrt{x}-5}\)
  \item \(\displaystyle\lim_{x\to4}\frac{\sqrt{x}-2}{x-4}\)
  \item \(\displaystyle\lim_{x\to0}\frac{\sqrt{1+5x}-1}{x}\)
  \item \(\displaystyle\lim_{x\to0}\frac{\sqrt{4+x}-\sqrt{4-x}}{x}\)
  \item \(\displaystyle\lim_{x\to3}\frac{x-3}{\sqrt{x+6}-3}\)
  \item \(\displaystyle\lim_{x\to0}\frac{\sqrt{1+2x}-\sqrt{1-x}}{x}\)
\end{bgexercises}

\subsection*{D. Complex fractions}
\begin{bgexercises}[label=\textbf{2.\arabic*},resume]
  \item \(\displaystyle\lim_{x\to0}\frac{\frac1{x+1}-1}{x}\)
  \item \(\displaystyle\lim_{x\to0}\frac{\frac1{x+2}-\frac12}{x}\)
  \item \(\displaystyle\lim_{x\to0}\frac{\frac1{3+x}-\frac13}{x}\)
  \item \(\displaystyle\lim_{x\to2}\frac{\frac1x-\frac12}{x-2}\)
  \item \(\displaystyle\lim_{x\to1}\frac{\frac1{x+1}-\frac12}{x-1}\)
  \item \(\displaystyle\lim_{h\to0}\frac{\frac1{a+h}-\frac1a}{h}\), where \(a\ne0\).
  \item \(\displaystyle\lim_{x\to0}\frac{\frac{x}{x+1}}{x}\)
  \item \(\displaystyle\lim_{x\to0}\frac{\frac1{1-x}-1}{x}\)
\end{bgexercises}

\subsection*{E. Absolute values and piecewise forms}
\begin{bgexercises}[label=\textbf{2.\arabic*},resume]
  \item \(\displaystyle\lim_{x\to0^-}\frac{|x|}{x}\)
  \item \(\displaystyle\lim_{x\to0^+}\frac{|x|}{x}\)
  \item \(\displaystyle\lim_{x\to0}\frac{|x|}{x}\)
  \item \(\displaystyle\lim_{x\to2}\frac{|x-2|}{x-2}\)
  \item \(\displaystyle\lim_{x\to-1^-}\frac{|x+1|}{x+1}\)
  \item \(\displaystyle\lim_{x\to3}\frac{|x-3|}{\sqrt{|x-3|}}\)
  \item Let \(f(x)=\begin{cases}x+2,&x<1,\\x^2+1,&x\ge1.\end{cases}\) Find the limit at \(1\).
  \item Let \(g(x)=\begin{cases}2x,&x<2,\\x+3,&x\ge2.\end{cases}\) Find both one-sided limits at \(2\).
\end{bgexercises}

\subsection*{F. Mixed exam practice}
\begin{bgexercises}[label=\textbf{2.\arabic*},resume]
  \item \(\displaystyle\lim_{x\to2}\frac{x^3-4x}{x^2-4}\)
  \item \(\displaystyle\lim_{x\to1}\frac{x^5-1}{x^2-1}\)
  \item \(\displaystyle\lim_{x\to0}\frac{\sqrt{9+2x}-3}{x}\)
  \item \(\displaystyle\lim_{x\to4}\frac{\frac1x-\frac14}{x-4}\)
  \item \(\displaystyle\lim_{x\to0}\frac{|x|}{\sqrt{|x|}}\)
  \item A student cancels \(x\) in \((x^2+x)/x\) and obtains \(x+1\). Explain why this cancellation is valid for \(x\ne0\), and distinguish it from cancelling a term across addition without factoring.
  \item Design a \(0/0\) limit whose value is \(5\).
  \item Design a \(0/0\) limit whose value is \(0\).
  \item Design a \(0/0\) limit that does not exist.
  \item Write a short decision procedure for determining whether to factor, rationalize, combine fractions, or split into one-sided cases.
\end{bgexercises}

Answers begin in \cref{app:chapter-answers}.
\webnext{calculus/limits/squeeze-theorem}{The Squeeze Theorem}
\end{webpage}


% --- END INLINE chapters/ch02_finite ---


% --- BEGIN INLINE chapters/ch03_squeeze_trig ---
\begin{webpage}{calculus/limits/squeeze-theorem}{The Squeeze Theorem}{lesson}
\chapter{The Squeeze Theorem and Trigonometric Limits}

\section{When Direct Algebra Is Not Enough}

\lessonobjective{
Use upper and lower bounds to determine a limit even when the middle function oscillates or is difficult to evaluate directly.
}

Some functions refuse to simplify into a friendly formula. The Squeeze Theorem handles them by trapping their outputs between two simpler functions.

\begin{bgconcept}
Suppose three people walk through a doorway side by side. The person in the middle cannot end up above the person on the top or below the person on the bottom. If the top and bottom people are forced toward the same height, the middle person is forced there too. That is the Squeeze Theorem.
\end{bgconcept}

\begin{bgtheorem}[title={Squeeze Theorem}]
Suppose that for all \(x\) sufficiently close to \(a\),
\[
  g(x)\le f(x)\le h(x),
\]
and suppose
\[
  \lim_{x\to a}g(x)=L
  \qquad\text{and}\qquad
  \lim_{x\to a}h(x)=L.
\]
Then
\[
  \boxed{\lim_{x\to a}f(x)=L}.
\]
\end{bgtheorem}

\begin{bgguided}[title={A number trapped between shrinking walls}]
Suppose \(f(x)\) satisfies
\[
  -x^2\le f(x)\le x^2
\]
near \(x=0\). Find \(\lim_{x\to0}f(x)\).

\begin{bgsolution}
The lower wall approaches
\[
  \lim_{x\to0}(-x^2)=0.
\]
The upper wall approaches
\[
  \lim_{x\to0}x^2=0.
\]
Because \(f(x)\) is trapped between two functions that both approach \(0\),
\[
  \boxed{\lim_{x\to0}f(x)=0}.
\]
We do not need a formula for \(f\). The trap is enough.
\end{bgsolution}
\end{bgguided}


% --- BEGIN INLINE checks/squeeze-flow-01 ---
\begin{bgcheck}{squeeze-flow-01}{integer}{0}[Check your understanding]
Evaluate \(\lim_{x\to0}x^2\sin(1/x)\).
\begin{bghint}
Use \(|\sin(1/x)|\le1\).
\end{bghint}
\begin{bgreveal}
The absolute value is at most \(x^2\), which tends to \(0\).
\end{bgreveal}
\end{bgcheck}

% --- END INLINE checks/squeeze-flow-01 ---


\subsection{Bounded oscillation times a shrinking factor}

The inequalities
\[
  -1\le\sin u\le1
\]
and
\[
  -1\le\cos u\le1
\]
hold for every real \(u\). If a bounded trigonometric factor is multiplied by something that approaches zero, the product is often squeezed to zero.

\begin{bgexample}[title={An oscillating function forced to zero}]
Evaluate
\[
  \lim_{x\to0}x^2\sin\left(\frac1x\right).
\]

\begin{bgsolution}
Start with the universal sine bound:
\[
  -1\le\sin\left(\frac1x\right)\le1.
\]
Since \(x^2\ge0\), multiply every part by \(x^2\) without reversing the inequalities:
\[
  -x^2\le x^2\sin\left(\frac1x\right)\le x^2.
\]
Now take limits of the outer functions:
\[
  \lim_{x\to0}(-x^2)=0,
  \qquad
  \lim_{x\to0}x^2=0.
\]
Therefore,
\[
  \boxed{\lim_{x\to0}x^2\sin\left(\frac1x\right)=0}.
\]
The sine factor continues oscillating. The shrinking factor \(x^2\) reduces the size of every oscillation until the graph is trapped near zero.
\end{bgsolution}
\end{bgexample}

\begin{figure}[htbp]
\centering
\begin{tikzpicture}
\begin{axis}[
  width=0.88\textwidth,height=0.46\textwidth,
  xmin=-0.35,xmax=0.35,ymin=-0.13,ymax=0.13,
  axis lines=middle,xlabel={$x$},ylabel={$y$},
  grid=major,major grid style={draw=BGLine},
  legend style={draw=none,fill=none,font=\small,at={(0.98,0.97)},anchor=north east}
]
\addplot[BGBlue,thick,domain=-0.35:-0.008,samples=140]{x^2*sin(deg(1/x))};
\addplot[BGBlue,thick,domain=0.008:0.35,samples=140]{x^2*sin(deg(1/x))};
\addlegendentry{$x^2\sin(1/x)$}
\addplot[BGGreen,dashed,thick,domain=-0.35:0.35,samples=120]{x^2};
\addlegendentry{$x^2$}
\addplot[BGRed,dashed,thick,domain=-0.35:0.35,samples=120]{-x^2};
\addlegendentry{$-x^2$}
\end{axis}
\end{tikzpicture}
\caption{The oscillating function is trapped between \(-x^2\) and \(x^2\), both of which approach zero.}
\label{fig:squeeze-oscillation}
\end{figure}

\begin{bggraphspec}
\textbf{Functions:} \(y=x^2\sin(1/x)\), \(y=x^2\), and \(y=-x^2\). \textbf{Window:} \(-0.35\le x\le0.35\), \(-0.13\le y\le0.13\). Use at least 1500 samples on each side of zero for the oscillating curve. Do not connect across \(x=0\), where the displayed formula is undefined.
\end{bggraphspec}

\begin{bgexample}[title={Absolute-value squeeze}]
Evaluate
\[
  \lim_{x\to0}x\cos\left(\frac{5}{x}\right).
\]

\begin{bgsolution}
Because \(|\cos u|\le1\),
\[
  \left|x\cos\left(\frac5x\right)\right|
  \le |x|.
\]
This means
\[
  -|x|\le x\cos\left(\frac5x\right)\le|x|.
\]
Both outer functions approach zero, so
\[
  \boxed{\lim_{x\to0}x\cos\left(\frac5x\right)=0}.
\]
\end{bgsolution}
\end{bgexample}

\begin{bghomework}
When a function contains \(\sin(1/x)\), \(\cos(1/x)\), or another bounded oscillating factor, ask whether the rest of the expression approaches zero. The inequality
\[
  |\text{bounded factor}\times\text{shrinking factor}|
  \le C|\text{shrinking factor}|
\]
is often the entire argument.
\end{bghomework}
\webnext{calculus/limits/sin-x-over-x}{Why sin(x)/x Approaches 1}
\end{webpage}

\begin{webpage}{calculus/limits/sin-x-over-x}{Why sin(x)/x Approaches 1}{lesson}
\section{The Fundamental Trigonometric Limit}

\lessonobjective{
Understand geometrically why \(\lim_{x\to0}\sin x/x=1\) when angles are measured in radians, and use it to evaluate scaled trigonometric limits.
}

The most important trigonometric limit in first-semester calculus is
\[
  \boxed{\lim_{x\to0}\frac{\sin x}{x}=1}.
\]
It is not a random formula. It follows from the geometry of the unit circle and the Squeeze Theorem.

\begin{bgmistake}
The formula requires radians. If \(x\) is measured in degrees, \(\sin x/x\) approaches \(\pi/180\), not \(1\). Calculus uses radians because arc length on the unit circle equals the angle measure itself.
\end{bgmistake}
\webnext{calculus/limits/sin-x-over-x-proof}{Geometric Proof of the sin(x)/x Limit}
\end{webpage}

\begin{webpage}{calculus/limits/sin-x-over-x-proof}{Geometric Proof of the sin(x)/x Limit}{lesson}
\subsection{The unit-circle comparison}

Take an angle \(x\) with
\[
  0<x<\frac{\pi}{2}.
\]
In a unit circle:

\begin{itemize}
  \item the area of the inner triangle is \(\frac12\sin x\cos x\);
  \item the area of the circular sector is \(\frac12x\);
  \item the area of the outer tangent triangle is \(\frac12\tan x\).
\end{itemize}

Therefore,
\[
  \frac12\sin x\cos x
  <\frac12x
  <\frac12\tan x.
\]
Multiply by \(2\):
\[
  \sin x\cos x<x<\tan x.
\]
Since \(\sin x>0\), divide by \(\sin x\):
\[
  \cos x<\frac{x}{\sin x}<\frac1{\cos x}.
\]
Taking positive reciprocals reverses the inequalities:
\[
  \cos x<\frac{\sin x}{x}<1.
\]
As \(x\to0^+\), both outer expressions approach \(1\). Hence
\[
  \lim_{x\to0^+}\frac{\sin x}{x}=1.
\]
The function \(\sin x/x\) is even because
\[
  \frac{\sin(-x)}{-x}=\frac{\sin x}{x}.
\]
Thus the left-hand limit is also \(1\), and the two-sided limit is \(1\).

\begin{figure}[htbp]
\centering
\begin{tikzpicture}[scale=3.2]
\coordinate (O) at (0,0);
\coordinate (A) at (1,0);
\coordinate (B) at ({cos(32)},{sin(32)});
\coordinate (T) at (1,{tan(32)});
\draw[BGInk,thick] (O) circle (1);
\draw[BGBlue,very thick] (O)--(A);
\draw[BGBlue,very thick] (O)--(B);
\draw[BGGreen,thick] (A)--(B);
\draw[BGRed,thick] (A)--(T)--(O);
\draw pic[draw=BGGold,fill=BGGold!18,angle radius=0.27cm,"$x$"{font=\small}] {angle=A--O--B};
\fill[BGInk] (O) circle (0.015) (A) circle (0.015) (B) circle (0.015) (T) circle (0.015);
\node[below left,font=\small] at (O) {$O$};
\node[below right,font=\small] at (A) {$A$};
\node[above,font=\small] at (B) {$B$};
\node[right,font=\small] at (T) {$T$};
\node[font=\small,BGGreen] at (0.78,0.23) {inner triangle};
\node[font=\small,BGRed,rotate=25] at (0.61,0.62) {outer triangle};
\end{tikzpicture}
\caption{For a unit circle and \(0<x<\pi/2\), the inner triangle lies inside the sector, which lies inside the outer tangent triangle.}
\label{fig:unit-circle-squeeze}
\end{figure}

\begin{bgconcept}
Near zero, the arc length \(x\) and the vertical height \(\sin x\) become nearly indistinguishable. Their ratio approaches \(1\). The unit-circle proof turns ``nearly'' into inequalities strong enough for the Squeeze Theorem.
\end{bgconcept}

\begin{figure}[htbp]
\centering
\begin{tikzpicture}
\begin{axis}[
  width=0.84\textwidth,height=0.46\textwidth,
  xmin=-6.3,xmax=6.3,ymin=-0.35,ymax=1.18,
  axis lines=middle,xlabel={$x$},ylabel={$\sin x/x$},
  grid=major,major grid style={draw=BGLine},
  xtick={-6.283,-3.142,0,3.142,6.283},
  xticklabels={$-2\pi$,$-\pi$,$0$,$\pi$,$2\pi$}
]
\addplot[BGBlue,very thick,domain=-6.3:-0.035,samples=120]{sin(deg(x))/x};
\addplot[BGBlue,very thick,domain=0.035:6.3,samples=120]{sin(deg(x))/x};
\addplot[only marks,mark=o,mark size=4pt,very thick,BGRed] coordinates {(0,1)};
\draw[dashed,BGGray] (axis cs:-6.3,1)--(axis cs:6.3,1);
\end{axis}
\end{tikzpicture}
\caption{The graph of \(\sin x/x\) has a removable hole at \((0,1)\) and approaches \(1\) from both sides.}
\label{fig:sinc}
\end{figure}

\begin{bgguided}[title={Use the exact pattern}]
Evaluate
\[
  \lim_{x\to0}\frac{\sin x}{x}.
\]

\begin{bgsolution}
This is the fundamental trigonometric limit itself:
\[
  \boxed{1}.
\]
No algebra is needed. Recognizing the pattern is the method.
\end{bgsolution}
\end{bgguided}


% --- BEGIN INLINE checks/trig-flow-01 ---
\begin{bgcheck}{trig-flow-01}{integer}{4}[Check your understanding]
Evaluate \(\lim_{x\to0}\frac{\sin(4x)}x\).
\begin{bghint}
Multiply and divide by \(4\).
\end{bghint}
\begin{bgreveal}
The limit is \(4\).
\end{bgreveal}
\end{bgcheck}

% --- END INLINE checks/trig-flow-01 ---

\webnext{calculus/limits/scaled-sine-tangent}{Scaled Sine and Tangent Limits}
\end{webpage}

\begin{webpage}{calculus/limits/scaled-sine-tangent}{Scaled Sine and Tangent Limits}{lesson}
\section{Scaled Sine and Tangent Limits}

\lessonobjective{
Rewrite trigonometric limits so that a factor has the standard form \(\sin u/u\) or \(\tan u/u\).
}

\subsection{The substitution pattern}

If \(u=kx\), then \(u\to0\) whenever \(x\to0\). We want to create the denominator \(kx\) that matches the sine argument.

\begin{bgexample}[title={\(\sin(5x)/x\)}]
Evaluate
\[
  \lim_{x\to0}\frac{\sin(5x)}{x}.
\]

\begin{bgsolution}
The sine argument is \(5x\), but the denominator is only \(x\). Multiply and divide by \(5\):
\[
  \frac{\sin(5x)}{x}
  =5\frac{\sin(5x)}{5x}.
\]
As \(x\to0\), \(5x\to0\), so
\[
  \frac{\sin(5x)}{5x}\to1.
\]
Therefore,
\[
  \boxed{5}.
\]
\end{bgsolution}
\end{bgexample}

\begin{bgguided}[title={Keep the outside constant}]
Evaluate
\[
  \lim_{x\to0}\frac{\sin(3x)}{2x}.
\]

\begin{bgsolution}
Create \(3x\) in the denominator:
\[
  \frac{\sin(3x)}{2x}
  =\frac32\frac{\sin(3x)}{3x}.
\]
The standard factor approaches \(1\), so
\[
  \boxed{\frac32}.
\]
\end{bgsolution}
\end{bgguided}
\webnext{calculus/limits/ratio-of-sines}{Limits of Ratios of Sine Functions}
\end{webpage}

\begin{webpage}{calculus/limits/ratio-of-sines}{Limits of Ratios of Sine Functions}{lesson}
\subsection{Ratio of two sine expressions}

\begin{bgexample}[title={Sine over sine}]
Evaluate
\[
  \lim_{x\to0}\frac{\sin(4x)}{\sin(7x)}.
\]

\begin{bgsolution}
Insert factors that create both standard limits:
\begin{align*}
\frac{\sin(4x)}{\sin(7x)}
&=\frac{\sin(4x)}{4x}
  \cdot\frac{4x}{7x}
  \cdot\frac{7x}{\sin(7x)}\\
&=\frac{\sin(4x)}{4x}
  \cdot\frac47
  \cdot\frac{7x}{\sin(7x)}.
\end{align*}
The first and third factors approach \(1\), so
\[
  \boxed{\frac47}.
\]
\end{bgsolution}
\end{bgexample}

\subsection{Tangent}

Since
\[
  \frac{\tan x}{x}
  =\frac{\sin x}{x\cos x},
\]
we have
\[
  \boxed{\lim_{x\to0}\frac{\tan x}{x}=1}.
\]

\begin{bgexample}[title={A tangent limit}]
Evaluate
\[
  \lim_{x\to0}\frac{\tan(6x)}{2x}.
\]

\begin{bgsolution}
Create \(6x\) in the denominator:
\[
  \frac{\tan(6x)}{2x}
  =3\frac{\tan(6x)}{6x}.
\]
The standard tangent factor approaches \(1\), so
\[
  \boxed{3}.
\]
\end{bgsolution}
\end{bgexample}
\webnext{calculus/limits/cosine-limits}{Cosine Limits and Trigonometric Identities}
\end{webpage}

\begin{webpage}{calculus/limits/cosine-limits}{Cosine Limits and Trigonometric Identities}{lesson}
\section{Cosine Limits and Trigonometric Identities}

\lessonobjective{
Use conjugates and identities to reduce limits involving \(1-\cos x\), \(\sin^2x\), or other trigonometric expressions to standard forms.
}

\subsection{A first cosine limit}

\[
  \boxed{\lim_{x\to0}\frac{1-\cos x}{x}=0}.
\]
To show this, rationalize with \(1+\cos x\):
\begin{align*}
\frac{1-\cos x}{x}
&\cdot\frac{1+\cos x}{1+\cos x}\\
&=\frac{1-\cos^2x}{x(1+\cos x)}\\
&=\frac{\sin^2x}{x(1+\cos x)}\\
&=\frac{\sin x}{x}\cdot\frac{\sin x}{1+\cos x}.
\end{align*}
The first factor approaches \(1\), and the second approaches \(0/2=0\). Therefore the product approaches \(0\).

\subsection{The second important cosine limit}

\[
  \boxed{\lim_{x\to0}\frac{1-\cos x}{x^2}=\frac12}.
\]
Again rationalize:
\begin{align*}
\frac{1-\cos x}{x^2}
&=\frac{\sin^2x}{x^2(1+\cos x)}\\
&=\left(\frac{\sin x}{x}\right)^2\frac1{1+\cos x}.
\end{align*}
Taking limits gives
\[
  1^2\cdot\frac12=\frac12.
\]

\begin{bgguided}[title={Recognize the squared pattern}]
Evaluate
\[
  \lim_{x\to0}\frac{\sin^2x}{x^2}.
\]

\begin{bgsolution}
Rewrite as a square:
\[
  \frac{\sin^2x}{x^2}
  =\left(\frac{\sin x}{x}\right)^2.
\]
The inside approaches \(1\), so the square approaches
\[
  \boxed{1}.
\]
\end{bgsolution}
\end{bgguided}

\begin{bgexample}[title={Identity plus a standard limit}]
Evaluate
\[
  \lim_{x\to0}\frac{1-\cos(3x)}{x^2}.
\]

\begin{bgsolution}
Create \((3x)^2\) in the denominator:
\[
  \frac{1-\cos(3x)}{x^2}
  =9\frac{1-\cos(3x)}{(3x)^2}.
\]
As \(x\to0\), \(3x\to0\), so the standard cosine limit gives
\[
  9\cdot\frac12=\boxed{\frac92}.
\]
\end{bgsolution}
\end{bgexample}

\begin{bgexample}[title={Exam-level: use an identity before the standard limit}]
Evaluate
\[
  \lim_{x\to0}\frac{\sin(2x)\sin(5x)}{x^2}.
\]

\begin{bgsolution}
Separate the two standard factors:
\[
\frac{\sin(2x)\sin(5x)}{x^2}
=\left(\frac{\sin(2x)}{x}\right)
 \left(\frac{\sin(5x)}{x}\right).
\]
Each factor can be rewritten:
\[
  \frac{\sin(2x)}{x}=2\frac{\sin(2x)}{2x}\to2,
\]
\[
  \frac{\sin(5x)}{x}=5\frac{\sin(5x)}{5x}\to5.
\]
Therefore the product approaches
\[
  \boxed{10}.
\]
\end{bgsolution}
\end{bgexample}

\begin{bgexamnote}
Do not replace \(\sin x\) by \(x\) as an algebraic identity. The statement \(\sin x\sim x\) near zero means their ratio approaches \(1\); it does not mean the two expressions are equal for nonzero \(x\). Preserve the limit argument.
\end{bgexamnote}
\webnext{calculus/limits/trig-limit-decision-guide}{Trigonometric Limit Decision Guide}
\end{webpage}

\begin{webpage}{calculus/limits/trig-limit-decision-guide}{Trigonometric Limit Decision Guide}{reference}
\section{Trig-Limit Decision Guide}

\begin{bgmethod}
For a trigonometric limit near zero:
\begin{enumerate}
  \item Substitute first.
  \item If the result is \(0/0\), look for \(\sin u/u\), \(\tan u/u\), or \((1-\cos u)/u^2\).
  \item Multiply by constants to make the denominator match the trigonometric argument.
  \item Use identities such as \(1-\cos^2u=\sin^2u\) when needed.
  \item If a bounded trig factor is multiplied by something tending to zero, consider the Squeeze Theorem.
\end{enumerate}
\end{bgmethod}
\webnext{calculus/limits/infinite-limits}{Infinite Limits Explained}
\end{webpage}

\begin{webpage}{calculus/limits/trig-limits-review}{Squeeze and Trigonometric Limits Review}{review}
\section*{Section 3 Summary}
\addcontentsline{toc}{section}{Section 3 Summary}

\begin{bgsummary}
\begin{itemize}
  \item If \(g\le f\le h\) and both outer limits equal \(L\), then the middle limit is \(L\).
  \item Bounded oscillation multiplied by a shrinking factor is often squeezed to zero.
  \item In radians, \(\lim_{x\to0}\sin x/x=1\) and \(\lim_{x\to0}\tan x/x=1\).
  \item Match the denominator to the trigonometric argument.
  \item \(\lim_{x\to0}(1-\cos x)/x^2=1/2\).
\end{itemize}
\end{bgsummary}
\webnext{calculus/limits/infinite-limits}{Infinite Limits Explained}
\end{webpage}

\begin{webpage}{calculus/limits/trig-limits-concept-quiz}{Trigonometric Limits Concept Quiz}{quiz}
\section*{Section 3 Concept Quiz}
\addcontentsline{toc}{section}{Section 3 Concept Quiz}

Use this as a short retrieval check after finishing the section. On the website, each item should accept an answer before revealing the hint and worked explanation.

\begin{bgcheck}{squeeze-q1}{integer}{0}[Concept check]
If \(-x^2\le f(x)\le x^2\), find \(\lim_{x\to0}f(x)\).
\begin{bghint}
Both outer functions approach the same value.
\end{bghint}
\begin{bgreveal}
Both bounds approach \(0\), so the Squeeze Theorem gives \(0\).
\end{bgreveal}
\end{bgcheck}

\begin{bgcheck}{sine-q1}{integer}{5}[Concept check]
Evaluate \(\lim_{x\to0}\frac{\sin(5x)}{x}\).
\begin{bghint}
Make the denominator match the angle: multiply and divide by \(5\).
\end{bghint}
\begin{bgreveal}
\(5\cdot\lim_{x\to0}\frac{\sin(5x)}{5x}=5\).
\end{bgreveal}
\end{bgcheck}

\begin{bgcheck}{sine-ratio-q1}{rational}{3/7}[Concept check]
Evaluate \(\lim_{x\to0}\frac{\sin(3x)}{\sin(7x)}\).
\begin{bghint}
Build one \(\sin u/u\) factor in the numerator and one in the denominator.
\end{bghint}
\begin{bgreveal}
The limit is the ratio of the coefficients, \(3/7\).
\end{bgreveal}
\end{bgcheck}

\begin{bgcheck}{cosine-q1}{rational}{1/2}[Concept check]
Evaluate \(\lim_{x\to0}\frac{1-\cos x}{x^2}\).
\begin{bghint}
Use the standard cosine limit or multiply by \(1+\cos x\).
\end{bghint}
\begin{bgreveal}
The limit is \(1/2\).
\end{bgreveal}
\end{bgcheck}

\begin{bgsummary}[title={Quiz use on the website}]
This page should report completion by concept, not merely a total score. A wrong answer should link back to the exact lesson that teaches the missing method. The same questions may appear in randomized numerical variants, but the canonical explanation should remain stable.
\end{bgsummary}
\webnext{calculus/limits/infinite-limits}{Infinite Limits Explained}
\end{webpage}

\begin{webpage}{calculus/limits/trig-limits-practice}{Trigonometric Limit Practice Problems}{practice}
\section*{Section 3 Exercises}
\addcontentsline{toc}{section}{Section 3 Exercises}

\subsection*{A. Squeeze Theorem}
\begin{bgexercises}[label=\textbf{3.\arabic*}]
  \item If \(-x^4\le f(x)\le x^4\) near zero, find \(\lim_{x\to0}f(x)\).
  \item If \(2-x^2\le g(x)\le2+x^2\) near zero, find the limit.
  \item Evaluate \(\lim_{x\to0}x^3\sin(1/x)\).
  \item Evaluate \(\lim_{x\to0}x^2\cos(7/x)\).
  \item Evaluate \(\lim_{x\to0}|x|\sin(4/x)\).
  \item Evaluate \(\lim_{x\to0}x\cos(1/x^2)\).
  \item Suppose \(|f(x)|\le5|x-2|\). Find \(\lim_{x\to2}f(x)\).
  \item Explain why boundedness of \(f\) alone is not enough to conclude \(\lim_{x\to0}f(x)=0\).
\end{bgexercises}

\subsection*{B. Fundamental sine limits}
\begin{bgexercises}[label=\textbf{3.\arabic*},resume]
  \item \(\displaystyle\lim_{x\to0}\frac{\sin x}{x}\)
  \item \(\displaystyle\lim_{x\to0}\frac{\sin(2x)}{x}\)
  \item \(\displaystyle\lim_{x\to0}\frac{\sin(7x)}{3x}\)
  \item \(\displaystyle\lim_{x\to0}\frac{\sin(5x)}{\sin(2x)}\)
  \item \(\displaystyle\lim_{x\to0}\frac{\sin(3x)}{\sin(8x)}\)
  \item \(\displaystyle\lim_{t\to0}\frac{\sin(\pi t)}{t}\)
  \item \(\displaystyle\lim_{x\to0}\frac{x}{\sin(4x)}\)
  \item \(\displaystyle\lim_{x\to0}\frac{3x}{\sin(5x)}\)
\end{bgexercises}

\subsection*{C. Tangent and cosine limits}
\begin{bgexercises}[label=\textbf{3.\arabic*},resume]
  \item \(\displaystyle\lim_{x\to0}\frac{\tan x}{x}\)
  \item \(\displaystyle\lim_{x\to0}\frac{\tan(4x)}{x}\)
  \item \(\displaystyle\lim_{x\to0}\frac{\tan(3x)}{2x}\)
  \item \(\displaystyle\lim_{x\to0}\frac{1-\cos x}{x}\)
  \item \(\displaystyle\lim_{x\to0}\frac{1-\cos x}{x^2}\)
  \item \(\displaystyle\lim_{x\to0}\frac{1-\cos(4x)}{x^2}\)
  \item \(\displaystyle\lim_{x\to0}\frac{\sin^2x}{x^2}\)
  \item \(\displaystyle\lim_{x\to0}\frac{\sin^2(3x)}{x^2}\)
\end{bgexercises}

\subsection*{D. Mixed trigonometric limits}
\begin{bgexercises}[label=\textbf{3.\arabic*},resume]
  \item \(\displaystyle\lim_{x\to0}\frac{\sin(2x)\sin(3x)}{x^2}\)
  \item \(\displaystyle\lim_{x\to0}\frac{\tan(5x)}{\sin(2x)}\)
  \item \(\displaystyle\lim_{x\to0}\frac{1-\cos(2x)}{\sin^2x}\)
  \item \(\displaystyle\lim_{x\to0}\frac{\sin(4x)}{x\cos x}\)
  \item \(\displaystyle\lim_{x\to0}\frac{\sin x}{x(1+\cos x)}\)
  \item \(\displaystyle\lim_{x\to0}\frac{\tan x-\sin x}{x}\)
  \item \(\displaystyle\lim_{x\to0}\frac{1-\cos x}{x\sin x}\)
  \item \(\displaystyle\lim_{x\to0}\frac{\sin(3x)}{\tan(5x)}\)
\end{bgexercises}

\subsection*{E. Reasoning and exam practice}
\begin{bgexercises}[label=\textbf{3.\arabic*},resume]
  \item Explain why radians are required for \(\lim_{x\to0}\sin x/x=1\).
  \item Derive \(\lim_{x\to0}\tan x/x=1\) from the sine limit.
  \item Derive \(\lim_{x\to0}(1-\cos x)/x^2=1/2\).
  \item A student writes \(\sin(5x)=5x\). Explain what is wrong and how to use the limit correctly.
  \item Give a squeeze argument for \(\lim_{x\to0}x^4\sin(1/x^3)\).
  \item Create a trigonometric limit near zero whose value is \(7/4\).
\end{bgexercises}

Answers begin in \cref{app:chapter-answers}.
\webnext{calculus/limits/infinite-limits}{Infinite Limits Explained}
\end{webpage}


% --- END INLINE chapters/ch03_squeeze_trig ---


% --- BEGIN INLINE chapters/ch04_infinite ---
\begin{webpage}{calculus/limits/infinite-limits}{Infinite Limits Explained}{lesson}
\chapter{Infinite Behavior and Asymptotes}

\section{Infinite Limits}

\lessonobjective{
Interpret infinite-limit notation; determine one-sided signs near a vertical asymptote; distinguish unbounded behavior from an ordinary finite limit.
}

A finite limit asks whether outputs approach one real number. An infinite limit describes outputs whose magnitude grows without bound.

\begin{bgdefinition}[title={Infinite-Limit Notation}]
The statement
\[
  \lim_{x\to a}f(x)=+\infty
\]
means that \(f(x)\) becomes arbitrarily large and positive as \(x\) approaches \(a\).

The statement
\[
  \lim_{x\to a}f(x)=-\infty
\]
means that \(f(x)\) becomes arbitrarily negative as \(x\) approaches \(a\).
\end{bgdefinition}

Infinity is not a real number. The notation records a direction of unbounded growth.

\begin{bgconcept}
Imagine a thermometer with no top. Saying the reading approaches \(+\infty\) does not mean it arrives at a final number called infinity. It means that no matter how large a height you name, the reading eventually exceeds it near the target input.
\end{bgconcept}

\begin{bgguided}[title={\(1/x\) near zero}]
Consider \(f(x)=1/x\).

From the right of zero, use small positive numbers:
\[
\begin{array}{c|c}
x&1/x\\
\hline
0.1&10\\
0.01&100\\
0.001&1000
\end{array}
\]
Thus,
\[
  \lim_{x\to0^+}\frac1x=+\infty.
\]

From the left, use small negative numbers:
\[
\begin{array}{c|c}
x&1/x\\
\hline
-0.1&-10\\
-0.01&-100\\
-0.001&-1000
\end{array}
\]
Thus,
\[
  \lim_{x\to0^-}\frac1x=-\infty.
\]
The one-sided behaviors differ, so the ordinary two-sided limit does not exist.
\end{bgguided}

\begin{figure}[htbp]
\centering
\begin{tikzpicture}
\begin{groupplot}[
  group style={group size=2 by 1,horizontal sep=1.2cm},
  width=0.45\textwidth,height=0.42\textwidth,
  xmin=-4,xmax=4,ymin=-8,ymax=8,
  axis lines=middle,grid=major,major grid style={draw=BGLine},
  xlabel={$x$},ylabel={$y$}
]
\nextgroupplot[title={$y=1/(x-1)$}]
\addplot[BGBlue,very thick,domain=-4:0.94,samples=140]{1/(x-1)};
\addplot[BGBlue,very thick,domain=1.06:4,samples=140]{1/(x-1)};
\draw[dashed,BGRed] (axis cs:1,-8)--(axis cs:1,8);
\nextgroupplot[title={$y=1/(x-1)^2$}]
\addplot[BGGreen,very thick,domain=-4:0.94,samples=140]{1/(x-1)^2};
\addplot[BGGreen,very thick,domain=1.06:4,samples=140]{1/(x-1)^2};
\draw[dashed,BGRed] (axis cs:1,-8)--(axis cs:1,8);
\end{groupplot}
\end{tikzpicture}
\caption{An odd power changes sign across the vertical asymptote; an even power remains positive on both sides.}
\label{fig:odd-even-asymptote}
\end{figure}
\webnext{calculus/limits/vertical-asymptotes}{Vertical Asymptotes and One-Sided Sign Analysis}
\end{webpage}

\begin{webpage}{calculus/limits/vertical-asymptotes}{Vertical Asymptotes and One-Sided Sign Analysis}{lesson}
\section{Vertical Asymptotes and Sign Analysis}

\lessonobjective{
Find vertical asymptote candidates, cancel removable factors before classifying behavior, and use a sign chart to determine each one-sided infinite limit.
}

The line \(x=a\) is a vertical asymptote if the function becomes unbounded on at least one side of \(a\).

For a rational function
\[
  f(x)=\frac{p(x)}{q(x)},
\]
zeros of \(q(x)\) are candidates. But first simplify common factors.

\begin{bgmistake}
A denominator zero does not automatically create a vertical asymptote. If the zero factor cancels with the numerator, the graph may have a removable hole instead.
\end{bgmistake}

\begin{bgguided}[title={Identify the sign before saying infinity}]
Evaluate
\[
  \lim_{x\to2^-}\frac1{x-2}.
\]

\begin{bgsolution}
As \(x\to2^-\), \(x-2\) is a tiny negative number. Dividing \(1\) by a tiny negative number gives a very large negative number. Therefore,
\[
  \boxed{\lim_{x\to2^-}\frac1{x-2}=-\infty}.
\]
\end{bgsolution}
\end{bgguided}


% --- BEGIN INLINE checks/infinite-sign-01 ---
\begin{bgcheck}{infinite-sign-01}{choice}{+infinity}[Check your understanding]
Evaluate \(\lim_{x\to1^+}\frac1{x-1}\).
\begin{bghint}
From the right, the denominator is positive and tiny.
\end{bghint}
\begin{bgreveal}
The limit is \(+\infty\).
\end{bgreveal}
\end{bgcheck}

% --- END INLINE checks/infinite-sign-01 ---


\begin{bgmethod}[title={One-Sided Sign Chart}]
To analyze \(p(x)/q(x)\) near \(x=a\):
\begin{enumerate}
  \item Factor numerator and denominator.
  \item Cancel any common factor, while remembering the corresponding hole.
  \item Determine the sign of every remaining factor just left of \(a\).
  \item Determine the sign just right of \(a\).
  \item Combine signs and decide whether the magnitude grows without bound.
\end{enumerate}
\end{bgmethod}

\begin{bgexample}[title={A sign chart with several factors}]
Find both one-sided limits of
\[
  f(x)=\frac{x+2}{(x-3)(x+1)}
\]
at \(x=3\).

\begin{bgsolution}
Near \(x=3\), the factors \(x+2\) and \(x+1\) are positive. Only \(x-3\) changes sign.

From the left of \(3\):
\[
  x+2>0,\qquad x-3<0,\qquad x+1>0.
\]
The quotient is negative. Its denominator approaches zero in magnitude, so
\[
  \boxed{\lim_{x\to3^-}f(x)=-\infty}.
\]

From the right of \(3\):
\[
  x+2>0,\qquad x-3>0,\qquad x+1>0.
\]
The quotient is positive, so
\[
  \boxed{\lim_{x\to3^+}f(x)=+\infty}.
\]
Thus \(x=3\) is a vertical asymptote.
\end{bgsolution}
\end{bgexample}

\begin{bgexample}[title={Even multiplicity}]
Analyze
\[
  g(x)=\frac{x-4}{(x+2)^2}
\]
near \(x=-2\).

\begin{bgsolution}
Near \(-2\), the numerator \(x-4\) is negative. The denominator \((x+2)^2\) is positive on both sides and approaches zero.

Therefore, on both sides the quotient is negative with unbounded magnitude:
\[
  \boxed{\lim_{x\to-2^-}g(x)=-\infty},
\]
\[
  \boxed{\lim_{x\to-2^+}g(x)=-\infty}.
\]
The two one-sided infinite behaviors agree, so it is also common to write
\[
  \lim_{x\to-2}g(x)=-\infty.
\]
This does not mean the limit is a real number; it means both sides decrease without bound.
\end{bgsolution}
\end{bgexample}

\begin{bgexample}[title={Exam-level: hole versus asymptote}]
Find and classify every discontinuity of
\[
  h(x)=\frac{(x-1)(x+4)}{(x-1)(x-3)}.
\]

\begin{bgsolution}
The original denominator is zero at \(x=1\) and \(x=3\). Cancel the common factor for \(x\ne1\):
\[
  h(x)=\frac{x+4}{x-3}.
\]
At \(x=1\), the simplified formula is finite:
\[
  \frac{1+4}{1-3}=-\frac52.
\]
Therefore \(x=1\) is a removable hole, not a vertical asymptote.

At \(x=3\), the remaining denominator is zero while the numerator is \(7\ne0\). Thus \(x=3\) is a vertical asymptote.

For signs near \(3\), the numerator is positive. The denominator \(x-3\) is negative from the left and positive from the right:
\[
  \lim_{x\to3^-}h(x)=-\infty,
  \qquad
  \lim_{x\to3^+}h(x)=+\infty.
\]
\end{bgsolution}
\end{bgexample}
\webnext{calculus/limits/limits-at-infinity}{Limits at Infinity and Horizontal Asymptotes}
\end{webpage}

\begin{webpage}{calculus/limits/limits-at-infinity}{Limits at Infinity and Horizontal Asymptotes}{lesson}
\section{Limits at Infinity}

\lessonobjective{
Interpret \(x\to\pm\infty\); find horizontal asymptotes by comparing dominant terms; evaluate rational-function end behavior.
}

The notation
\[
  \lim_{x\to\infty}f(x)=L
\]
asks what happens to outputs as inputs become arbitrarily large and positive. The notation
\[
  \lim_{x\to-\infty}f(x)=L
\]
asks what happens far to the left.

If either end limit equals a finite number \(L\), then \(y=L\) is a horizontal asymptote on that end.

\begin{bgconcept}
A horizontal asymptote describes the far-away behavior of a graph. The graph may cross it nearby. An asymptote is not an electric fence. It is a long-term trend.
\end{bgconcept}

\subsection{Reciprocal powers}

For every positive integer \(n\),
\[
  \lim_{x\to\infty}\frac1{x^n}=0,
  \qquad
  \lim_{x\to-\infty}\frac1{x^n}=0.
\]
These facts drive the degree rules for rational functions.

\begin{bgguided}[title={One over a growing number}]
Evaluate
\[
  \lim_{x\to\infty}\frac1x.
\]

\begin{bgsolution}
As \(x\) becomes \(10,100,1000,\ldots\), the values \(1/x\) become
\[
  0.1,0.01,0.001,\ldots
\]
which approach zero. Therefore,
\[
  \boxed{0}.
\]
\end{bgsolution}
\end{bgguided}

\subsection{Equal degrees}

\begin{bgexample}[title={Divide by the largest denominator power}]
Evaluate
\[
  \lim_{x\to\infty}\frac{3x^2-2x+5}{x^2+4}.
\]

\begin{bgsolution}
Divide every term by \(x^2\), the largest power in the denominator:
\[
\frac{3x^2-2x+5}{x^2+4}
=\frac{3-\frac2x+\frac5{x^2}}{1+\frac4{x^2}}.
\]
As \(x\to\infty\), all reciprocal terms approach zero:
\[
  \frac{3-0+0}{1+0}=\boxed{3}.
\]
Thus \(y=3\) is a horizontal asymptote.
\end{bgsolution}


% --- BEGIN INLINE checks/end-degree-01 ---
\begin{bgcheck}{end-degree-01}{rational}{3/2}[Check your understanding]
Evaluate \(\lim_{x\to\infty}\frac{3x^2+1}{2x^2-5}\).
\begin{bghint}
Equal degrees: use the leading coefficients.
\end{bghint}
\begin{bgreveal}
The limit is \(3/2\).
\end{bgreveal}
\end{bgcheck}

% --- END INLINE checks/end-degree-01 ---

\end{bgexample}

\subsection{Numerator degree smaller than denominator degree}

\begin{bgexample}[title={The denominator wins}]
Evaluate
\[
  \lim_{x\to-\infty}\frac{5x-1}{x^2+3}.
\]

\begin{bgsolution}
Divide by \(x^2\):
\[
  \frac{\frac5x-\frac1{x^2}}{1+\frac3{x^2}}.
\]
Every reciprocal term approaches zero, so
\[
  \boxed{0}.
\]
The quadratic denominator grows faster in magnitude than the linear numerator.
\end{bgsolution}
\end{bgexample}

\subsection{Numerator degree larger than denominator degree}

If the numerator degree is larger, there is no finite horizontal asymptote. The quotient may grow like a polynomial.

\begin{bgexample}[title={Unbounded end behavior}]
Evaluate
\[
  \lim_{x\to\infty}\frac{2x^3-x}{x^2+1}.
\]

\begin{bgsolution}
The leading behavior is approximately
\[
  \frac{2x^3}{x^2}=2x,
\]
which grows to \(+\infty\). More formally, divide by \(x^2\):
\[
  \frac{2x-\frac1x}{1+\frac1{x^2}}.
\]
The numerator grows without bound while the denominator approaches \(1\), so
\[
  \boxed{+\infty}.
\]
\end{bgsolution}
\end{bgexample}

\begin{bgtheorem}[title={Rational-Function Degree Rules}]
For \(f(x)=p(x)/q(x)\):
\begin{itemize}
  \item If \(\deg p<\deg q\), then \(f(x)\to0\) as \(x\to\pm\infty\).
  \item If \(\deg p=\deg q\), the limit is the ratio of leading coefficients.
  \item If \(\deg p>\deg q\), there is no finite horizontal asymptote; use division or dominant terms to determine the end behavior.
\end{itemize}
\end{bgtheorem}

\begin{figure}[htbp]
\centering
\begin{tikzpicture}
\begin{axis}[
  width=0.86\textwidth,height=0.46\textwidth,
  xmin=-12,xmax=12,ymin=0,ymax=4.3,
  axis lines=middle,xlabel={$x$},ylabel={$f(x)$},
  grid=major,major grid style={draw=BGLine}
]
\addplot[BGBlue,very thick,domain=-12:12,samples=120]{(3*x^2-2*x+5)/(x^2+4)};
\draw[dashed,BGRed,thick] (axis cs:-12,3)--(axis cs:12,3);
\node[anchor=south west,font=\small,BGRed] at (axis cs:5,3) {$y=3$};
\end{axis}
\end{tikzpicture}
\caption{The rational function approaches its horizontal asymptote \(y=3\) at both ends.}
\label{fig:horizontal-asymptote}
\end{figure}
\webnext{calculus/limits/slant-asymptotes}{Polynomial and Slant Asymptotes}
\end{webpage}

\begin{webpage}{calculus/limits/slant-asymptotes}{Polynomial and Slant Asymptotes}{lesson}
\section{Polynomial and Slant Asymptotes}

\lessonobjective{
Use polynomial division when the numerator degree exceeds the denominator degree; identify a slant asymptote when the degree difference is one.
}

When the numerator degree is exactly one larger than the denominator degree, long division produces
\[
  f(x)=\text{linear quotient}+\frac{\text{remainder}}{q(x)}.
\]
The remainder term approaches zero, leaving a slant or oblique asymptote.

\begin{bgexample}[title={Slant asymptote by division}]
Find the end behavior of
\[
  f(x)=\frac{x^2+1}{x-1}.
\]

\begin{bgsolution}
Divide \(x^2+1\) by \(x-1\):
\[
  \frac{x^2+1}{x-1}=x+1+\frac2{x-1}.
\]
As \(x\to\pm\infty\),
\[
  \frac2{x-1}\to0.
\]
Therefore the graph approaches
\[
  \boxed{y=x+1}.
\]
This is a slant asymptote. The function does not approach a finite number, but its difference from \(x+1\) approaches zero:
\[
  \lim_{x\to\pm\infty}[f(x)-(x+1)]=0.
\]
\end{bgsolution}
\end{bgexample}
\webnext{calculus/limits/radicals-at-infinity}{Radical Limits at Infinity}
\end{webpage}

\begin{webpage}{calculus/limits/radicals-at-infinity}{Radical Limits at Infinity}{lesson}
\section{Radicals at Infinity and the Absolute-Value Trap}

\lessonobjective{
Factor the highest power from a radical correctly; use \(\sqrt{x^2}=|x|\); evaluate radical differences by rationalization.
}

The identity
\[
  \boxed{\sqrt{x^2}=|x|}
\]
is essential. At positive infinity, \(|x|=x\). At negative infinity, \(|x|=-x\).

\begin{bgguided}[title={Why the sign changes}]
Evaluate
\[
  \lim_{x\to-\infty}\frac{\sqrt{x^2}}{x}.
\]

\begin{bgsolution}
Since \(\sqrt{x^2}=|x|\),
\[
  \frac{\sqrt{x^2}}{x}=\frac{|x|}{x}.
\]
For negative \(x\), \(|x|=-x\). Therefore,
\[
  \frac{|x|}{x}=\frac{-x}{x}=-1.
\]
Hence
\[
  \boxed{-1}.
\]
Writing \(\sqrt{x^2}=x\) would incorrectly give \(+1\).
\end{bgsolution}
\end{bgguided}

\begin{bgexample}[title={The same radical on two ends}]
Evaluate
\[
  \lim_{x\to\infty}\frac{\sqrt{x^2+1}}{x}
  \qquad\text{and}\qquad
  \lim_{x\to-\infty}\frac{\sqrt{x^2+1}}{x}.
\]

\begin{bgsolution}
Factor \(x^2\) inside the radical:
\[
  \sqrt{x^2+1}=\sqrt{x^2\left(1+\frac1{x^2}\right)}
  =|x|\sqrt{1+\frac1{x^2}}.
\]
At positive infinity, \(|x|/x=1\), so
\[
  \boxed{1}.
\]
At negative infinity, \(|x|/x=-1\), so
\[
  \boxed{-1}.
\]
\end{bgsolution}
\end{bgexample}
\webnext{calculus/limits/infinity-minus-infinity}{Infinity Minus Infinity Limits}
\end{webpage}

\begin{webpage}{calculus/limits/infinity-minus-infinity}{Infinity Minus Infinity Limits}{lesson}
\subsection{The form \(\infty-\infty\)}

A difference of two large expressions is indeterminate. Rationalization can reveal the finite remainder.

\begin{bgexample}[title={Rationalize at infinity}]
Evaluate
\[
  \lim_{x\to\infty}\left(\sqrt{x^2+3x}-x\right).
\]

\begin{bgsolution}
Multiply by the conjugate:
\begin{align*}
\sqrt{x^2+3x}-x
&=\frac{(x^2+3x)-x^2}{\sqrt{x^2+3x}+x}\\
&=\frac{3x}{\sqrt{x^2+3x}+x}.
\end{align*}
Since \(x>0\) for large positive \(x\), divide numerator and denominator by \(x\):
\[
  \frac3{\sqrt{1+\frac3x}+1}.
\]
Now take the limit:
\[
  \frac3{1+1}=\boxed{\frac32}.
\]
\end{bgsolution}
\end{bgexample}

\begin{bgexample}[title={Exam-level: negative infinity changes the conjugate calculation}]
Evaluate
\[
  \lim_{x\to-\infty}\left(\sqrt{x^2+4x}+x\right).
\]

\begin{bgsolution}
Rationalize:
\begin{align*}
\sqrt{x^2+4x}+x
&=\frac{(x^2+4x)-x^2}{\sqrt{x^2+4x}-x}\\
&=\frac{4x}{\sqrt{x^2+4x}-x}.
\end{align*}
For \(x<0\), factor \(|x|=-x\) from the radical:
\[
  \sqrt{x^2+4x}=(-x)\sqrt{1+\frac4x}.
\]
Then the denominator becomes
\[
  -x\sqrt{1+\frac4x}-x
  =-x\left(\sqrt{1+\frac4x}+1\right).
\]
Therefore,
\[
  \frac{4x}{-x\left(\sqrt{1+\frac4x}+1\right)}
  =\frac{-4}{\sqrt{1+\frac4x}+1}.
\]
Taking the limit gives
\[
  \boxed{-2}.
\]
\end{bgsolution}
\end{bgexample}
\webnext{calculus/limits/complete-asymptote-analysis}{Complete Rational Function Asymptote Analysis}
\end{webpage}

\begin{webpage}{calculus/limits/complete-asymptote-analysis}{Complete Rational Function Asymptote Analysis}{lesson}
\section{Complete Asymptote Analysis}

\begin{bgmethod}[title={Rational-Function Analysis Checklist}]
For a rational function:
\begin{enumerate}
  \item Factor numerator and denominator.
  \item Cancel common factors and record holes.
  \item Set the remaining denominator equal to zero to find vertical asymptotes.
  \item Use one-sided sign analysis at each vertical asymptote.
  \item Compare degrees or divide polynomials to find horizontal or polynomial asymptotes.
  \item Check both \(x\to\infty\) and \(x\to-\infty\); radical functions may behave differently on the two ends.
\end{enumerate}
\end{bgmethod}
\webnext{calculus/continuity/continuity-at-a-point}{Continuity at a Point}
\end{webpage}

\begin{webpage}{calculus/limits/infinite-behavior-review}{Infinite Limits and Asymptotes Review}{review}
\section*{Section 4 Summary}
\addcontentsline{toc}{section}{Section 4 Summary}

\begin{bgsummary}
\begin{itemize}
  \item Infinite limits describe unbounded behavior; infinity is not a real number.
  \item Remaining denominator zeros after cancellation give vertical asymptote candidates.
  \item Sign charts determine \(+\infty\) versus \(-\infty\) on each side.
  \item Rational end behavior is controlled by degrees and leading coefficients.
  \item If the degree difference is one, polynomial division may reveal a slant asymptote.
  \item At negative infinity, never replace \(\sqrt{x^2}\) by \(x\); use \(|x|=-x\).
\end{itemize}
\end{bgsummary}
\webnext{calculus/continuity/continuity-at-a-point}{Continuity at a Point}
\end{webpage}

\begin{webpage}{calculus/limits/infinite-behavior-concept-quiz}{Infinite Limits Concept Quiz}{quiz}
\section*{Section 4 Concept Quiz}
\addcontentsline{toc}{section}{Section 4 Concept Quiz}

Use this as a short retrieval check after finishing the section. On the website, each item should accept an answer before revealing the hint and worked explanation.

\begin{bgcheck}{infinite-q1}{choice}{-infinity}[Concept check]
Evaluate \(\lim_{x\to2^-}\frac1{x-2}\).
\begin{bghint}
From the left, \(x-2\) is a small negative number.
\end{bghint}
\begin{bgreveal}
A positive numerator divided by a small negative number tends to \(-\infty\).
\end{bgreveal}
\end{bgcheck}

\begin{bgcheck}{vertical-q1}{integer}{3}[Concept check]
For \(f(x)=1/(x-3)^2\), enter the \(x\)-value of the vertical asymptote.
\begin{bghint}
A vertical asymptote occurs where the denominator is zero and does not cancel.
\end{bghint}
\begin{bgreveal}
\(x=3\).
\end{bgreveal}
\end{bgcheck}

\begin{bgcheck}{end-q1}{rational}{2/5}[Concept check]
Evaluate \(\lim_{x\to\infty}\frac{2x^2+1}{5x^2-3}\).
\begin{bghint}
The degrees match, so compare leading coefficients.
\end{bghint}
\begin{bgreveal}
The limit is \(2/5\).
\end{bgreveal}
\end{bgcheck}

\begin{bgcheck}{radical-infinity-q1}{integer}{-1}[Concept check]
Evaluate \(\lim_{x\to-\infty}\frac{\sqrt{x^2+1}}x\).
\begin{bghint}
Remember \(\sqrt{x^2}=|x|=-x\) when \(x<0\).
\end{bghint}
\begin{bgreveal}
The expression approaches \(-1\).
\end{bgreveal}
\end{bgcheck}

\begin{bgsummary}[title={Quiz use on the website}]
This page should report completion by concept, not merely a total score. A wrong answer should link back to the exact lesson that teaches the missing method. The same questions may appear in randomized numerical variants, but the canonical explanation should remain stable.
\end{bgsummary}
\webnext{calculus/continuity/continuity-at-a-point}{Continuity at a Point}
\end{webpage}

\begin{webpage}{calculus/limits/infinite-behavior-practice}{Infinite Limits and Asymptote Practice}{practice}
\section*{Section 4 Exercises}
\addcontentsline{toc}{section}{Section 4 Exercises}

\subsection*{A. One-sided infinite limits}
\begin{bgexercises}[label=\textbf{4.\arabic*}]
  \item \(\displaystyle\lim_{x\to4^-}\frac1{x-4}\)
  \item \(\displaystyle\lim_{x\to4^+}\frac1{x-4}\)
  \item \(\displaystyle\lim_{x\to0^-}\frac1{x^2}\)
  \item \(\displaystyle\lim_{x\to0^+}\frac1{x^3}\)
  \item \(\displaystyle\lim_{x\to-2^+}\frac{x-1}{x+2}\)
  \item \(\displaystyle\lim_{x\to1^-}\frac{x+3}{(x-1)^2}\)
  \item \(\displaystyle\lim_{x\to3^-}\frac{x+2}{(x-3)(x+1)}\)
  \item \(\displaystyle\lim_{x\to3^+}\frac{x+2}{(x-3)(x+1)}\)
  \item \(\displaystyle\lim_{x\to-1^-}\frac{2-x}{(x+1)^3}\)
  \item \(\displaystyle\lim_{x\to-1^+}\frac{2-x}{(x+1)^3}\)
\end{bgexercises}

\subsection*{B. Holes and vertical asymptotes}
\begin{bgexercises}[label=\textbf{4.\arabic*},resume]
  \item Find and classify discontinuities of \(\dfrac{x-2}{(x-2)(x+1)}\).
  \item Find and classify discontinuities of \(\dfrac{x^2-9}{(x-3)(x+2)}\).
  \item Find all holes and vertical asymptotes of \(\dfrac{(x+1)(x-4)}{(x+1)(x-2)^2}\).
  \item Determine one-sided behavior at every vertical asymptote of \(\dfrac{x+3}{(x-1)(x+2)}\).
  \item Determine one-sided behavior at every vertical asymptote of \(\dfrac{x-5}{(x+1)^2}\).
  \item Explain why a cancelled denominator factor creates a hole rather than a vertical asymptote.
\end{bgexercises}

\subsection*{C. Rational limits at infinity}
\begin{bgexercises}[label=\textbf{4.\arabic*},resume]
  \item \(\displaystyle\lim_{x\to\infty}\frac{4x+1}{x^2+7}\)
  \item \(\displaystyle\lim_{x\to-\infty}\frac{2x^2-3}{5x^2+x}\)
  \item \(\displaystyle\lim_{x\to\infty}\frac{7x^3+x}{2x^3-5}\)
  \item \(\displaystyle\lim_{x\to-\infty}\frac{3x^4-x}{x^4+9}\)
  \item \(\displaystyle\lim_{x\to\infty}\frac{x^2+1}{x^3+2}\)
  \item \(\displaystyle\lim_{x\to-\infty}\frac{5x^3+1}{x^2-4}\)
  \item \(\displaystyle\lim_{x\to\infty}\frac{-2x^5+x}{x^5+3}\)
  \item Find all horizontal asymptotes of \(\dfrac{3x^2+1}{x^2-5x}\).
  \item Find the slant asymptote of \(\dfrac{x^2+3}{x-2}\).
  \item Use division to describe the end behavior of \(\dfrac{2x^3+x}{x^2-1}\).
\end{bgexercises}

\subsection*{D. Radicals at infinity}
\begin{bgexercises}[label=\textbf{4.\arabic*},resume]
  \item \(\displaystyle\lim_{x\to\infty}\frac{\sqrt{x^2+4}}{x}\)
  \item \(\displaystyle\lim_{x\to-\infty}\frac{\sqrt{x^2+4}}{x}\)
  \item \(\displaystyle\lim_{x\to\infty}\frac{\sqrt{9x^2+1}}{x}\)
  \item \(\displaystyle\lim_{x\to-\infty}\frac{\sqrt{9x^2+1}}{x}\)
  \item \(\displaystyle\lim_{x\to\infty}\left(\sqrt{x^2+8x}-x\right)\)
  \item \(\displaystyle\lim_{x\to\infty}\left(\sqrt{4x^2+x}-2x\right)\)
  \item \(\displaystyle\lim_{x\to-\infty}\left(\sqrt{x^2+6x}+x\right)\)
  \item Explain the exact point where \(|x|\) must be used in a radical-at-infinity problem.
\end{bgexercises}

\subsection*{E. Mixed exam practice}
\begin{bgexercises}[label=\textbf{4.\arabic*},resume]
  \item Analyze all discontinuities and asymptotes of \(\dfrac{x^2-1}{x^2-3x+2}\).
  \item Analyze all discontinuities and asymptotes of \(\dfrac{x^2-4}{x^2-x-2}\).
  \item Find the exact one-sided limits at every vertical asymptote of \(\dfrac{x}{x^2-9}\).
  \item Find the end behavior of \(\dfrac{2x^3-x+1}{x^2+4}\), including any polynomial asymptote.
  \item Construct a rational function with a hole at \(x=1\), a vertical asymptote at \(x=-2\), and horizontal asymptote \(y=3\).
  \item Explain why a graph may cross a horizontal asymptote but cannot take a finite value on a vertical asymptote where its formula is undefined.
\end{bgexercises}

Answers begin in \cref{app:chapter-answers}.
\webnext{calculus/continuity/continuity-at-a-point}{Continuity at a Point}
\end{webpage}


% --- END INLINE chapters/ch04_infinite ---


% --- BEGIN INLINE chapters/ch05_continuity ---
\begin{webpage}{calculus/continuity/continuity-at-a-point}{Continuity at a Point}{lesson}
\chapter{Continuity and the Intermediate Value Theorem}

\section{Continuity at a Point}

\lessonobjective{
Use the three-part continuity test at a point; distinguish continuity from merely having a limit; explain the graphical meaning of continuity.
}

A limit describes nearby behavior. Continuity connects that nearby behavior to the function's actual value.

\begin{bgdefinition}[title={Continuity at \(x=a\)}]
A function \(f\) is continuous at \(x=a\) if all three conditions hold:
\begin{enumerate}
  \item \(f(a)\) is defined;
  \item \(\displaystyle\lim_{x\to a}f(x)\) exists;
  \item \(\displaystyle\lim_{x\to a}f(x)=f(a)\).
\end{enumerate}
Equivalently,
\[
  \boxed{\lim_{x\to a}f(x)=f(a)}.
\]
\end{bgdefinition}

\begin{bgconcept}
Continuity means the graph, the nearby trend, and the actual dot all agree at the point. You approach one height, and the function is actually located at that height.
\end{bgconcept}

\begin{bgguided}[title={A continuous line}]
Is \(f(x)=2x+1\) continuous at \(x=3\)?

\begin{bgsolution}
Check the three conditions.

\textbf{1. The function value exists:}
\[
  f(3)=2(3)+1=7.
\]

\textbf{2. The limit exists:}
\[
  \lim_{x\to3}(2x+1)=7.
\]

\textbf{3. They agree:}
\[
  \lim_{x\to3}f(x)=f(3)=7.
\]
Therefore, \(f\) is continuous at \(3\).
\end{bgsolution}
\end{bgguided}

\begin{bgexample}[title={A limit exists but continuity fails}]
Let
\[
  g(x)=
  \begin{cases}
    x^2+1,&x\ne2,\\
    9,&x=2.
  \end{cases}
\]
Is \(g\) continuous at \(2\)?

\begin{bgsolution}
\textbf{Condition 1:} \(g(2)=9\), so the function value exists.

\textbf{Condition 2:} Nearby values use \(x^2+1\), so
\[
  \lim_{x\to2}g(x)=5.
\]
The limit exists.

\textbf{Condition 3:} Compare:
\[
  \lim_{x\to2}g(x)=5\ne9=g(2).
\]
Therefore, \(g\) is not continuous at \(2\).

This discontinuity is removable. Redefining \(g(2)=5\) would repair it.
\end{bgsolution}
\end{bgexample}

\begin{bgexample}[title={Continuity fails because no two-sided limit exists}]
Let
\[
  h(x)=
  \begin{cases}
    x+1,&x<2,\\
    6-x,&x\ge2.
  \end{cases}
\]
Is \(h\) continuous at \(2\)?

\begin{bgsolution}
The function value exists:
\[
  h(2)=6-2=4.
\]
But
\[
  \lim_{x\to2^-}h(x)=3,
\]
while
\[
  \lim_{x\to2^+}h(x)=4.
\]
The one-sided limits disagree, so the two-sided limit does not exist. Condition 2 fails. Therefore, \(h\) is not continuous at \(2\).
\end{bgsolution}
\end{bgexample}

\begin{bgexamnote}
Do not use only the slogan ``draw without lifting your pencil.'' It is a helpful picture, not a complete test. At endpoints, one-sided continuity is allowed; on disconnected domains, a graph can be continuous at every point of its domain even though you cannot draw all pieces in one stroke.
\end{bgexamnote}
\webnext{calculus/continuity/intervals-endpoints}{Continuity on Intervals and at Endpoints}
\end{webpage}

\begin{webpage}{calculus/continuity/intervals-endpoints}{Continuity on Intervals and at Endpoints}{lesson}
\section{Continuity on Intervals and at Endpoints}

\lessonobjective{
Determine intervals of continuity; use one-sided limits at endpoints; state the natural domains on which familiar functions are continuous.
}

A function is continuous on an open interval if it is continuous at every point in the interval.

At the left endpoint \(a\) of a closed interval \([a,b]\), continuity requires
\[
  \lim_{x\to a^+}f(x)=f(a).
\]
At the right endpoint \(b\), continuity requires
\[
  \lim_{x\to b^-}f(x)=f(b).
\]

\begin{bgconcept}
At the left edge of a road, there is no road to the left. It would be silly to demand traffic from a direction that does not exist. Endpoint continuity checks only from inside the interval.
\end{bgconcept}

\begin{bgguided}[title={Square root at its endpoint}]
Is \(f(x)=\sqrt{x}\) continuous at \(x=0\) on its domain \([0,\infty)\)?

\begin{bgsolution}
The function value is
\[
  f(0)=0.
\]
Only the right-hand limit is relevant because negative inputs are outside the real domain:
\[
  \lim_{x\to0^+}\sqrt{x}=0.
\]
The right-hand limit equals the function value, so \(f\) is continuous at the endpoint \(0\).
\end{bgsolution}
\end{bgguided}
\webnext{calculus/continuity/continuous-functions}{Which Functions Are Continuous?}
\end{webpage}

\begin{webpage}{calculus/continuity/continuous-functions}{Which Functions Are Continuous?}{reference}
\subsection{Families of continuous functions}

The following functions are continuous everywhere in their natural domains:

\begin{itemize}
  \item polynomials;
  \item rational functions where the denominator is nonzero;
  \item root functions where the root is real;
  \item trigonometric functions wherever defined;
  \item exponential functions;
  \item logarithmic functions for positive arguments.
\end{itemize}

Sums, differences, products, constant multiples, and valid compositions of continuous functions are continuous. Quotients are continuous where their denominators are nonzero.

\begin{bgexample}[title={Find intervals of continuity}]
Find the intervals on which
\[
  f(x)=\frac{\sqrt{x+2}}{x^2-9}
\]
is continuous.

\begin{bgsolution}
The square root requires
\[
  x+2\ge0,
\]
so
\[
  x\ge-2.
\]
The denominator must be nonzero:
\[
  x^2-9\ne0
\]
so
\[
  x\ne-3,3.
\]
The value \(-3\) is already outside the square-root domain. The only excluded domain point within \([-2,\infty)\) is \(3\).

Therefore, the function is continuous on
\[
  \boxed{[-2,3)\cup(3,\infty)}.
\]
At \(-2\), continuity is interpreted from the right.
\end{bgsolution}
\end{bgexample}

\begin{bgexample}[title={Composition}]
Where is
\[
  g(x)=\ln(4-x^2)
\]
continuous?

\begin{bgsolution}
The logarithm requires a positive argument:
\[
  4-x^2>0.
\]
This means
\[
  x^2<4,
\]
so
\[
  -2<x<2.
\]
The polynomial \(4-x^2\) is continuous everywhere, and \(\ln u\) is continuous for \(u>0\). Therefore the composition is continuous on
\[
  \boxed{(-2,2)}.
\]
\end{bgsolution}
\end{bgexample}
\webnext{calculus/continuity/types-of-discontinuity}{Types of Discontinuity}
\end{webpage}

\begin{webpage}{calculus/continuity/types-of-discontinuity}{Types of Discontinuity}{lesson}
\section{Classifying Discontinuities}

\lessonobjective{
Distinguish removable, jump, infinite, and oscillatory discontinuities; connect each classification to the relevant limit behavior.
}

\begin{center}
\begin{tabularx}{\textwidth}{@{}p{0.24\textwidth}X X@{}}
\toprule
\textbf{Type} & \textbf{Limit behavior} & \textbf{Typical graph feature}\\
\midrule
Removable & Finite two-sided limit exists, but value is missing or wrong & Hole or displaced dot\\
Jump & Finite one-sided limits exist but disagree & Sudden step\\
Infinite & Function becomes unbounded & Vertical asymptote\\
Oscillatory & Outputs do not settle to one value & Endless oscillation\\
\bottomrule
\end{tabularx}
\end{center}

\begin{figure}[htbp]
\centering
\begin{tikzpicture}
\begin{groupplot}[
 group style={group size=2 by 2,horizontal sep=1cm,vertical sep=1cm},
 width=0.43\textwidth,height=0.32\textwidth,
 xmin=-2,xmax=2,ymin=-2.2,ymax=3.2,
 axis lines=middle,grid=major,major grid style={draw=BGLine},
 xtick=\empty,ytick=\empty
]
\nextgroupplot[title={Removable}]
\addplot[BGBlue,very thick,domain=-2:-0.05]{x+1};
\addplot[BGBlue,very thick,domain=0.05:2]{x+1};
\addplot[only marks,mark=o,mark size=4pt,very thick,BGRed] coordinates {(0,1)};
\nextgroupplot[title={Jump}]
\addplot[BGBlue,very thick,domain=-2:0]{0};
\addplot[BGGreen,very thick,domain=0:2]{2};
\addplot[only marks,mark=o,mark size=4pt,very thick,BGBlue] coordinates {(0,0)};
\addplot[only marks,mark=*,mark size=3pt,BGGreen] coordinates {(0,2)};
\nextgroupplot[title={Infinite}]
\addplot[BGBlue,very thick,domain=-2:-0.06,samples=120]{1/x};
\addplot[BGBlue,very thick,domain=0.06:2,samples=120]{1/x};
\nextgroupplot[title={Oscillatory}]
\addplot[BGBlue,thick,domain=-2:-0.08,samples=120]{sin(deg(1/x))};
\addplot[BGBlue,thick,domain=0.08:2,samples=120]{sin(deg(1/x))};
\end{groupplot}
\end{tikzpicture}
\caption{The four principal discontinuity types encountered in first-semester calculus.}
\label{fig:discontinuities}
\end{figure}

\begin{bgexample}[title={Classify a rational function's discontinuities}]
Classify every discontinuity of
\[
  f(x)=\frac{x^2-1}{x^2-x-2}.
\]

\begin{bgsolution}
Factor numerator and denominator:
\[
  x^2-1=(x-1)(x+1),
\]
\[
  x^2-x-2=(x-2)(x+1).
\]
For \(x\ne-1\),
\[
  f(x)=\frac{x-1}{x-2}.
\]
At \(x=-1\), the common factor cancelled. The simplified function is finite:
\[
  \frac{-1-1}{-1-2}=\frac{-2}{-3}=\frac23.
\]
Thus \(x=-1\) is a removable discontinuity with a hole at \((-1,2/3)\).

At \(x=2\), the remaining denominator is zero and the numerator is nonzero. Thus \(x=2\) is an infinite discontinuity and vertical asymptote.
\end{bgsolution}
\end{bgexample}
\webnext{calculus/continuity/repair-removable-discontinuity}{How to Repair a Removable Discontinuity}
\end{webpage}

\begin{webpage}{calculus/continuity/repair-removable-discontinuity}{How to Repair a Removable Discontinuity}{lesson}
\section{Repairing Removable Discontinuities}

\lessonobjective{
Choose a missing function value so that a function becomes continuous at a point.
}

If a finite limit exists at \(a\), define the missing or incorrect function value to equal that limit.

\begin{bgguided}[title={Fill the hole with the approached value}]
Let
\[
  f(x)=
  \begin{cases}
    \dfrac{x^2-4}{x-2},&x\ne2,\\
    c,&x=2.
  \end{cases}
\]
Find \(c\) so that \(f\) is continuous at \(2\).

\begin{bgsolution}
For \(x\ne2\),
\[
  \frac{x^2-4}{x-2}=x+2.
\]
Therefore,
\[
  \lim_{x\to2}f(x)=4.
\]
Continuity requires
\[
  f(2)=\lim_{x\to2}f(x).
\]
Since \(f(2)=c\), choose
\[
  \boxed{c=4}.
\]
\end{bgsolution}
\end{bgguided}


% --- BEGIN INLINE checks/continuity-flow-01 ---
\begin{bgcheck}{continuity-flow-01}{integer}{8}[Check your understanding]
Choose \(c\) so \(f(x)=\frac{x^2-16}{x-4}\) for \(x\ne4\), and \(f(4)=c\), is continuous.
\begin{bghint}
Find the limit at \(4\) after factoring.
\end{bghint}
\begin{bgreveal}
The limit is \(8\), so set \(c=8\).
\end{bgreveal}
\end{bgcheck}

% --- END INLINE checks/continuity-flow-01 ---


\begin{bgexample}[title={Repair a more complicated hole}]
Choose \(c\) so that
\[
  g(x)=
  \begin{cases}
    \dfrac{x^3-27}{x-3},&x\ne3,\\
    c,&x=3
  \end{cases}
\]
is continuous at \(3\).

\begin{bgsolution}
Factor the difference of cubes:
\[
  x^3-27=(x-3)(x^2+3x+9).
\]
Thus, for \(x\ne3\),
\[
  g(x)=x^2+3x+9.
\]
Take the limit:
\[
  3^2+3(3)+9=9+9+9=27.
\]
Therefore,
\[
  \boxed{c=27}.
\]
\end{bgsolution}
\end{bgexample}

\begin{bgmistake}
Only removable discontinuities can be repaired by changing one function value. A jump or vertical asymptote reflects the behavior of infinitely many nearby points, not one misplaced dot.
\end{bgmistake}
\webnext{calculus/continuity/piecewise-parameters}{Choosing Parameters for Piecewise Continuity}
\end{webpage}

\begin{webpage}{calculus/continuity/piecewise-parameters}{Choosing Parameters for Piecewise Continuity}{lesson}
\section{Choosing Parameters in Piecewise Functions}

\lessonobjective{
Set the left-hand expression, right-hand expression, and function value equal at a joining point; solve for unknown parameters.
}

At a piecewise join \(x=a\), continuity requires
\[
  \lim_{x\to a^-}f(x)
  =f(a)
  =\lim_{x\to a^+}f(x).
\]
For ordinary polynomial pieces, substitute \(x=a\) into both formulas and set the results equal.

\begin{bgguided}[title={One parameter}]
Find \(k\) so that
\[
  f(x)=
  \begin{cases}
    kx+1,&x<2,\\
    x+3,&x\ge2
  \end{cases}
\]
is continuous at \(2\).

\begin{bgsolution}
Left side at the join:
\[
  2k+1.
\]
Right side and function value:
\[
  2+3=5.
\]
Set them equal:
\[
  2k+1=5.
\]
Subtract \(1\):
\[
  2k=4.
\]
Divide by \(2\):
\[
  \boxed{k=2}.
\]
\end{bgsolution}
\end{bgguided}

\begin{bgexample}[title={The parameter appears in both pieces}]
Find \(a\) so that
\[
  g(x)=
  \begin{cases}
    ax+1,&x<2,\\
    x^2-a,&x\ge2
  \end{cases}
\]
is continuous at \(2\).

\begin{bgsolution}
The left-hand limit is
\[
  2a+1.
\]
The right-hand limit and function value are
\[
  2^2-a=4-a.
\]
Set them equal:
\[
  2a+1=4-a.
\]
Add \(a\) to both sides:
\[
  3a+1=4.
\]
Subtract \(1\):
\[
  3a=3.
\]
Therefore,
\[
  \boxed{a=1}.
\]
\end{bgsolution}
\end{bgexample}

\begin{bgexample}[title={Exam-level: no parameter works}]
Find \(k\) so that
\[
  h(x)=
  \begin{cases}
    kx+2,&x<1,\\
    x^2+k,&x\ge1
  \end{cases}
\]
is continuous at \(1\).

\begin{bgsolution}
The left-hand expression approaches
\[
  k+2.
\]
The right-hand expression and function value equal
\[
  1+k.
\]
Continuity would require
\[
  k+2=k+1.
\]
Subtracting \(k\) gives
\[
  2=1,
\]
which is impossible. Therefore,
\[
  \boxed{\text{no value of }k\text{ works}}.
\]
A parameter problem is not guaranteed to have a parameter solution. The algebra is allowed to reject the premise, rude though that may seem to a worksheet.
\end{bgsolution}
\end{bgexample}
\webnext{calculus/continuity/two-parameter-continuity}{Two-Parameter Continuity Problems}
\end{webpage}

\begin{webpage}{calculus/continuity/two-parameter-continuity}{Two-Parameter Continuity Problems}{lesson}
\subsection{Two conditions and two parameters}

\begin{bgexample}[title={Solve a two-parameter continuity system}]
Choose \(a\) and \(b\) so that
\[
  f(x)=
  \begin{cases}
    ax+b,&x<1,\\
    x^2+1,&1\le x<3,\\
    2x+a,&x\ge3
  \end{cases}
\]
is continuous at both \(1\) and \(3\).

\begin{bgsolution}
At \(x=1\), continuity requires
\[
  a+b=1^2+1=2.
\]
At \(x=3\), continuity requires
\[
  3^2+1=2(3)+a.
\]
Thus,
\[
  10=6+a,
\]
so
\[
  a=4.
\]
Use \(a+b=2\):
\[
  4+b=2,
\]
so
\[
  b=-2.
\]
Therefore,
\[
  \boxed{a=4,\quad b=-2}.
\]
\end{bgsolution}
\end{bgexample}
\webnext{calculus/continuity/intermediate-value-theorem}{Intermediate Value Theorem}
\end{webpage}

\begin{webpage}{calculus/continuity/intermediate-value-theorem}{Intermediate Value Theorem}{lesson}
\section{The Intermediate Value Theorem}

\lessonobjective{
Use continuity and endpoint values to prove that a function attains an intermediate output or has a root in an interval; distinguish existence from calculation.
}

\begin{bgconcept}
A continuous path from below a horizontal line to above that line must cross it somewhere. You may not know exactly where the crossing occurs, but skipping it would require a jump.
\end{bgconcept}

\begin{bgtheorem}[title={Intermediate Value Theorem}]
Suppose \(f\) is continuous on \([a,b]\). If \(N\) lies between \(f(a)\) and \(f(b)\), then there exists at least one \(c\in[a,b]\) such that
\[
  \boxed{f(c)=N}.
\]
In particular, if \(f(a)\) and \(f(b)\) have opposite signs, then there exists \(c\in(a,b)\) such that \(f(c)=0\).
\end{bgtheorem}

\begin{bgguided}[title={Crossing ground level}]
Suppose a continuous function has
\[
  f(1)=-2
  \qquad\text{and}\qquad
  f(4)=5.
\]
The graph starts below the \(x\)-axis and ends above it. Because it is continuous, it must cross the axis somewhere between \(1\) and \(4\). Thus there is some \(c\in(1,4)\) with \(f(c)=0\).
\end{bgguided}


% --- BEGIN INLINE checks/ivt-flow-01 ---
\begin{bgcheck}{ivt-flow-01}{choice}{yes}[Check your understanding]
Does the IVT guarantee a root of \(x^3+x-1\) on \([0,1]\)?
\begin{bghint}
Check continuity and endpoint signs.
\end{bghint}
\begin{bgreveal}
Yes. The values are \(-1\) and \(1\).
\end{bgreveal}
\end{bgcheck}

% --- END INLINE checks/ivt-flow-01 ---


\begin{bgexample}[title={Prove a polynomial has a root}]
Show that
\[
  x^3+x-1=0
\]
has at least one solution in \((0,1)\).

\begin{bgsolution}
Define
\[
  f(x)=x^3+x-1.
\]
Polynomials are continuous everywhere, so \(f\) is continuous on \([0,1]\).

Evaluate the endpoints:
\[
  f(0)=0+0-1=-1,
\]
\[
  f(1)=1+1-1=1.
\]
Because \(-1<0<1\), the Intermediate Value Theorem guarantees some \(c\in(0,1)\) such that
\[
  f(c)=0.
\]
Therefore, the equation has at least one solution in \((0,1)\).
\end{bgsolution}
\end{bgexample}

\begin{figure}[htbp]
\centering
\begin{tikzpicture}
\begin{axis}[
  width=0.8\textwidth,height=0.48\textwidth,
  xmin=0,xmax=1,ymin=-1.15,ymax=1.15,
  axis lines=middle,xlabel={$x$},ylabel={$f(x)$},
  grid=major,major grid style={draw=BGLine}
]
\addplot[BGBlue,very thick,domain=0:1,samples=140]{x^3+x-1};
\addplot[only marks,mark=*,mark size=2.7pt,BGGold] coordinates {(0,-1) (1,1)};
\addplot[only marks,mark=*,mark size=2.7pt,BGGreen] coordinates {(0.6823278,0)};
\node[anchor=south,font=\small] at (axis cs:0.6823278,0) {$c$};
\end{axis}
\end{tikzpicture}
\caption{Continuity and opposite endpoint signs guarantee at least one root between \(0\) and \(1\).}
\label{fig:ivt}
\end{figure}

\begin{bgexamnote}
The Intermediate Value Theorem does not give the exact root, does not prove the root is unique, and cannot be used unless continuity on the entire closed interval has been established.
\end{bgexamnote}
\webnext{calculus/continuity/bisection-method}{Bisection Method After the IVT}
\end{webpage}

\begin{webpage}{calculus/continuity/bisection-method}{Bisection Method After the IVT}{lesson}
\section{Bisection: Turning Existence Into an Approximation}

\lessonobjective{
Use repeated sign changes on smaller intervals to approximate a root guaranteed by the Intermediate Value Theorem.
}

For
\[
  f(x)=x^3+x-1,
\]
we know a root lies in \((0,1)\). Test the midpoint \(0.5\):
\[
  f(0.5)=0.125+0.5-1=-0.375.
\]
Since \(f(0.5)<0\) and \(f(1)>0\), a root lies in \((0.5,1)\).

Test the new midpoint \(0.75\):
\[
  f(0.75)=0.421875+0.75-1=0.171875>0.
\]
Now the root lies in \((0.5,0.75)\).

Continue:
\[
\begin{array}{c|c|c}
\text{interval}&\text{midpoint}&\text{sign at midpoint}\\
\hline
(0,1)&0.5&-\\
(0.5,1)&0.75&+\\
(0.5,0.75)&0.625&-\\
(0.625,0.75)&0.6875&+\\
(0.625,0.6875)&0.65625&-
\end{array}
\]
The root is being trapped in shorter intervals. This is the bisection method.

\begin{bgconcept}
The Intermediate Value Theorem says a root is in the box. Bisection keeps cutting the box in half and throwing away the half that cannot contain the sign change.
\end{bgconcept}
\webnext{calculus/limits/epsilon-delta-introduction}{Epsilon-Delta Definition: An Introduction}
\end{webpage}

\begin{webpage}{calculus/continuity/continuity-review}{Continuity and IVT Review}{review}
\section*{Section 5 Summary}
\addcontentsline{toc}{section}{Section 5 Summary}

\begin{bgsummary}
\begin{itemize}
  \item Continuity at \(a\) requires a defined value, an existing limit, and equality between them.
  \item At endpoints, use the one-sided limit from inside the interval.
  \item Removable discontinuities can be repaired with one value; jumps and asymptotes cannot.
  \item For piecewise continuity, set the left and right expressions equal at the join.
  \item The Intermediate Value Theorem guarantees an attained value, not its exact location or uniqueness.
  \item Bisection refines an IVT interval into a numerical approximation.
\end{itemize}
\end{bgsummary}
\webnext{calculus/limits/epsilon-delta-introduction}{Epsilon-Delta Definition: An Introduction}
\end{webpage}

\begin{webpage}{calculus/continuity/continuity-concept-quiz}{Continuity Concept Quiz}{quiz}
\section*{Section 5 Concept Quiz}
\addcontentsline{toc}{section}{Section 5 Concept Quiz}

Use this as a short retrieval check after finishing the section. On the website, each item should accept an answer before revealing the hint and worked explanation.

\begin{bgcheck}{continuity-q1}{choice}{yes}[Concept check]
Is \(f(x)=x^2+1\) continuous at \(x=3\)? Enter yes or no.
\begin{bghint}
Polynomials are continuous everywhere.
\end{bghint}
\begin{bgreveal}
Yes. The function is defined, the limit exists, and both equal \(10\).
\end{bgreveal}
\end{bgcheck}

\begin{bgcheck}{repair-q1}{integer}{6}[Concept check]
Choose \(c\) so \(f(x)=\frac{x^2-9}{x-3}\) for \(x\ne3\), and \(f(3)=c\), is continuous.
\begin{bghint}
Factor, cancel, and find the approached value at \(3\).
\end{bghint}
\begin{bgreveal}
The nearby formula is \(x+3\), so choose \(c=6\).
\end{bgreveal}
\end{bgcheck}

\begin{bgcheck}{piecewise-q1}{integer}{2}[Concept check]
Choose \(a\) so \(2x+a\) for \(x<1\) joins continuously to \(x^2+3\) for \(x\ge1\).
\begin{bghint}
Set the left value at \(1\) equal to the right value at \(1\).
\end{bghint}
\begin{bgreveal}
\(2+a=4\), so \(a=2\).
\end{bgreveal}
\end{bgcheck}

\begin{bgcheck}{ivt-q1}{choice}{yes}[Concept check]
Does the IVT guarantee a zero of \(x^3+x-1\) on \([0,1]\)? Enter yes or no.
\begin{bghint}
Check continuity and the signs at both endpoints.
\end{bghint}
\begin{bgreveal}
Yes. The polynomial is continuous, with values \(-1\) and \(1\).
\end{bgreveal}
\end{bgcheck}

\begin{bgsummary}[title={Quiz use on the website}]
This page should report completion by concept, not merely a total score. A wrong answer should link back to the exact lesson that teaches the missing method. The same questions may appear in randomized numerical variants, but the canonical explanation should remain stable.
\end{bgsummary}
\webnext{calculus/limits/epsilon-delta-introduction}{Epsilon-Delta Definition: An Introduction}
\end{webpage}

\begin{webpage}{calculus/continuity/continuity-practice}{Continuity and IVT Practice Problems}{practice}
\section*{Section 5 Exercises}
\addcontentsline{toc}{section}{Section 5 Exercises}

\subsection*{A. Three-part continuity test}
\begin{bgexercises}[label=\textbf{5.\arabic*}]
  \item Is \(f(x)=x^2+1\) continuous at \(x=2\)? Verify all three conditions.
  \item Is \(g(x)=1/(x-3)\) continuous at \(x=3\)? Identify the first failed condition.
  \item A function has \(f(1)=4\) and \(\lim_{x\to1}f(x)=4\). Is it continuous at \(1\)?
  \item A function has \(f(1)=4\) and \(\lim_{x\to1}f(x)=5\). Which condition fails?
  \item A function is undefined at \(2\), but its limit there is \(7\). Which conditions fail?
  \item A function has unequal one-sided limits at \(0\). Can it be continuous there?
  \item Explain why existence of \(f(a)\) alone says nothing about continuity.
  \item Explain why existence of a finite limit alone does not guarantee continuity.
\end{bgexercises}

\subsection*{B. Intervals of continuity}
\begin{bgexercises}[label=\textbf{5.\arabic*},resume]
  \item Find intervals of continuity of \(1/(x^2-4)\).
  \item Find intervals of continuity of \(\sqrt{x-1}\).
  \item Find intervals of continuity of \(\sqrt{5-x}/(x+2)\).
  \item Find intervals of continuity of \(\ln(x+3)\).
  \item Find intervals of continuity of \(\tan x\) on \([-2\pi,2\pi]\).
  \item Find intervals of continuity of \(\sqrt{x+2}/(x^2-9)\).
  \item Find the domain and intervals of continuity of \(\ln(9-x^2)\).
  \item Explain endpoint continuity for \(f(x)=\sqrt{4-x^2}\) on \([-2,2]\).
\end{bgexercises}

\subsection*{C. Classifying discontinuities}
\begin{bgexercises}[label=\textbf{5.\arabic*},resume]
  \item Classify the discontinuity of \((x^2-1)/(x-1)\) at \(1\).
  \item Classify the discontinuity of \(|x|/x\) at \(0\).
  \item Classify the discontinuity of \(1/(x-2)^2\) at \(2\).
  \item Classify the discontinuity of \(\sin(1/x)\) at \(0\).
  \item Find and classify every discontinuity of \((x^2-4)/(x^2-x-2)\).
  \item Find and classify every discontinuity of \((x^2-9)/(x^2-6x+9)\).
  \item Can changing one function value repair a jump? Explain.
  \item Can changing one function value repair a vertical asymptote? Explain.
\end{bgexercises}

\subsection*{D. Repairing holes}
\begin{bgexercises}[label=\textbf{5.\arabic*},resume]
  \item Find \(c\) so \(f(2)=c\) makes \((x^2-4)/(x-2)\) continuous.
  \item Find \(c\) so \(f(3)=c\) makes \((x^2-9)/(x-3)\) continuous.
  \item Find \(c\) so \(f(1)=c\) makes \((x^3-1)/(x-1)\) continuous.
  \item Find \(c\) so \(f(4)=c\) makes \((\sqrt{x}-2)/(x-4)\) continuous.
  \item Explain why no value at \(x=0\) makes \(|x|/x\) continuous there.
  \item Explain why no value at \(x=2\) makes \(1/(x-2)\) continuous there.
\end{bgexercises}

\subsection*{E. Piecewise parameters}
\begin{bgexercises}[label=\textbf{5.\arabic*},resume]
  \item Find \(k\) so \(\begin{cases}kx+1,&x<2,\\x+3,&x\ge2\end{cases}\) is continuous at \(2\).
  \item Find \(a\) so \(\begin{cases}2x+a,&x<1,\\x^2+3,&x\ge1\end{cases}\) is continuous at \(1\).
  \item Find \(b\) so \(\begin{cases}x^2+b,&x\le-1,\\3x+1,&x>-1\end{cases}\) is continuous at \(-1\).
  \item Find \(m\) so \(\begin{cases}mx-4,&x<3,\\2x+m,&x\ge3\end{cases}\) is continuous at \(3\).
  \item Determine whether any \(k\) makes \(\begin{cases}kx+5,&x<0,\\kx+2,&x\ge0\end{cases}\) continuous at \(0\).
  \item Find \(a,b\) so \(\begin{cases}ax+b,&x<1,\\x^2+1,&1\le x<3,\\2x+a,&x\ge3\end{cases}\) is continuous at both joins.
  \item Design a piecewise function with one parameter that becomes continuous when the parameter is \(7\).
  \item Explain why setting left and right formulas equal is necessary but not always sufficient when the actual function value is defined by a third rule.
\end{bgexercises}

\subsection*{F. Intermediate Value Theorem and bisection}
\begin{bgexercises}[label=\textbf{5.\arabic*},resume]
  \item Show that \(x^3-2x-2=0\) has a root in \((1,2)\).
  \item Show that \(x^5+x-2=0\) has a root in \((0,2)\).
  \item Show that \(\cos x=x\) has a solution in \((0,1)\).
  \item Explain why a sign change is sufficient but not necessary for a root to exist.
  \item Give a continuous function with a root in \((-1,1)\) whose endpoint values have the same sign.
  \item Explain why IVT cannot be applied to \(1/x\) on \([-1,1]\), despite opposite endpoint signs.
  \item Use two bisection steps to narrow a root of \(x^3+x-1\) from \((0,1)\).
  \item Use three bisection steps to narrow a root of \(x^3-2x-2\) from \((1,2)\).
  \item Does IVT prove uniqueness? Give an example supporting your answer.
  \item State every hypothesis and conclusion in a complete IVT solution.
\end{bgexercises}

Answers begin in \cref{app:chapter-answers}.
\webnext{calculus/limits/epsilon-delta-introduction}{Epsilon-Delta Definition: An Introduction}
\end{webpage}


% --- END INLINE chapters/ch05_continuity ---


% --- BEGIN INLINE chapters/ch06_formal ---
\begin{webpage}{calculus/limits/epsilon-delta-introduction}{Epsilon-Delta Definition: An Introduction}{lesson}
\chapter{Making Limits Precise}

\section{Why the Informal Definition Needs Numbers}

\lessonobjective{
Translate ``arbitrarily close'' into inequalities; understand the roles of \(\eps\) and \(\delta\); verify simple formal limit proofs.
}

The informal definition says that \(f(x)\) can be made as close as desired to \(L\) by taking \(x\) sufficiently close to \(a\). The phrase ``as close as desired'' must be made numerical.

\begin{bgconcept}
Suppose a customer says, ``Make the board close to 10 feet long.'' A builder needs a tolerance: within one inch? one millimeter? The symbol \(\eps\) names the allowed output error. The symbol \(\delta\) tells how close the input must be to guarantee that error.
\end{bgconcept}

The distance between \(x\) and \(a\) is
\[
  |x-a|.
\]
The distance between \(f(x)\) and \(L\) is
\[
  |f(x)-L|.
\]

\begin{bgdefinition}[title={Formal \(\eps\)-\(\delta\) Definition}]
We say
\[
  \lim_{x\to a}f(x)=L
\]
if
\[
  \boxed{\forall\eps>0,\ \exists\delta>0\text{ such that }
  0<|x-a|<\delta\implies|f(x)-L|<\eps.}
\]
\end{bgdefinition}

In words:

\begin{quote}
For every positive output tolerance \(\eps\), there is a positive input tolerance \(\delta\) such that every allowed input within \(\delta\) of \(a\), except \(a\) itself, produces an output within \(\eps\) of \(L\).
\end{quote}

\begin{center}
\begin{tabularx}{0.9\textwidth}{@{}p{0.24\textwidth}X@{}}
\toprule
\textbf{Expression} & \textbf{Meaning}\\
\midrule
\(|x-a|<\delta\) & The input lies in the horizontal interval \((a-\delta,a+\delta)\).\\
\(|f(x)-L|<\eps\) & The output lies in the vertical band \((L-\eps,L+\eps)\).\\
\(0<|x-a|\) & The target input itself is excluded.\\
\(\forall\eps>0\) & The argument must work for every requested accuracy.\\
\(\exists\delta>0\) & We may choose an input tolerance depending on \(\eps\).\\
\bottomrule
\end{tabularx}
\end{center}
\webnext{calculus/limits/epsilon-delta-linear}{Epsilon-Delta Proofs for Linear Functions}
\end{webpage}

\begin{webpage}{calculus/limits/epsilon-delta-linear}{Epsilon-Delta Proofs for Linear Functions}{lesson}
\section{Linear Functions: The Cleanest Proofs}

\lessonobjective{
Construct a \(\delta\) directly from a requested \(\eps\) for a linear function.
}

\begin{bgguided}[title={\(f(x)=2x\)}]
Prove
\[
  \lim_{x\to3}2x=6.
\]

\begin{bgsolution}
We want the output error to be less than \(\eps\):
\[
  |2x-6|<\eps.
\]
Factor:
\[
  |2x-6|=2|x-3|.
\]
So it is enough to require
\[
  2|x-3|<\eps.
\]
Divide by \(2\):
\[
  |x-3|<\frac\eps2.
\]
Choose
\[
  \boxed{\delta=\frac\eps2}.
\]
Then if \(0<|x-3|<\delta\),
\[
  |2x-6|=2|x-3|<2\delta=2\left(\frac\eps2\right)=\eps.
\]
Therefore,
\[
  \lim_{x\to3}2x=6.
\]
\end{bgsolution}
\end{bgguided}


% --- BEGIN INLINE checks/epsilon-flow-01 ---
\begin{bgcheck}{epsilon-flow-01}{expression}{epsilon/2}[Check your understanding]
For \(f(x)=2x\) at \(x=3\), choose \(\delta\) in terms of \(\varepsilon\).
\begin{bghint}
\(|2x-6|=2|x-3|\).
\end{bghint}
\begin{bgreveal}
Choose \(\delta=\varepsilon/2\).
\end{bgreveal}
\end{bgcheck}

% --- END INLINE checks/epsilon-flow-01 ---


This proof has two phases.

\begin{description}
  \item[Discovery phase.] Start from the desired output inequality and work backward to guess a useful \(\delta\).
  \item[Proof phase.] State the chosen \(\delta\), assume the input inequality, and work forward to the output inequality.
\end{description}

\begin{bgexample}[title={A general linear function}]
Prove
\[
  \lim_{x\to2}(3x+1)=7.
\]

\begin{bgsolution}
\textbf{Discovery.} We need
\[
  |(3x+1)-7|<\eps.
\]
Simplify:
\[
  |3x-6|=3|x-2|.
\]
Thus, requiring
\[
  |x-2|<\frac\eps3
\]
is enough.

\textbf{Proof.} Let \(\eps>0\). Choose
\[
  \delta=\frac\eps3.
\]
If
\[
  0<|x-2|<\delta,
\]
then
\[
  |(3x+1)-7|=3|x-2|<3\delta=\eps.
\]
Therefore,
\[
  \boxed{\lim_{x\to2}(3x+1)=7}.
\]
\end{bgsolution}
\end{bgexample}

\begin{bgsummary}
For \(f(x)=mx+b\), a standard choice is
\[
  \delta=\frac\eps{|m|}
\]
when \(m\ne0\). The output changes \(|m|\) times as much as the input, so the input tolerance must be \(|m|\) times smaller.
\end{bgsummary}
\webnext{calculus/limits/epsilon-delta-quadratic}{Epsilon-Delta Proofs for Quadratic Functions}
\end{webpage}

\begin{webpage}{calculus/limits/epsilon-delta-quadratic}{Epsilon-Delta Proofs for Quadratic Functions}{lesson}
\section{Nonlinear Functions: Controlling an Extra Factor}

\lessonobjective{
Use a preliminary restriction such as \(|x-a|<1\) to bound a factor that depends on \(x\); choose \(\delta\) as the minimum of two requirements.
}

For a quadratic, factoring the output error creates an extra factor involving \(x\).

\begin{bgexample}[title={Prove \(\lim_{x\to2}x^2=4\)}]
\begin{bgsolution}
We need to make
\[
  |x^2-4|
\]
small. Factor:
\[
  |x^2-4|=|x-2||x+2|.
\]
The first factor is controlled by \(\delta\), but \(|x+2|\) still depends on \(x\). First impose the simple restriction
\[
  |x-2|<1.
\]
Then
\[
  1<x<3,
\]
so
\[
  3<x+2<5,
\]
and therefore
\[
  |x+2|<5.
\]
Now
\[
  |x^2-4|=|x-2||x+2|<5|x-2|.
\]
To make this less than \(\eps\), require
\[
  |x-2|<\frac\eps5.
\]
We need both restrictions, so choose
\[
  \boxed{\delta=\min\left\{1,\frac\eps5\right\}}.
\]

Now write the forward proof. Let \(\eps>0\), and choose the \(\delta\) above. If
\[
  0<|x-2|<\delta,
\]
then \(|x-2|<1\), so \(|x+2|<5\). Also, \(|x-2|<\eps/5\). Hence
\[
  |x^2-4|
  =|x-2||x+2|
  <5\left(\frac\eps5\right)
  =\eps.
\]
Therefore,
\[
  \boxed{\lim_{x\to2}x^2=4}.
\]
\end{bgsolution}
\end{bgexample}

\begin{bgconcept}
The minimum notation means, ``Use whichever restriction is smaller.'' If \(\eps/5\) is tiny, use it. If \(\eps/5\) is larger than \(1\), keep the local bound \(1\). Both jobs must be done at once.
\end{bgconcept}

\begin{bgexample}[title={A shifted square}]
Prove
\[
  \lim_{x\to1}x^2=1.
\]

\begin{bgsolution}
Factor the output error:
\[
  |x^2-1|=|x-1||x+1|.
\]
If \(|x-1|<1\), then \(0<x<2\), so \(|x+1|<3\). Therefore,
\[
  |x^2-1|<3|x-1|.
\]
Choose
\[
  \boxed{\delta=\min\left\{1,\frac\eps3\right\}}.
\]
Then
\[
  |x^2-1|<3\delta\le\eps.
\]
\end{bgsolution}
\end{bgexample}
\webnext{calculus/limits/epsilon-delta-graph}{Reading Epsilon and Delta From a Graph}
\end{webpage}

\begin{webpage}{calculus/limits/epsilon-delta-graph}{Reading Epsilon and Delta From a Graph}{lesson}
\section{Reading the Definition From a Graph}

\begin{figure}[htbp]
\centering
\begin{tikzpicture}
\begin{axis}[
  width=0.84\textwidth,height=0.50\textwidth,
  xmin=-0.5,xmax=4.5,ymin=-0.5,ymax=8.5,
  axis lines=middle,xlabel={$x$},ylabel={$f(x)$},
  grid=major,major grid style={draw=BGLine}
]
\addplot[BGBlue,very thick,domain=-0.3:4.2,samples=140]{x^2/2+1};
\addplot[draw=none,name path=upper,domain=-0.5:4.5]{5.5};
\addplot[draw=none,name path=lower,domain=-0.5:4.5]{3.5};
\addplot[BGPaleGold] fill between[of=upper and lower];
\draw[BGGold,thick,dashed] (axis cs:-0.5,5.5)--(axis cs:4.5,5.5);
\draw[BGGold,thick,dashed] (axis cs:-0.5,3.5)--(axis cs:4.5,3.5);
\draw[BGGreen,thick,dashed] (axis cs:1.75,-0.5)--(axis cs:1.75,8.5);
\draw[BGGreen,thick,dashed] (axis cs:2.25,-0.5)--(axis cs:2.25,8.5);
\node[anchor=west,font=\small,BGGold] at (axis cs:3.65,5.5) {$L+\eps$};
\node[anchor=west,font=\small,BGGold] at (axis cs:3.65,3.5) {$L-\eps$};
\node[anchor=north,font=\small,BGGreen] at (axis cs:1.75,-0.2) {$a-\delta$};
\node[anchor=north,font=\small,BGGreen] at (axis cs:2.25,-0.2) {$a+\delta$};
\end{axis}
\end{tikzpicture}
\caption{An \(\eps\)-band around \(L\) and a corresponding \(\delta\)-window around \(a\). A valid \(\delta\) keeps the nearby graph inside the band.}
\label{fig:epsilon-delta}
\end{figure}

\begin{bggraphspec}
Use a continuous nonlinear function such as \(f(x)=x^2/2+1\) near \(a=2\), where \(L=3\). Provide an adjustable \(\eps\)-band and show the largest symmetric \(\delta\)-window whose graph segment remains inside the band. Display \(|x-a|<\delta\) and \(|f(x)-L|<\eps\) simultaneously.
\end{bggraphspec}
\webnext{calculus/limits/disprove-limit-epsilon-delta}{How to Disprove a Claimed Limit}
\end{webpage}

\begin{webpage}{calculus/limits/disprove-limit-epsilon-delta}{How to Disprove a Claimed Limit}{lesson}
\section{Proving a Limit Is Not a Claimed Number}

To show
\[
  \lim_{x\to a}f(x)\ne L,
\]
it is enough to find one positive tolerance \(\eps_0\) for which no possible \(\delta\) works.

\begin{bgexample}[title={A jump fails the formal definition}]
Let
\[
  f(x)=
  \begin{cases}
    0,&x<0,\\
    1,&x\ge0.
  \end{cases}
\]
Show that \(\lim_{x\to0}f(x)\) does not equal \(1/2\).

\begin{bgsolution}
Choose
\[
  \eps_0=\frac14.
\]
The band around \(1/2\) is
\[
  \left(\frac14,\frac34\right).
\]
But every left-side input near zero has output \(0\), which lies outside this band, and every right-side input has output \(1\), also outside the band.

No matter how small \(\delta>0\) is chosen, the point \(x=-\delta/2\) satisfies
\[
  0<|x|<\delta,
\]
but
\[
  \left|f(x)-\frac12\right|
  =\left|0-\frac12\right|
  =\frac12
  >\frac14.
\]
Therefore the claimed limit cannot be \(1/2\). In fact, the left and right limits disagree, so no two-sided limit exists.
\end{bgsolution}
\end{bgexample}
\webnext{calculus/limits/unit-review}{Limits and Continuity Unit Review}
\end{webpage}

\begin{webpage}{calculus/limits/formal-infinite-limits}{Formal Definition of an Infinite Limit}{extension}
\section{Formal Infinite Limits: Optional Extension}

The statement
\[
  \lim_{x\to a}f(x)=+\infty
\]
can be formalized as follows:

\begin{bgdefinition}[title={Formal Positive Infinite Limit}]
For every \(M>0\), there exists \(\delta>0\) such that
\[
  0<|x-a|<\delta\implies f(x)>M.
\]
\end{bgdefinition}

Instead of asking for an output band around a finite \(L\), the challenger names an arbitrarily high threshold \(M\). The proof must force the function above it.

\begin{bgguided}[title={\(1/x^2\to+\infty\)}]
Show formally that
\[
  \lim_{x\to0}\frac1{x^2}=+\infty.
\]

\begin{bgsolution}
Let \(M>0\). We want
\[
  \frac1{x^2}>M.
\]
This is equivalent to
\[
  x^2<\frac1M,
\]
which is guaranteed by
\[
  |x|<\frac1{\sqrt M}.
\]
Choose
\[
  \delta=\frac1{\sqrt M}.
\]
Then \(0<|x|<\delta\) implies
\[
  x^2<\delta^2=\frac1M,
\]
so
\[
  \frac1{x^2}>M.
\]
\end{bgsolution}
\end{bgguided}
\webnext{calculus/limits/unit-review}{Limits and Continuity Unit Review}
\end{webpage}

\begin{webpage}{calculus/limits/epsilon-delta-review}{Epsilon-Delta Review}{review}
\section*{Section 6 Summary}
\addcontentsline{toc}{section}{Section 6 Summary}

\begin{bgsummary}
\begin{itemize}
  \item \(\eps\) is an allowed output error; \(\delta\) is an input tolerance that guarantees it.
  \item Formal proofs work backward to discover \(\delta\) and forward to verify it.
  \item Linear proofs usually use \(\delta=\eps/|m|\).
  \item Nonlinear proofs often need a preliminary local bound and a minimum of two restrictions.
  \item To disprove a claimed limit, one bad \(\eps_0\) is enough.
\end{itemize}
\end{bgsummary}
\webnext{calculus/limits/unit-review}{Limits and Continuity Unit Review}
\end{webpage}

\begin{webpage}{calculus/limits/epsilon-delta-concept-quiz}{Epsilon-Delta Concept Quiz}{quiz}
\section*{Section 6 Concept Quiz}
\addcontentsline{toc}{section}{Section 6 Concept Quiz}

Use this as a short retrieval check after finishing the section. On the website, each item should accept an answer before revealing the hint and worked explanation.

\begin{bgcheck}{epsilon-linear-q1}{expression}{epsilon/3}[Concept check]
For \(f(x)=3x+1\) at \(x=2\), give a working choice of \(\delta\) in terms of \(\varepsilon\).
\begin{bghint}
\(|f(x)-7|=3|x-2|\).
\end{bghint}
\begin{bgreveal}
Choose \(\delta=\varepsilon/3\).
\end{bgreveal}
\end{bgcheck}

\begin{bgcheck}{epsilon-linear-q2}{expression}{epsilon/5}[Concept check]
For \(f(x)=5x\) at \(x=1\), give a working \(\delta\).
\begin{bghint}
\(|5x-5|=5|x-1|\).
\end{bghint}
\begin{bgreveal}
Choose \(\delta=\varepsilon/5\).
\end{bgreveal}
\end{bgcheck}

\begin{bgcheck}{epsilon-quadratic-q1}{expression}{min(1,epsilon/5)}[Concept check]
For the standard proof that \(x^2\to4\) as \(x\to2\), enter the usual safe choice of \(\delta\).
\begin{bghint}
First require \(|x-2|<1\), which makes \(|x+2|<5\).
\end{bghint}
\begin{bgreveal}
Choose \(\delta=\min\{1,\varepsilon/5\}\).
\end{bgreveal}
\end{bgcheck}

\begin{bgcheck}{formal-reading-q1}{choice}{output}[Concept check]
In \(|f(x)-L|<\varepsilon\), does \(\varepsilon\) control input distance or output distance?
\begin{bghint}
Look at which quantities are being compared.
\end{bghint}
\begin{bgreveal}
It controls output distance.
\end{bgreveal}
\end{bgcheck}

\begin{bgsummary}[title={Quiz use on the website}]
This page should report completion by concept, not merely a total score. A wrong answer should link back to the exact lesson that teaches the missing method. The same questions may appear in randomized numerical variants, but the canonical explanation should remain stable.
\end{bgsummary}
\webnext{calculus/limits/unit-review}{Limits and Continuity Unit Review}
\end{webpage}

\begin{webpage}{calculus/limits/epsilon-delta-practice}{Epsilon-Delta Practice Problems}{practice}
\section*{Section 6 Exercises}
\addcontentsline{toc}{section}{Section 6 Exercises}

\begin{bgexercises}[label=\textbf{6.\arabic*}]
  \item In your own words, explain the roles of \(\eps\) and \(\delta\).
  \item Rewrite \(|x-3|<0.2\) as an ordinary interval.
  \item Rewrite \(|f(x)-5|<0.1\) as a vertical output interval.
  \item For \(f(x)=4x\), find a \(\delta\) in terms of \(\eps\) that proves \(\lim_{x\to2}f(x)=8\).
  \item Prove \(\lim_{x\to1}(5x-2)=3\).
  \item Prove \(\lim_{x\to-2}(3x+4)=-2\).
  \item Prove \(\lim_{x\to4}(x/2)=2\).
  \item Prove \(\lim_{x\to1}x^2=1\) using \(\delta=\min\{1,\eps/3\}\).
  \item Prove \(\lim_{x\to3}x^2=9\) by first bounding \(|x+3|\).
  \item Find a suitable \(\delta\) to prove \(\lim_{x\to2}(x^2+x)=6\).
  \item Explain why a proof may choose a smaller \(\delta\) than necessary.
  \item Explain why \(\delta\) may depend on \(\eps\), but not on the particular \(x\) chosen later.
  \item Use \(\eps_0=1/4\) to show the step function in the section does not approach \(1/2\) at zero.
  \item Give a formal \(M\)-\(\delta\) proof that \(1/x^2\to+\infty\) as \(x\to0\).
  \item Give a formal \(N\)-style proof that \(1/x\to0\) as \(x\to\infty\).
  \item Identify the first incorrect step in the claim: ``Choose \(\delta=\eps\). Then \(|x^2-4|<\eps\) whenever \(|x-2|<\delta\).''
  \item Why is \(0<|x-a|\) included in the formal finite-limit definition?
  \item Explain how the graphical \(\eps\)-band and \(\delta\)-window represent the quantified definition.
\end{bgexercises}

Answers begin in \cref{app:chapter-answers}.
\webnext{calculus/limits/unit-review}{Limits and Continuity Unit Review}
\end{webpage}


% --- END INLINE chapters/ch06_formal ---


% --- BEGIN INLINE chapters/ch07_review ---
\begin{webpage}{calculus/limits/unit-review}{Limits and Continuity Unit Review}{review}
\chapter{Mastery Review and Practice Exams}

\section{The Complete Limit Decision Process}

A strong student does not merely know many techniques. A strong student can decide which technique applies when the problem is not labeled.

\begin{bgmethod}[title={Master Decision Tree}]
When evaluating a limit:
\begin{enumerate}
  \item \textbf{Read the direction.} Is the input approaching a finite number, one side of a finite number, \(+\infty\), or \(-\infty\)?
  \item \textbf{Substitute when meaningful.}
  \begin{itemize}
    \item Real number: usually finished.
    \item \(0/0\): algebraic indeterminate form.
    \item Nonzero over zero: one-sided sign analysis.
    \item \(\infty/\infty\): compare powers or divide by the dominant power.
    \item \(\infty-\infty\): combine or rationalize.
  \end{itemize}
  \item \textbf{Inspect the structure.}
  \begin{itemize}
    \item polynomial factors \(\to\) factor and cancel;
    \item radicals \(\to\) conjugate;
    \item fractions inside fractions \(\to\) common denominator;
    \item absolute values/piecewise rules \(\to\) one-sided cases;
    \item bounded oscillation times zero \(\to\) squeeze;
    \item trig near zero \(\to\) standard trig limits;
    \item rational function at infinity \(\to\) degrees/dominant terms;
    \item radical at negative infinity \(\to\) use \(|x|\).
  \end{itemize}
  \item \textbf{State the conclusion precisely.} Distinguish a finite number, \(\DNE\), \(+\infty\), and \(-\infty\).
\end{enumerate}
\end{bgmethod}
\end{webpage}

\begin{webpage}{calculus/limits/common-exam-errors}{Common Limits Exam Errors}{review}
\section{Common Exam Errors and Their Repairs}

\begin{center}
\begin{longtable}{@{}p{0.30\textwidth}p{0.31\textwidth}p{0.31\textwidth}@{}}
\toprule
\textbf{Error} & \textbf{Why it fails} & \textbf{Repair}\\
\midrule
\endhead
Using the filled point as the limit & The point gives \(f(a)\), not nearby behavior & Trace the graph from both sides first\\
Calling \(0/0\) the answer & It is an indeterminate form & Simplify, then substitute again\\
Cancelling terms across addition & Cancellation applies to factors & Factor the entire numerator or denominator\\
Ignoring left and right sides & Two-sided limits require agreement & Compute both one-sided limits\\
Writing \(\sqrt{x^2}=x\) at negative infinity & The principal square root is nonnegative & Use \(\sqrt{x^2}=|x|=-x\) for \(x<0\)\\
Using degrees in \(\sin x/x\) & The fundamental limit is a radian statement & Convert or use radian measure\\
Applying IVT without continuity & Jumps can skip intermediate values & State continuity on the entire closed interval\\
Claiming IVT gives the exact root & It gives existence only & Use bisection or another numerical method for location\\
Using l'Hospital's Rule too early & It may be unavailable, circular, or hide the needed algebra & Use first-unit methods requested by the course\\
\bottomrule
\end{longtable}
\end{center}
\webnext{calculus/limits/unit-review}{Limits and Continuity Unit Review}
\end{webpage}

\begin{webpage}{calculus/study/correct-practice-test}{How to Correct a Calculus Practice Test}{study}
\section{How to Correct a Practice Test}

A practice test is useful only if corrections identify the cause of each error.

For every missed problem, record:

\begin{enumerate}
  \item the first incorrect line, not merely the final answer;
  \item the category of error: concept, method selection, algebra, sign, notation, or time pressure;
  \item the smallest skill that would have prevented it;
  \item one new problem of the same type, worked correctly without looking back.
\end{enumerate}

\begin{bgconcept}
``I made a silly mistake'' is not a diagnosis. A sign error near a vertical asymptote, a failed factorization, and forgetting that one-sided limits must agree are different problems. Name the species before trying to exterminate it.
\end{bgconcept}
\webnext{calculus/limits/unit-review}{Limits and Continuity Unit Review}
\end{webpage}

\begin{webpage}{calculus/limits/cumulative-practice}{Cumulative Limits and Continuity Practice}{practice}
\section{Cumulative Review Set}

Work this set without section labels telling you the method. Unless a calculator is explicitly permitted, use exact values.

\subsection*{Part A: Concepts and graphs}
\begin{bgexercises}[label=\textbf{R\arabic*.}]
  \item Explain the difference between \(f(a)\) and \(\lim_{x\to a}f(x)\).
  \item State the one-sided criterion for a two-sided limit.
  \item List four reasons a limit may fail to exist.
  \item State the three continuity conditions at \(a\).
  \item Explain why changing one value can repair a hole but not a jump.
  \item Explain why infinity is not treated as an ordinary real-number limit.
  \item State the Squeeze Theorem.
  \item State the fundamental sine limit and its unit requirement.
  \item State the Intermediate Value Theorem, including every hypothesis.
  \item Explain what IVT does not guarantee.
\end{bgexercises}

\subsection*{Part B: Finite limits}
\begin{bgexercises}[label=\textbf{R\arabic*.},resume]
  \item \(\displaystyle\lim_{x\to2}(3x^2-x+4)\)
  \item \(\displaystyle\lim_{x\to-1}\frac{x^2+2}{x+4}\)
  \item \(\displaystyle\lim_{x\to4}\frac{x^2-16}{x-4}\)
  \item \(\displaystyle\lim_{x\to-2}\frac{x^3+8}{x+2}\)
  \item \(\displaystyle\lim_{x\to1}\frac{x^4-1}{x^2-1}\)
  \item \(\displaystyle\lim_{x\to0}\frac{\sqrt{x+25}-5}{x}\)
  \item \(\displaystyle\lim_{x\to9}\frac{x-9}{\sqrt{x}-3}\)
  \item \(\displaystyle\lim_{x\to0}\frac{\frac1{x+2}-\frac12}{x}\)
  \item \(\displaystyle\lim_{x\to3}\frac{|x-3|}{x-3}\)
  \item \(\displaystyle\lim_{x\to2}\frac{|x-2|}{\sqrt{|x-2|}}\)
\end{bgexercises}

\subsection*{Part C: Squeeze and trigonometry}
\begin{bgexercises}[label=\textbf{R\arabic*.},resume]
  \item \(\displaystyle\lim_{x\to0}x^3\sin(1/x)\)
  \item \(\displaystyle\lim_{x\to0}x^2\cos(4/x)\)
  \item \(\displaystyle\lim_{x\to0}\frac{\sin(6x)}x\)
  \item \(\displaystyle\lim_{x\to0}\frac{\sin(2x)}{\sin(9x)}\)
  \item \(\displaystyle\lim_{x\to0}\frac{\tan(5x)}{3x}\)
  \item \(\displaystyle\lim_{x\to0}\frac{1-\cos(4x)}{x^2}\)
  \item \(\displaystyle\lim_{x\to0}\frac{\sin(3x)\sin(7x)}{x^2}\)
  \item \(\displaystyle\lim_{x\to0}\frac{1-\cos x}{x\sin x}\)
\end{bgexercises}

\subsection*{Part D: Infinite behavior}
\begin{bgexercises}[label=\textbf{R\arabic*.},resume]
  \item \(\displaystyle\lim_{x\to1^-}\frac{x+2}{x-1}\)
  \item \(\displaystyle\lim_{x\to1^+}\frac{x+2}{x-1}\)
  \item \(\displaystyle\lim_{x\to-2}\frac{-3}{(x+2)^2}\)
  \item Find every hole and vertical asymptote of \(\dfrac{x^2-1}{x^2-3x+2}\).
  \item \(\displaystyle\lim_{x\to\infty}\frac{4x^2-x}{2x^2+7}\)
  \item \(\displaystyle\lim_{x\to-\infty}\frac{3x+1}{x^2+5}\)
  \item \(\displaystyle\lim_{x\to\infty}\frac{x^3+1}{x^2-1}\)
  \item Find the slant asymptote of \(\dfrac{x^2+2}{x-1}\).
  \item \(\displaystyle\lim_{x\to\infty}\frac{\sqrt{4x^2+1}}x\)
  \item \(\displaystyle\lim_{x\to-\infty}\frac{\sqrt{4x^2+1}}x\)
  \item \(\displaystyle\lim_{x\to\infty}(\sqrt{x^2+12x}-x)\)
  \item \(\displaystyle\lim_{x\to-\infty}(\sqrt{x^2+8x}+x)\)
\end{bgexercises}

\subsection*{Part E: Continuity and the IVT}
\begin{bgexercises}[label=\textbf{R\arabic*.},resume]
  \item Find intervals of continuity of \(\sqrt{x+1}/(x-2)\).
  \item Classify all discontinuities of \((x^2-4)/(x^2-x-2)\).
  \item Find \(c\) so \(\begin{cases}(x^2-9)/(x-3),&x\ne3,\\c,&x=3\end{cases}\) is continuous.
  \item Find \(k\) so \(\begin{cases}kx+2,&x<1,\\x^2+4,&x\ge1\end{cases}\) is continuous at \(1\).
  \item Determine whether any \(a\) makes \(\begin{cases}ax+1,&x<0,\\ax+4,&x\ge0\end{cases}\) continuous at \(0\).
  \item Show that \(x^3-x-1=0\) has a root in \((1,2)\).
  \item Explain why \(1/x\) cannot use IVT on \([-1,1]\).
  \item Use two bisection steps to narrow a root of \(x^3-x-1\) from \((1,2)\).
\end{bgexercises}

\subsection*{Part F: Formal limits}
\begin{bgexercises}[label=\textbf{R\arabic*.},resume]
  \item Find a \(\delta\) in terms of \(\eps\) proving \(\lim_{x\to2}(4x-1)=7\).
  \item Prove \(\lim_{x\to1}x^2=1\) using a local bound.
  \item Explain why \(\delta=\eps\) does not automatically work for every function.
  \item Translate \(|x-5|<0.01\) and \(|f(x)-3|<0.02\) into interval language.
\end{bgexercises}
\webnext{calculus/limits/unit-review}{Limits and Continuity Unit Review}
\end{webpage}

\begin{webpage}{calculus/limits/practice-exam-a}{Limits and Continuity Practice Exam A}{exam}
\section{Practice Examination A}

\textbf{Suggested time:} 75 minutes. \textbf{Calculator:} none unless your course normally permits one. Show reasoning; unsupported answers may receive little credit.

\subsection*{Part I: Concepts, 20 points}
\begin{bgexercises}[label=\textbf{A\arabic*.}]
  \item A graph approaches \(4\) from both sides at \(x=2\), but the filled point is \((2,7)\). State \(f(2)\) and the limit. Explain the distinction.
  \item State all three continuity conditions at \(x=a\).
  \item Explain why \(0/0\) is indeterminate rather than equal to zero.
  \item State what the Intermediate Value Theorem guarantees and what it does not guarantee.
\end{bgexercises}

\subsection*{Part II: Finite limits, 35 points}
\begin{bgexercises}[label=\textbf{A\arabic*.},resume]
  \item \(\displaystyle\lim_{x\to-2}(2x^3-x+1)\)
  \item \(\displaystyle\lim_{x\to3}\frac{x^2-9}{x-3}\)
  \item \(\displaystyle\lim_{x\to0}\frac{\sqrt{x+9}-3}{x}\)
  \item \(\displaystyle\lim_{x\to2}\frac{\frac1x-\frac12}{x-2}\)
  \item \(\displaystyle\lim_{x\to0}\frac{\sin(7x)}{2x}\)
  \item \(\displaystyle\lim_{x\to0}x^2\cos(1/x)\)
  \item \(\displaystyle\lim_{x\to1}\frac{|x-1|}{x-1}\)
\end{bgexercises}

\subsection*{Part III: Infinite behavior, 25 points}
\begin{bgexercises}[label=\textbf{A\arabic*.},resume]
  \item Find both one-sided limits of \(\dfrac{x+1}{x-2}\) at \(x=2\).
  \item \(\displaystyle\lim_{x\to-\infty}\frac{5x^2+1}{2x^2-3x}\)
  \item \(\displaystyle\lim_{x\to-\infty}\frac{\sqrt{x^2+4}}x\)
  \item Find the hole, vertical asymptote, and horizontal asymptote of \(\dfrac{x^2-1}{x^2-3x+2}\).
\end{bgexercises}

\subsection*{Part IV: Continuity and existence, 20 points}
\begin{bgexercises}[label=\textbf{A\arabic*.},resume]
  \item Find \(k\) so \(\begin{cases}kx+1,&x<2,\\x^2-1,&x\ge2\end{cases}\) is continuous at \(2\).
  \item Use IVT to show \(x^3+x-3=0\) has a root in \((1,2)\).
  \item Give a valid \(\delta\) in terms of \(\eps\) for \(\lim_{x\to3}(2x+1)=7\).
\end{bgexercises}
\webnext{calculus/limits/unit-review}{Limits and Continuity Unit Review}
\end{webpage}

\begin{webpage}{calculus/limits/practice-exam-b}{Limits and Continuity Practice Exam B}{exam}
\section{Practice Examination B}

This exam is deliberately less patterned. It is meant to test whether you can identify methods rather than imitate section headings.

\textbf{Suggested time:} 90 minutes.

\begin{bgexercises}[label=\textbf{B\arabic*.}]
  \item Let
  \[
    f(x)=\begin{cases}
      x^2+1,&x<1,\\
      5,&x=1,\\
      3x-1,&x>1.
    \end{cases}
  \]
  Find both one-sided limits, the two-sided limit, \(f(1)\), and determine continuity at \(1\).
  \item \(\displaystyle\lim_{x\to2}\frac{x^3-8}{x^2-4}\)
  \item \(\displaystyle\lim_{x\to0}\frac{\sqrt{1+4x}-\sqrt{1-x}}x\)
  \item \(\displaystyle\lim_{x\to0}\frac{1-\cos(3x)}{x^2}\)
  \item \(\displaystyle\lim_{x\to0}\frac{\tan(2x)}{\sin(5x)}\)
  \item \(\displaystyle\lim_{x\to0}x\sin(1/x^2)\)
  \item Find every discontinuity of \(\dfrac{(x-1)(x+2)}{(x-1)(x-3)^2}\), classify it, and state all one-sided infinite limits.
  \item \(\displaystyle\lim_{x\to\infty}\frac{2x^3-x}{x^2+1}\)
  \item Find the polynomial asymptote of the function in the previous problem.
  \item \(\displaystyle\lim_{x\to-\infty}(\sqrt{x^2+10x}+x)\)
  \item Find \(a,b\) so
  \[
    g(x)=\begin{cases}
      ax+b,&x<1,\\
      x^2+2,&1\le x<2,\\
      3x+a,&x\ge2
    \end{cases}
  \]
  is continuous at \(1\) and \(2\).
  \item Show that \(x^5-3x+1=0\) has a root in \((0,1)\), then perform two bisection steps.
  \item Prove \(\lim_{x\to1}(4x-3)=1\) using the formal definition.
  \item A student claims that because \(f(-1)<0<f(1)\), every function has a root in \((-1,1)\). Give a counterexample and state the missing hypothesis.
\end{bgexercises}

Full answers and selected worked solutions are in the appendices. Do not read them until you have completed an honest attempt. The universe already contains enough answer keys masquerading as education.
\webnext{calculus/limits/unit-review}{Limits and Continuity Unit Review}
\end{webpage}


% --- END INLINE chapters/ch07_review ---

\appendix

% --- BEGIN INLINE appendices/answers ---
\chapter{Answers and Selected Solutions}

\section{Diagnostic Answers and Skill Map}
\label{app:diagnostic-answers}

\begin{enumerate}
  \item \((x-5)(x+5)\).
  \item \((x-2)(x^2+2x+4)\).
  \item \((2x+1)(x-3)\).
  \item \(x+3\), with \(x\ne3\).
  \item \(-1/[2(x+2)]\), with the relevant exclusions inherited from the original expression.
  \item \(1/(\sqrt{x+9}+3)\).
  \item \(3.7<x<4.3\).
  \item \(|x|\).
  \item \(x\ne-2,2\); domain \(( -\infty,-2)\cup(-2,2)\cup(2,\infty)\).
  \item \(x\le5\).
  \item \(1/2,1/2,1\).
  \item \(\sin^2\theta\).
  \item \(f(2+h)=h^2+h-2\).
  \item \(h+1\), for \(h\ne0\).
  \item \(x=3\).
  \item \(a=1\).
  \item \(2\).
  \item \(x=2\) is equality; \(x\to2\) describes nearby values approaching \(2\).
  \item Division by zero is undefined; no number multiplied by zero can recover the required numerator in a unique way.
  \item Any sketch with an open circle at \((1,3)\) and a filled circle at \((1,5)\).
\end{enumerate}

\begin{center}
\begin{tabularx}{\textwidth}{@{}p{0.26\textwidth}X@{}}
\toprule
\textbf{Missed questions} & \textbf{Review before continuing}\\
\midrule
1--4 & Factoring and rational-expression simplification\\
5--6 & Complex fractions and conjugates\\
7--8 & Absolute value and distance notation\\
9--10 & Domains\\
11--12 & Unit-circle trigonometry and identities\\
13--14 & Function substitution and difference quotients\\
15--17 & Solving equations and slope\\
18--20 & Limit language and graph notation\\
\bottomrule
\end{tabularx}
\end{center}

\section{Section Exercise Answers}
\label{app:chapter-answers}

\subsection*{Section 1}

\begin{enumerate}[label=\textbf{1.\arabic*}]
  \item Input target \(4\); output target \(9\).
  \item \(\lim_{t\to2}g(t)=-5\).
  \item ``The limit of \(Q(h)\) as \(h\) approaches zero is \(12\).''
  \item False. Continuity would be needed to guarantee equality.
  \item False. For example, \((x^2-4)/(x-2)\) is undefined at \(2\) but has limit \(4\).
  \item Approach using inputs less than \(3\).
  \item Approach using inputs greater than \(3\).
  \item It describes a process through nearby values, not equality at the target.
  \item \(7\).
  \item \(f(-2)\) is undefined; the limit is \(4\).
  \item \(5\) ft/s.
  \item Average velocity \(6+h\); instantaneous velocity \(6\).
  \item Average velocity \(16-16h\); instantaneous velocity \(16\) ft/s.
  \item Average rate \(66+2h\); limiting rate \(66\) people per time unit.
  \item The slope of a secant line.
  \item A tangent line, when the limiting slope exists.
  \item \(6\).
  \item \(12\).
  \item \(5\).
  \item \(p(1)\) is undefined; the limit is \(2\).
  \item \(q(1)=100\); the limit is \(2\).
  \item Example: \(f(x)=x\) with \(f(0)=0\), and \(g(x)=x\) for \(x\ne0\) with \(g(0)=7\). Both limits are zero.
  \item Left \(3\), right \(4\), so \(\DNE\).
  \item Left \(4\), right \(5\), so \(\DNE\).
  \item Both sides equal \(4\); limit \(4\).
  \item Left \(-1\), right \(2\); limit \(\DNE\).
  \item Left \(-1\), right \(1\).
  \item Left \(-1\), right \(1\).
  \item Opposite unbounded one-sided behavior.
  \item Endless oscillation.
  \item The right-hand limit is \(0\); a real two-sided limit is unavailable because the function has no domain immediately left of zero.
  \item The left- and right-hand behavior is missing; the filled dot alone gives only \(f(2)\).
  \item Example: \((x^2-1)/(x-1)\) at \(x=1\).
  \item One construction: \(f(x)=2\) for \(x<1\), \(f(1)=-3\), and \(f(x)=5\) for \(x>1\).
  \item Example: \(f(x)=(x^2-16)/(x-4)+1\) for \(x\ne4\). Then the limit is \(9\).
  \item A finite table checks only finitely many inputs and may miss oscillation or exceptional sequences.
  \item Left \(3\), right \(3\), two-sided limit \(3\), and \(f(2)=4\).
  \item The limit is \(3\). The value \(f(-1)\) may be any real number or may be left undefined without changing the limit.
  \item \(18\).
  \item Open circle at \((0,1)\), filled point at \((0,2)\), nearby graph approaching \(1\) from both sides.
  \item Left approach \(1\), right approach \(3\), filled point at \((0,2)\).
  \item The first has matching one-sided limits; the second does not.
\end{enumerate}

\subsection*{Section 2}

\begin{enumerate}[label=\textbf{2.\arabic*}]
  \item \(13\).
  \item \(13\).
  \item \(56\).
  \item \(4/3\).
  \item \(6/5\).
  \item \(3\).
  \item \(3\).
  \item \(11\).
  \item \(-6\).
  \item \(1/5\).
  \item Division by a limit of zero is not defined by the quotient law.
  \item Answers vary; for example \(\lim_{x\to2}(x^2+1)=5\).
  \item \(10\).
  \item \(-1\).
  \item \(1\).
  \item \(12\).
  \item \(12\).
  \item \(4\).
  \item \(-4\).
  \item \(3\).
  \item \(3/2\).
  \item \(4\).
  \item \(27\).
  \item The expressions agree at every sufficiently close input except possibly the target, which the limit excludes.
  \item \(1/6\).
  \item \(1/8\).
  \item \(10\).
  \item \(1/4\).
  \item \(5/2\).
  \item \(1/2\).
  \item \(6\).
  \item \(3/2\).
  \item \(-1\).
  \item \(-1/4\).
  \item \(-1/9\).
  \item \(-1/4\).
  \item \(-1/4\).
  \item \(-1/a^2\).
  \item \(1\).
  \item \(1\).
  \item \(-1\).
  \item \(1\).
  \item \(\DNE\).
  \item \(\DNE\).
  \item \(-1\).
  \item \(0\).
  \item \(\DNE\).
  \item \(\DNE\).
  \item \(2\).
  \item \(5/2\).
  \item \(1/3\).
  \item \(-1/16\).
  \item \(0\).
  \item Valid cancellation occurs after factoring \(x(x+1)\); terms cannot be cancelled separately across addition.
  \item Example: \(\lim_{x\to1}(x^2+3x-4)/(x-1)=5\).
  \item Example: \(\lim_{x\to0}x^2/x=0\).
  \item Example: \(\lim_{x\to0}|x|/x\) is \(0/0\) under substitution but does not exist.
  \item Substitute, identify the form, then choose structure-specific algebra as summarized in the section decision tree.
\end{enumerate}

\subsection*{Section 3}

\begin{enumerate}[label=\textbf{3.\arabic*}]
  \item \(0\).
  \item \(2\).
  \item \(0\).
  \item \(0\).
  \item \(0\).
  \item \(0\).
  \item \(0\).
  \item A bounded function need not approach zero; for example, \(\sin(1/x)\) remains bounded but has no limit.
  \item \(1\).
  \item \(2\).
  \item \(7/3\).
  \item \(5/2\).
  \item \(3/8\).
  \item \(\pi\).
  \item \(1/4\).
  \item \(3/5\).
  \item \(1\).
  \item \(4\).
  \item \(3/2\).
  \item \(0\).
  \item \(1/2\).
  \item \(8\).
  \item \(1\).
  \item \(9\).
  \item \(6\).
  \item \(5/2\).
  \item \(2\).
  \item \(4\).
  \item \(1/2\).
  \item \(0\).
  \item \(1/2\).
  \item \(3/5\).
  \item Radian measure equals unit-circle arc length; the geometric squeeze produces ratio \(1\) only in radians.
  \item \(\tan x/x=(\sin x/x)(1/\cos x)\to1\).
  \item Rationalize and use \((\sin x/x)^2/(1+\cos x)\to1/2\).
  \item \(\sin(5x)\) is not equal to \(5x\); only their ratio approaches \(1\) near zero.
  \item Use \(-x^4\le x^4\sin(1/x^3)\le x^4\); the limit is zero.
  \item Example: \(\lim_{x\to0}\sin(7x)/(4x)=7/4\).
\end{enumerate}

\subsection*{Section 4}

\begin{enumerate}[label=\textbf{4.\arabic*}]
  \item \(-\infty\).
  \item \(+\infty\).
  \item \(+\infty\).
  \item \(+\infty\).
  \item \(-\infty\).
  \item \(+\infty\).
  \item \(-\infty\).
  \item \(+\infty\).
  \item \(-\infty\).
  \item \(+\infty\).
  \item Hole at \(2\); vertical asymptote at \(-1\).
  \item Hole at \(3\); vertical asymptote at \(-2\).
  \item Hole at \(-1\); vertical asymptote at \(2\).
  \item At \(x=1\): left \(-\infty\), right \(+\infty\). At \(x=-2\): left \(+\infty\), right \(-\infty\).
  \item Both one-sided limits at \(-1\) are \(-\infty\).
  \item Cancellation removes the zero factor for nearby inputs, leaving a finite limiting value.
  \item \(0\).
  \item \(2/5\).
  \item \(7/2\).
  \item \(3\).
  \item \(0\).
  \item \(-\infty\).
  \item \(-2\).
  \item Horizontal asymptote \(y=3\).
  \item Slant asymptote \(y=x+2\).
  \item Polynomial asymptote \(y=2x\).
  \item \(1\).
  \item \(-1\).
  \item \(3\).
  \item \(-3\).
  \item \(4\).
  \item \(1/4\).
  \item \(-3\).
  \item Use \(\sqrt{x^2}=|x|\); at negative infinity \(|x|=-x\).
  \item Hole at \(1\), vertical asymptote \(x=2\), horizontal asymptote \(y=1\); one-sided limits at \(2\): \(-\infty,+\infty\).
  \item Hole at \(2\), vertical asymptote \(x=-1\), horizontal asymptote \(y=1\); one-sided limits at \(-1\): \(-\infty,+\infty\).
  \item At \(-3\): left \(-\infty\), right \(+\infty\). At \(3\): left \(-\infty\), right \(+\infty\).
  \item Polynomial asymptote \(y=2x\).
  \item Example: \(\dfrac{(x-1)(3x+1)}{(x-1)(x+2)}\).
  \item Horizontal asymptotes describe end behavior and may be crossed; a vertical asymptote occurs at an excluded finite input where the function is unbounded.
\end{enumerate}

\subsection*{Section 5}

\begin{enumerate}[label=\textbf{5.\arabic*}]
  \item Yes; value and limit both equal \(5\).
  \item No; the function value is undefined and behavior is unbounded.
  \item Yes.
  \item The equality condition fails.
  \item The function-value condition fails, and therefore the equality condition cannot hold.
  \item No.
  \item A dot may be isolated from the nearby graph.
  \item The value may be missing or different from the limit.
  \item \(( -\infty,-2)\cup(-2,2)\cup(2,\infty)\).
  \item \([1,\infty)\).
  \item \(( -\infty,-2)\cup(-2,5]\).
  \item \((-3,\infty)\).
  \item Break at \(-3\pi/2,-\pi/2,\pi/2,3\pi/2\); continuous on the resulting subintervals.
  \item \([-2,3)\cup(3,\infty)\).
  \item \((-3,3)\).
  \item Continuous on \([-2,2]\), using right continuity at \(-2\) and left continuity at \(2\).
  \item Removable.
  \item Jump.
  \item Infinite.
  \item Oscillatory.
  \item Hole at \(2\); vertical asymptote at \(-1\).
  \item Infinite discontinuity at \(3\) after one factor cancels.
  \item No; nearby one-sided behavior remains unequal.
  \item No; nearby values remain unbounded.
  \item \(4\).
  \item \(6\).
  \item \(3\).
  \item \(1/4\).
  \item No value works because the one-sided limits are \(-1\) and \(1\).
  \item No value works because the function is unbounded.
  \item \(k=2\).
  \item \(a=2\).
  \item \(b=-3\).
  \item \(m=5\).
  \item No value works.
  \item \(a=4,b=-2\).
  \item Answers vary.
  \item The third rule must assign the common limiting value at the join.
  \item Polynomial continuity and signs \(-3,2\) at \(1,2\).
  \item Polynomial continuity and signs \(-2,32\) at \(0,2\).
  \item \(f(x)=\cos x-x\) is continuous; \(f(0)=1\), \(f(1)=\cos1-1<0\).
  \item A function can touch zero and return without an endpoint sign change.
  \item Example: \(f(x)=x^2\) on \([-1,1]\).
  \item \(1/x\) is not continuous on the interval because it is undefined at zero.
  \item After two steps, the root lies in \((0.5,0.75)\).
  \item After three steps, the root lies in \((1.75,1.875)\).
  \item No. Example: \(x^2-1\) has two roots on \([-2,2]\).
  \item State continuity on \([a,b]\), an intermediate target between endpoint outputs, and the existence of \(c\in[a,b]\) with \(f(c)=N\).
\end{enumerate}

\subsection*{Section 6}

\begin{enumerate}[label=\textbf{6.\arabic*}]
  \item \(\eps\) is the allowed output error; \(\delta\) is the input distance guaranteeing it.
  \item \(2.8<x<3.2\).
  \item \(4.9<f(x)<5.1\).
  \item \(\delta=\eps/4\).
  \item \(\delta=\eps/5\).
  \item \(\delta=\eps/3\).
  \item \(\delta=2\eps\).
  \item Use the proof in the section.
  \item One valid choice is \(\delta=\min\{1,\eps/7\}\).
  \item One valid choice is \(\delta=\min\{1,\eps/6\}\).
  \item A smaller positive \(\delta\) preserves the implication.
  \item The proof must guarantee the conclusion uniformly for every eligible \(x\) in the window.
  \item Choose \(\eps_0=1/4\); points on either side remain at distance \(1/2\) from the claimed limit.
  \item Choose \(\delta=1/\sqrt M\).
  \item Given \(\eps>0\), choose \(N>1/\eps\); then \(x>N\) implies \(1/x<\eps\).
  \item The factor \(|x+2|\) was not bounded, so \(|x-2|<\eps\) alone does not control the product.
  \item A limit ignores the target point's value and studies punctured neighborhoods.
  \item The horizontal window must send the graph into the requested vertical band.
\end{enumerate}

\section{Cumulative Review Answers}

\begin{enumerate}[label=\textbf{R\arabic*.}]
  \item The function value concerns the target point; the limit concerns nearby points.
  \item The one-sided limits must both exist and equal the same \(L\).
  \item Jump, unbounded behavior, oscillation, or missing domain on a required side.
  \item Defined value, existing limit, equality.
  \item A hole is caused by one missing or wrong point; a jump is caused by different neighborhoods.
  \item Infinity describes unbounded behavior rather than a final real output.
  \item See Section 3 theorem statement.
  \item \(\lim_{x\to0}\sin x/x=1\), in radians.
  \item See Section 5 theorem statement.
  \item It does not provide exact location or uniqueness.
  \item \(14\).
  \item \(1\).
  \item \(8\).
  \item \(12\).
  \item \(2\).
  \item \(1/10\).
  \item \(6\).
  \item \(-1/4\).
  \item \(\DNE\).
  \item \(0\).
  \item \(0\).
  \item \(0\).
  \item \(6\).
  \item \(2/9\).
  \item \(5/3\).
  \item \(8\).
  \item \(21\).
  \item \(1/2\).
  \item \(-\infty\).
  \item \(+\infty\).
  \item \(-\infty\).
  \item Hole at \(1\); vertical asymptote at \(2\).
  \item \(2\).
  \item \(0\).
  \item \(+\infty\).
  \item \(y=x+1\).
  \item \(2\).
  \item \(-2\).
  \item \(6\).
  \item \(-4\).
  \item \([-1,2)\cup(2,\infty)\).
  \item Hole at \(2\); vertical asymptote at \(-1\).
  \item \(c=6\).
  \item \(k=3\).
  \item No value of \(a\) works.
  \item Polynomial continuity; \(f(1)=-1\), \(f(2)=5\).
  \item The function is not continuous at zero.
  \item After two steps, \((1.25,1.5)\).
  \item \(\delta=\eps/4\).
  \item \(\delta=\min\{1,\eps/3\}\) works.
  \item Nonlinear output error may amplify input error by an additional factor.
  \item \((4.99,5.01)\) and \((2.98,3.02)\).
\end{enumerate}

\section{Practice Examination A Solutions}

\begin{enumerate}[label=\textbf{A\arabic*.}]
  \item \(f(2)=7\), while the limit is \(4\). The filled point gives the first; nearby graph behavior gives the second.
  \item Defined value, existing two-sided limit, equality between them.
  \item It records simultaneous numerator and denominator convergence to zero but does not determine their relative rates.
  \item IVT guarantees at least one attained intermediate value under continuity; it does not calculate or prove uniqueness.
  \item \(-13\).
  \item Factor and cancel to obtain \(x+3\); answer \(6\).
  \item Rationalization gives \(1/(\sqrt{x+9}+3)\); answer \(1/6\).
  \item Combine fractions to obtain \(-1/(2x)\); at \(2\), answer \(-1/4\).
  \item \((7/2)(\sin7x/7x)\to7/2\).
  \item Squeezed between \(-x^2\) and \(x^2\); answer \(0\).
  \item Left \(-1\), right \(1\); answer \(\DNE\).
  \item Numerator approaches \(3>0\); denominator changes sign. Left \(-\infty\), right \(+\infty\).
  \item \(5/2\).
  \item \(-1\).
  \item Hole at \(1\), vertical asymptote \(x=2\), horizontal asymptote \(y=1\).
  \item \(2k+1=3\), so \(k=1\).
  \item The polynomial is continuous; values at \(1,2\) are \(-1,7\), so a root exists.
  \item \(\delta=\eps/2\).
\end{enumerate}

\section{Practice Examination B Solutions}

\begin{enumerate}[label=\textbf{B\arabic*.}]
  \item Left \(2\), right \(2\), two-sided limit \(2\), function value \(5\); not continuous, removable discontinuity.
  \item Cancel \(x-2\): \((x^2+2x+4)/(x+2)\to12/4=3\).
  \item Rationalization gives \(5/[\sqrt{1+4x}+\sqrt{1-x}]\to5/2\).
  \item \(9\cdot(1/2)=9/2\).
  \item \(2/5\).
  \item Squeeze gives \(0\).
  \item Hole at \(1\); vertical asymptote \(x=3\). Both one-sided limits at \(3\) are \(+\infty\).
  \item \(+\infty\).
  \item Division gives polynomial asymptote \(y=2x\).
  \item \(-5\).
  \item Continuity at \(2\) gives \(6=6+a\), so \(a=0\); continuity at \(1\) then gives \(a+b=3\), so \(b=3\).
  \item The polynomial is continuous, \(f(0)=1\), \(f(1)=-1\). First midpoint \(0.5\) is negative, so interval \((0,0.5)\). Second midpoint \(0.25\) is positive, so interval \((0.25,0.5)\).
  \item Choose \(\delta=\eps/4\); then \(|(4x-3)-1|=4|x-1|<4\delta=\eps\).
  \item Counterexample: \(f(x)=1/x\) on \([-1,1]\). Endpoint signs differ, but the function is not continuous on the interval. The missing hypothesis is continuity.
\end{enumerate}

% --- END INLINE appendices/answers ---


% --- BEGIN INLINE appendices/graph_specs ---
\chapter{Digital Graph Specifications}

This appendix translates the static textbook figures into implementation requirements for interactive BetterGrades graphics. Each interactive should support keyboard controls, screen-reader descriptions, a high-contrast mode, and a text alternative containing the relevant equations and conclusions.

\section{Secant Lines Approaching a Tangent}

\begin{bggraphspec}
\textbf{Function:}
\[
  s(t)=64t-16t^2.
\]
\textbf{Primary window:} \(0\le t\le4\), \(0\le s\le72\). Fix \(P=(1,48)\). Define a movable point
\[
  Q_h=(1+h,s(1+h)),\qquad -0.8\le h\le1.5,\ h\ne0.
\]
Display the secant slope
\[
  m_h=\frac{s(1+h)-s(1)}h=32-16h.
\]
Display the tangent line
\[
  s-48=32(t-1).
\]
Controls should allow logarithmic stepping of \(|h|\) through \(1,0.5,0.2,0.1,0.05,0.01\), with a separate direction toggle. Show a synchronized table containing \(h\), \(Q_h\), and \(m_h\).
\end{bggraphspec}

\section{Removable Hole}

\begin{bggraphspec}
\textbf{Function:}
\[
  f(x)=\frac{x^2-4}{x-2},\qquad x\ne2.
\]
\textbf{Equivalent nearby rule:} \(f(x)=x+2\). \textbf{Window:} \(-1\le x\le5\), \(0\le y\le7\). Draw an open circle at \((2,4)\), no filled point, and optional table markers at \(1.9,1.99,1.999,2.001,2.01,2.1\). A toggle should add an arbitrary value \(f(2)=c\) without changing the displayed limit.
\end{bggraphspec}

\section{Function Value Versus Limit}

\begin{bggraphspec}
\textbf{Function:}
\[
  g(x)=\begin{cases}x^2+1,&x\ne2,\\9,&x=2.\end{cases}
\]
\textbf{Window:} \(-1\le x\le4\), \(0\le y\le11\). Use an open circle at \((2,5)\) and a filled point at \((2,9)\). Label the open point as the limiting value and the filled point as the actual function value. An adjustable filled point should demonstrate that changing \(g(2)\) leaves the limit unchanged.
\end{bggraphspec}

\section{Jump Discontinuity}

\begin{bggraphspec}
\textbf{Function:}
\[
  g(x)=\begin{cases}x+1,&x<2,\\6-x,&x\ge2.\end{cases}
\]
\textbf{Window:} \(-1\le x\le6\), \(0\le y\le7\). Draw an open circle at \((2,3)\) and a filled point at \((2,4)\). Provide separate left-approach and right-approach animations with labels
\[
  \lim_{x\to2^-}g(x)=3,
  \qquad
  \lim_{x\to2^+}g(x)=4.
\]
The two-sided limit readout must state \(\DNE\), not an average of \(3\) and \(4\).
\end{bggraphspec}

\section{Oscillation Near Zero}

\begin{bggraphspec}
\textbf{Function:}
\[
  f(x)=\sin(1/x),\qquad x\ne0.
\]
Use a zoom control with symmetric windows \([-r,r]\) for \(r=1,0.5,0.25,0.1,0.05\). Increase sampling density as \(r\) decreases. The graphic should emphasize that zooming does not make the function settle. Include two input sequences approaching zero that produce different fixed output behaviors, such as
\[
  x_n=\frac{1}{\pi/2+2\pi n}\quad\Rightarrow\quad f(x_n)=1,
\]
\[
  y_n=\frac{1}{3\pi/2+2\pi n}\quad\Rightarrow\quad f(y_n)=-1.
\]
\end{bggraphspec}

\section{Squeeze Theorem Oscillation}

\begin{bggraphspec}
\textbf{Functions:}
\[
  y=x^2\sin(1/x),\qquad y=x^2,\qquad y=-x^2.
\]
\textbf{Window:} \(-0.35\le x\le0.35\), \(-0.13\le y\le0.13\). Shade the region between the two parabolas. Include a slider moving toward zero and a live inequality
\[
  -x^2\le x^2\sin(1/x)\le x^2.
\]
The screen-reader description should state that the amplitude is at most \(x^2\), which approaches zero.
\end{bggraphspec}

\section{Unit-Circle Squeeze}

\begin{bggraphspec}
Draw a unit circle with angle \(x\in(0,\pi/2)\). Display the inner triangle, circular sector, and outer tangent triangle. Show their areas dynamically:
\[
  A_{\mathrm{inner}}=\frac12\sin x\cos x,
\]
\[
  A_{\mathrm{sector}}=\frac12x,
\]
\[
  A_{\mathrm{outer}}=\frac12\tan x.
\]
Use a slider for \(x\), measured only in radians, and display the resulting inequality chain leading to
\[
  \cos x<\frac{\sin x}{x}<1.
\]
\end{bggraphspec}

\section{The Sinc Function}

\begin{bggraphspec}
\textbf{Function:}
\[
  f(x)=\frac{\sin x}{x},\qquad x\ne0.
\]
\textbf{Wide window:} \([-2\pi,2\pi]\times[-0.35,1.18]\). \textbf{Zoom window:} \([-0.5,0.5]\times[0.95,1.01]\). Draw an open circle at \((0,1)\) and a dashed line \(y=1\). Include a unit toggle locked to radians by default; if degrees are demonstrated, clearly show that the limiting value changes to \(\pi/180\).
\end{bggraphspec}

\section{Vertical Asymptote Multiplicity}

\begin{bggraphspec}
Compare
\[
  f(x)=\frac1{x-1}
  \qquad\text{and}\qquad
  g(x)=\frac1{(x-1)^2}.
\]
Use matched windows and a shared vertical asymptote \(x=1\). Show sign tables on each side. Emphasize that an odd power changes sign across the asymptote while an even power preserves sign.
\end{bggraphspec}

\section{Horizontal Asymptote}

\begin{bggraphspec}
\textbf{Function:}
\[
  f(x)=\frac{3x^2-2x+5}{x^2+4}.
\]
\textbf{Primary window:} \([-12,12]\times[0,4.3]\). Draw \(y=3\) as a dashed asymptote. Add a second view or zoom control extending to \(|x|=100\). Display
\[
  f(x)-3=\frac{-2x-7}{x^2+4}
\]
to show quantitatively why the difference approaches zero.
\end{bggraphspec}

\section{Discontinuity Gallery}

\begin{bggraphspec}
Provide four linked panels:
\begin{enumerate}
  \item removable: \((x^2-1)/(x-1)\);
  \item jump: \(0\) for \(x<0\), \(2\) for \(x\ge0\);
  \item infinite: \(1/x\);
  \item oscillatory: \(\sin(1/x)\).
\end{enumerate}
Each panel should list the left limit, right limit, function value, two-sided limit, and classification. Avoid relying on color alone; use distinct marker styles and labels.
\end{bggraphspec}

\section{Intermediate Value Theorem}

\begin{bggraphspec}
\textbf{Function:}
\[
  f(x)=x^3+x-1
\]
on \([0,1]\). Mark \((0,-1)\), \((1,1)\), and the root near \(0.6823278\). Include a draggable horizontal target \(N\) between \(-1\) and \(1\). The graphic should find at least one corresponding \(c\) visually while stating that the theorem guarantees existence, not uniqueness or a closed-form value.
\end{bggraphspec}

\section{Epsilon-Delta Window}

\begin{bggraphspec}
Use \(f(x)=x^2/2+1\), \(a=2\), and \(L=3\). Provide an \(\eps\) slider controlling the horizontal band \((L-\eps,L+\eps)\). Compute or approximate the largest symmetric \(\delta\)-window about \(a\) whose graph remains in the band. Display the inequalities
\[
  |x-a|<\delta
\]
and
\[
  |f(x)-L|<\eps
\]
with synchronized highlighting. Include a mode in which a learner proposes \(\delta\) and the graphic identifies any violating points.
\end{bggraphspec}

% --- END INLINE appendices/graph_specs ---


% --- BEGIN INLINE appendices/references ---
\chapter{Sources and Licensing Notes}

\section*{Primary pedagogical references}

\begin{thebibliography}{9}

\bibitem{activecalculus}
Matthew Boelkins, David Austin, Christina Safranski, and Steven Schlicker.
\emph{Active Calculus: Single Variable}, 2nd ed., 2025.
Creative Commons Attribution-ShareAlike 4.0 International.

\bibitem{thompson}
Silvanus P. Thompson.
\emph{Calculus Made Easy}, 2nd ed., 1914.
Public-domain United States edition distributed by Project Gutenberg.

\bibitem{granville}
William Anthony Granville and Percey F. Smith.
\emph{Elements of the Differential and Integral Calculus}.
Public-domain historical edition.

\bibitem{greenhill}
Alfred George Greenhill.
\emph{Differential and Integral Calculus, with Applications}, 1886.
Public domain.

\end{thebibliography}

\section*{How the sources were used}

This BetterGrades manuscript uses original prose, original organization, original worked examples, and newly composed exercises. The source texts were consulted to verify topic coverage, compare pedagogical sequences, and identify durable classes of calculus problems.

\emph{Active Calculus} was especially useful as a modern reference for moving among numerical, graphical, algebraic, and verbal representations and for emphasizing active student work. Because that book is licensed CC BY-SA 4.0, no exercise or extended passage from it is reproduced verbatim in this manuscript.

\emph{Calculus Made Easy}, Granville and Smith, and Greenhill are public-domain sources in the United States. Historical problem settings drawn from their tradition have been modernized and recomposed rather than mechanically transcribed. Their strongest contribution here is the insistence that calculus answer intelligible questions about motion, rates, approximation, and geometry before becoming a catalogue of symbols.

\section*{Rights-separation policy}

Any future BetterGrades ingestion pipeline should preserve source-level provenance. Public-domain exercises may be stored in a public-domain corpus. Material copied or adapted from CC BY-SA sources must remain separately identified, attributed, and distributed under compatible ShareAlike terms. Original BetterGrades exercises should be marked as original rather than quietly laundering provenance through a database, the sort of administrative miracle that eventually becomes a legal problem.

% --- END INLINE appendices/references ---

\end{document}
